% Oscillations électriques forcées en régime sinusoïdal — Physique Tle D — Cours signé OnBuch+
% Programme officiel MINESEC. Compiler avec : tectonic P12-oscillations-electriques-forcees.tex
\newcommand{\DOCMATIERE}{Physique}
\newcommand{\DOCNIVEAU}{Tle D}
\newcommand{\DOCMODULE}{Module 3 : Électricité — évolutions temporelles des circuits électriques et électroniques}
\newcommand{\DOCLECON}{P12}
\newcommand{\DOCTITRE}{Oscillations électriques forcées en régime sinusoïdal}
% OnBuch+ — préambule « cours signés » ENRICHI (compilé avec tectonic / XeLaTeX)
% Système visuel « Pop » d'OnBuch+ : fond crème (jamais blanc), bordures encre
% épaisses, ombres « dures » décalées, titres Archivo Black en capitales.
% Version MATHÉMATIQUES : graphes (pgfplots/tikz), boîtes pédagogiques
% (définition, propriété, méthode, attention, le savais-tu, exercice résolu,
% exercice, TP, bilan), page de garde et signature OnBuch+ ancrée en bas.
%
% Macros attendues dans main.tex AVANT \input{preamble} :
%   \DOCMATIERE  \DOCNIVEAU  \DOCMODULE  \DOCLECON (numéro)  \DOCTITRE
\documentclass[11pt]{article}

\usepackage{fontspec}
\usepackage[french]{babel}
\usepackage[a4paper,margin=2cm,top=3.4cm,bottom=2.8cm,headheight=1.4cm,footskip=1.3cm]{geometry}
\usepackage{amsmath,amssymb,amsfonts}
\usepackage{mathtools} % \xmapsto, \coloneqq... souvent utilisés par les modèles
\usepackage{cancel} % simplifications barrées (bilans de forces)
% \abs{z} (module, valeur absolue) et \norm{u} (norme), utilisés par les modèles
\providecommand{\abs}{}\let\abs\relax\DeclarePairedDelimiter{\abs}{\lvert}{\rvert}
\providecommand{\norm}{}\let\norm\relax\DeclarePairedDelimiter{\norm}{\lVert}{\rVert}
\usepackage[table]{xcolor} % [table] : fournit aussi \rowcolors (alternance de lignes)
\usepackage{pagecolor}
\usepackage{tikz}
\usetikzlibrary{arrows.meta,positioning,calc,shapes.geometric,shapes.misc,shapes.arrows,decorations.pathmorphing,decorations.markings,patterns,shadows,mindmap,trees,fit,backgrounds,matrix,intersections,angles,quotes}
\usepackage{pgfplots}
\usepackage[version=4]{mhchem} % équations nucléaires \ce{^{235}_{92}U}
\usepackage[european,siunitx]{circuitikz} % schémas électriques
\pgfplotsset{compat=1.18}
\usepgfplotslibrary{fillbetween} % aires entre courbes (calcul intégral)
\usepackage{enumitem}
\usepackage{fancyhdr}
% [most] charge toutes les bibliothèques tcolorbox (listings, minted, raster,
% documentation...) et gonfle inutilement la mémoire TeX sur un document long
% (~300 boîtes) : on ne charge que ce qui est réellement utilisé.
\usepackage[skins,breakable]{tcolorbox}
\usepackage{graphicx}
\usepackage{colortbl}
\usepackage{booktabs}
\usepackage{tabularx}
\usepackage{diagbox} % case d'en-tête diagonale des tableaux à double entrée
% Colonnes extensibles centrée / gauche / droite (C, L, R), souvent utilisées par les modèles
\newcolumntype{C}{>{\centering\arraybackslash}X}
\newcolumntype{L}{>{\raggedright\arraybackslash}X}
\newcolumntype{R}{>{\raggedleft\arraybackslash}X}
\usepackage{array}
\usepackage{multicol}
\usepackage{siunitx}
% parse-numbers=false : accepte aussi \SI{2,5 \times 10^{-3}}{...}
\sisetup{locale=FR,output-decimal-marker={,},group-minimum-digits=4,parse-numbers=false}
\DeclareSIUnit\ohm{\ensuremath{\symup{\Omega}}} % U+2126 absent de la police
\DeclareSIUnit\division{div} % réglages d'oscilloscope (ms/div, V/div)
\usepackage{qrcode}
% \degree : commande fréquemment utilisée par les modèles pour un angle ou une
% température en dehors de \SI{}{} (non fournie par les paquets chargés).
\providecommand{\degree}{^{\circ}}
% \qed (fin de démonstration), utilisé par les modèles sans amsthm
\providecommand{\qed}{\hfill$\square$}
% \barre : double barre des valeurs interdites dans un tableau de variations (tkz-tab)
\providecommand{\barre}{\ensuremath{\|}}
% environnement proof (amsthm non chargé) utilisé par les modèles
\ifdefined\proof\else\newenvironment{proof}[1][Démonstration]{\par\noindent\textit{#1.}\ }{\qed\par}\fi
\usepackage{lastpage}
\usepackage{hyperref}
\hypersetup{hidelinks}

% ============================================================
% Palette « Pop » OnBuch+ — mêmes hex que la classe P (plus_ui.dart)
% ============================================================
\definecolor{poporange}{HTML}{F59321}
\definecolor{popdark}{HTML}{E07A0C}
\definecolor{popcream}{HTML}{FFF6E8}
\definecolor{popink}{HTML}{1C1714}
\definecolor{poppurple}{HTML}{7A5AE0}
\definecolor{popgreen}{HTML}{2E9E6A}
\definecolor{poppink}{HTML}{E0457B}
\definecolor{popblue}{HTML}{2F7DE1}
\definecolor{popgold}{HTML}{C98A12}
\definecolor{popmuted}{HTML}{8A7F76}
\definecolor{popline}{HTML}{EADFD2}
\definecolor{popgray}{HTML}{8A7F76}
\colorlet{popgrey}{popgray}
% Alias tolérants (le modèle nomme parfois les couleurs différemment)
\colorlet{popwhite}{white}
\colorlet{popblack}{popink}
\colorlet{popred}{poppink}
\colorlet{popyellow}{popgold}
% Teintes claires pour fonds de tableaux / zones
\colorlet{popblueL}{popblue!10!white}
\colorlet{poporangeL}{poporange!14!white}
\colorlet{popgreenL}{popgreen!12!white}
\colorlet{poppurpleL}{poppurple!12!white}
\colorlet{poppinkL}{poppink!12!white}
\colorlet{popgoldL}{popgold!14!white}

\pagecolor{popcream}

% ============================================================
% Typographie
% ============================================================
\defaultfontfeatures{Ligatures=TeX}
\setmainfont{PlusJakartaSans-Medium.ttf}[Path=../../../fonts/,BoldFont=PlusJakartaSans-Bold.ttf]
% Maths en sans-serif (Fira Math) avec chiffres et lettres droites de Plus
% Jakarta, pour que les formules se fondent dans le texte.
\usepackage[math-style=ISO,bold-style=ISO]{unicode-math}
\setmathfont{FiraMath-Regular.otf}[Path=../../../fonts/]
\setmathfont{PlusJakartaSans-Medium.ttf}[Path=../../../fonts/,range={up/num,up/{Latin,latin},\mathup}]
\usepackage{icomma} % virgule décimale sans espace en mode math
% Caractères mathématiques Unicode absents des polices de texte (Plus Jakarta,
% Archivo Black) : rendus par leur équivalent mathématique, en texte comme en
% formule — sinon ils s'affichent en carré vide (ex. « Arithmétique dans ℤ »).
\usepackage{newunicodechar}
\newunicodechar{ℝ}{\ensuremath{\mathbb{R}}}
\newunicodechar{ℤ}{\ensuremath{\mathbb{Z}}}
\newunicodechar{ℂ}{\ensuremath{\mathbb{C}}}
\newunicodechar{ℕ}{\ensuremath{\mathbb{N}}}
\newunicodechar{ℚ}{\ensuremath{\mathbb{Q}}}
\newunicodechar{𝔻}{\ensuremath{\mathbb{D}}}
\newunicodechar{α}{\ensuremath{\alpha}}
\newunicodechar{β}{\ensuremath{\beta}}
\newunicodechar{γ}{\ensuremath{\gamma}}
\newunicodechar{δ}{\ensuremath{\delta}}
\newunicodechar{ε}{\ensuremath{\varepsilon}}
\newunicodechar{θ}{\ensuremath{\theta}}
\newunicodechar{λ}{\ensuremath{\lambda}}
\newunicodechar{μ}{\ensuremath{\mu}}
\newunicodechar{σ}{\ensuremath{\sigma}}
\newunicodechar{τ}{\ensuremath{\tau}}
\newunicodechar{φ}{\ensuremath{\varphi}}
\newunicodechar{ψ}{\ensuremath{\psi}}
\newunicodechar{ω}{\ensuremath{\omega}}
\newunicodechar{Σ}{\ensuremath{\Sigma}}
% majuscules produites par \MakeUppercase dans les titres (α-aminés -> Α)
\newunicodechar{Α}{\ensuremath{\alpha}}
\newunicodechar{Β}{\ensuremath{\beta}}
\newunicodechar{Γ}{\ensuremath{\Gamma}}
\newunicodechar{Δ}{\ensuremath{\Delta}}
\newunicodechar{Ε}{\ensuremath{\varepsilon}}
\newunicodechar{Θ}{\ensuremath{\Theta}}
\newunicodechar{Λ}{\ensuremath{\Lambda}}
\newunicodechar{Μ}{\ensuremath{\mu}}
\newunicodechar{Τ}{\ensuremath{\tau}}
\newunicodechar{Φ}{\ensuremath{\Phi}}
\newunicodechar{Ψ}{\ensuremath{\Psi}}
\newunicodechar{Ω}{\ensuremath{\Omega}}
\newunicodechar{↦}{\ensuremath{\mapsto}}
\newunicodechar{∈}{\ensuremath{\in}}
\newunicodechar{∉}{\ensuremath{\notin}}
\newunicodechar{∧}{\ensuremath{\wedge}}
\newunicodechar{⇔}{\ensuremath{\Leftrightarrow}}
\newunicodechar{⟺}{\ensuremath{\Longleftrightarrow}}
\newunicodechar{⇒}{\ensuremath{\Rightarrow}}
\newunicodechar{⟹}{\ensuremath{\Longrightarrow}}
\newunicodechar{∘}{\ensuremath{\circ}}
\newunicodechar{∪}{\ensuremath{\cup}}
\newunicodechar{∩}{\ensuremath{\cap}}
\newunicodechar{≡}{\ensuremath{\equiv}}
\newunicodechar{⊂}{\ensuremath{\subset}}
\newunicodechar{⊥}{\ensuremath{\perp}}
\newunicodechar{∃}{\ensuremath{\exists}}
\newunicodechar{∀}{\ensuremath{\forall}}
\newunicodechar{∛}{\ensuremath{\sqrt[3]{\,}}}
\newunicodechar{∜}{\ensuremath{\sqrt[4]{\,}}}
\newunicodechar{⃗}{\ensuremath{{}^{\to}}}
\newunicodechar{ⁿ}{\ifmmode^{n}\else\textsuperscript{\ensuremath{n}}\fi}
\newunicodechar{ˣ}{\ifmmode^{x}\else\textsuperscript{\ensuremath{x}}\fi}
\newunicodechar{ᵇ}{\ifmmode^{b}\else\textsuperscript{\ensuremath{b}}\fi}
\newunicodechar{ʳ}{\ifmmode^{r}\else\textsuperscript{\ensuremath{r}}\fi}
\newunicodechar{ᵘ}{\ifmmode^{u}\else\textsuperscript{\ensuremath{u}}\fi}
\newunicodechar{ʸ}{\ifmmode^{y}\else\textsuperscript{\ensuremath{y}}\fi}
\newunicodechar{ᵃ}{\ifmmode^{a}\else\textsuperscript{\ensuremath{a}}\fi}
\newunicodechar{ᵉ}{\ifmmode^{e}\else\textsuperscript{\ensuremath{e}}\fi}
\newunicodechar{ᵏ}{\ifmmode^{k}\else\textsuperscript{\ensuremath{k}}\fi}
\newunicodechar{ᵖ}{\ifmmode^{p}\else\textsuperscript{\ensuremath{p}}\fi}
\newunicodechar{ᴺ}{\ifmmode^{N}\else\textsuperscript{\ensuremath{N}}\fi}
\newunicodechar{ᶜ}{\ifmmode^{c}\else\textsuperscript{\ensuremath{c}}\fi}
\newunicodechar{ʰ}{\ifmmode^{h}\else\textsuperscript{\ensuremath{h}}\fi}
\newunicodechar{ᵠ}{\ifmmode^{q}\else\textsuperscript{\ensuremath{q}}\fi}
\newunicodechar{ₙ}{\ifmmode_{n}\else\textsubscript{\ensuremath{n}}\fi}
\newunicodechar{ₐ}{\ifmmode_{a}\else\textsubscript{\ensuremath{a}}\fi}
\newunicodechar{ᵢ}{\ifmmode_{i}\else\textsubscript{\ensuremath{i}}\fi}
\newunicodechar{ᵣ}{\ifmmode_{r}\else\textsubscript{\ensuremath{r}}\fi}
\newunicodechar{ₖ}{\ifmmode_{k}\else\textsubscript{\ensuremath{k}}\fi}
\newunicodechar{ᵦ}{\ifmmode_{\beta}\else\textsubscript{\ensuremath{\beta}}\fi}
\newunicodechar{ₓ}{\ifmmode_{x}\else\textsubscript{\ensuremath{x}}\fi}
\newfontfamily\popdisplay{ArchivoBlack-Regular.ttf}[Path=../../../fonts/]
\newfontfamily\poplogo{SpaceGrotesk-Bold.ttf}[Path=../../../fonts/]
\newfontfamily\popsemi{PlusJakartaSans-SemiBold.ttf}[Path=../../../fonts/]
\newfontfamily\popxbold{PlusJakartaSans-ExtraBold.ttf}[Path=../../../fonts/]
\newfontfamily\popxboldls{PlusJakartaSans-ExtraBold.ttf}[Path=../../../fonts/,LetterSpace=3.0]

\setlist[enumerate]{leftmargin=1.7em,itemsep=5pt,topsep=4pt}
\setlist[itemize]{leftmargin=1.7em,itemsep=4pt,topsep=4pt,label={\color{poporange}\textbullet}}
\setlength{\parindent}{0pt}
\setlength{\parskip}{5pt}
\renewcommand{\arraystretch}{1.25}

% pgfplots : style Pop par défaut
\pgfplotsset{
  every axis/.append style={axis line style={popink,line width=1pt},tick style={popink},
    grid style={popline,line width=0.5pt},label style={font=\small},tick label style={font=\footnotesize}},
  popcurve/.style={poporange,line width=1.8pt,no markers,smooth},
}
% Même style disponible dans un \draw TikZ simple (hors pgfplots)
\tikzset{popcurve/.style={poporange,line width=1.8pt}}
\tikzset{popgrid/.style={popline,very thin}}

% ============================================================
% Pill / squiggle
% ============================================================
\newcommand{\poppill}[3]{%
  \tikz[baseline=-0.55ex]\node[draw=popink,line width=1.2pt,rounded corners=2.6pt,%
    fill=#2,inner xsep=6.5pt,inner ysep=3pt]{%
    \popxboldls\fontsize{7}{7}\selectfont\color{#3}\MakeUppercase{#1}};%
}
\newcommand{\popsquiggle}[1]{%
  \tikz[baseline=-0.4ex]{
    \draw[#1,line width=2.6pt,line cap=round]
      plot [smooth] coordinates {(0,0.09) (0.34,0.01) (0.68,0.21) (1.02,0.01) (1.36,0.10)};
  }%
}
% Mot-clé surligné (marqueur orange sous le texte)
\newcommand{\cle}[1]{{\popxbold\color{popdark}#1}}

% ============================================================
% En-tête / pied
% ============================================================
\pagestyle{fancy}
\fancyhf{}
\renewcommand{\headrulewidth}{0pt}
\renewcommand{\footrulewidth}{0pt}
\lhead{\raisebox{-0.15ex}{\poplogo\fontsize{13}{13}\selectfont\color{popink}OnBuch\color{poporange}+}}
\rhead{\poppill{\DOCMATIERE\ \textbullet\ \DOCNIVEAU\ \textbullet\ Leçon \DOCLECON}{white}{popink}}
\lfoot{\popxbold\fontsize{7.6}{7.6}\selectfont\color{popmuted}\MakeUppercase{Cours signé OnBuch+}\ \textbullet\ {\popsemi\DOCTITRE}}
\rfoot{\tikz[baseline=-0.6ex]\node[rounded corners=8.5pt,fill=popink,minimum height=17pt,minimum width=17pt,inner xsep=5pt,inner sep=0]{\popxbold\fontsize{8}{8}\selectfont\color{popcream}\thepage\,/\,\pageref*{LastPage}};}

% ============================================================
% Titres
% ============================================================
\newcommand{\chaptitre}[2]{%
  \begin{tcolorbox}[enhanced,colback=poporange,colframe=popink,boxrule=2.6pt,arc=11pt,
      drop shadow=popink,left=15pt,right=15pt,top=12pt,bottom=11pt,before skip=3pt,after skip=18pt]
    \poppill{#2}{popink}{popcream}\par\vspace{9pt}
    {\raggedright\hyphenpenalty=10000\exhyphenpenalty=10000\popdisplay\fontsize{22}{24}\selectfont\color{popcream}\MakeUppercase{#1}\par}
  \end{tcolorbox}
}
% Section numérotée : pastille orange avec le numéro + titre Archivo
\newcounter{popsec}
\newcommand{\coursec}[1]{%
  \stepcounter{popsec}\setcounter{popsub}{0}%
  \par\vspace{16pt}\noindent
  \begin{minipage}{\linewidth}
  \tikz[baseline=-0.7ex]\node[draw=popink,line width=1.6pt,fill=poporange,rounded corners=4pt,minimum size=22pt,inner sep=0]{\popdisplay\fontsize{12}{12}\selectfont\color{popcream}\thepopsec};%
  \hspace{8pt}{\popdisplay\fontsize{15}{17}\selectfont\color{popink}\MakeUppercase{#1}}\par
  \vspace{4pt}\noindent{\color{poporange}\rule{36pt}{3.6pt}}
  \end{minipage}\par\nopagebreak\vspace{8pt}
}
\newcounter{popsub}
\newcommand{\courssub}[1]{%
  \stepcounter{popsub}%
  \par\vspace{9pt}\noindent{\popxbold\fontsize{11.5}{13}\selectfont\color{popink}{\color{poporange}\thepopsec.\thepopsub}\hspace{6pt}#1}\par\nopagebreak\vspace{4pt}
}

% ============================================================
% Boîtes pédagogiques — même recette : fond blanc, bordure épaisse colorée,
% ombre dure encre, badge plein. Titre optionnel [#1].
% ============================================================
\tcbset{popbase/.style={breakable,enhanced,colback=white,boxrule=2.3pt,arc=9pt,
  shadow={3pt}{-3pt}{0pt}{popink},left=14pt,right=14pt,top=11pt,bottom=11pt,before skip=11pt,after skip=15pt,
  fontupper=\popsemi\fontsize{10.6}{14.5}\selectfont\color{popink},
  fontlower=\popsemi\fontsize{10.6}{14.5}\selectfont\color{popink}}}
% Coche dessinée (le glyphe ✓ n'existe pas dans les polices du document)
\newcommand{\popcheck}[1]{\tikz[baseline=-0.5ex]\draw[#1,line width=1.6pt,line cap=round,line join=round](-0.13,0.0)--(-0.04,-0.09)--(0.13,0.1);}
\providecommand{\coche}{\popcheck{popgreen}}
\providecommand{\resultat}[1]{\par\smallskip\noindent\fcolorbox{popgreen}{popgreenL}{\parbox{\dimexpr\linewidth-2\fboxsep-2\fboxrule\relax}{\textbf{Résultat.} #1}}\par\smallskip}
\newcommand{\popboxhead}[3]{\poppill{#1}{#2}{white}\ifx\relax#3\relax\else\hspace{6pt}{\popxbold\fontsize{10.6}{12}\selectfont\color{popink}#3}\fi\par\vspace{7pt}}

\newtcolorbox{aretenir}[1][]{popbase,colframe=popgreen,before upper={\popboxhead{À retenir}{popgreen}{#1}}}
\newtcolorbox{exemplebox}[1][]{popbase,colframe=poporange,before upper={\popboxhead{Exemple}{poporange}{#1}}}
\newtcolorbox{definition}[1][]{popbase,colframe=popblue,before upper={\popboxhead{Définition}{popblue}{#1}}}
\newtcolorbox{propriete}[1][]{popbase,colframe=poppurple,before upper={\popboxhead{Propriété}{poppurple}{#1}}}
\newtcolorbox{methode}[1][]{popbase,colframe=popgold,before upper={\popboxhead{Méthode}{popgold}{#1}}}
\newtcolorbox{attention}[1][]{popbase,colframe=poppink,before upper={\popboxhead{Attention}{poppink}{#1}}}
\newtcolorbox{experience}[1][]{popbase,colframe=popblue,colback=popblueL,before upper={\popboxhead{Expérience}{popblue}{#1}}}
% « Le savais-tu ? » : carte violette pleine (contexte, culture, Cameroun)
\newtcolorbox{savaistu}[1][]{popbase,colback=poppurple,colframe=popink,
  fontupper=\popsemi\fontsize{10.4}{14.2}\selectfont\color{white},
  before upper={\poppill{Le savais-tu\,?}{white}{poppurple}\ifx\relax#1\relax\else\hspace{6pt}{\popxbold\color{white}#1}\fi\par\vspace{7pt}}}
% Exercice résolu : énoncé en haut, \tcblower puis la solution sur fond crème
\newcounter{popexo}\newcounter{popexr}
\newtcolorbox{exoresolu}[1][]{popbase,colframe=poporange,colbacklower=popcream,
  segmentation style={popink,line width=1.2pt,dashed},
  before upper={\stepcounter{popexr}\popboxhead{Exercice résolu \thepopexr}{poporange}{#1}},
  before lower={\poppill{Solution}{popink}{popcream}\par\vspace{6pt}}}
% Exercice (non résolu en place) — la correction est dans la partie Corrigés
\newtcolorbox{exercice}[1][]{popbase,colframe=popink,
  before upper={\stepcounter{popexo}\popboxhead{Exercice \thepopexo}{popink}{#1}}}
% Corrigé
\newtcolorbox{corrige}[1][]{popbase,colframe=popgreen,colback=popgreenL,
  before upper={\popboxhead{Corrigé}{popgreen}{#1}}}
% Objectifs / prérequis / bilan
\newtcolorbox{objectifs}{popbase,colframe=popink,colback=poporangeL,
  before upper={\popboxhead{Objectifs de la leçon}{popink}{}}}
\newtcolorbox{prerequis}{popbase,colframe=popink,colback=popblueL,
  before upper={\popboxhead{Prérequis}{popblue}{}}}
\newtcolorbox{bilan}{popbase,colframe=popink,colback=white,boxrule=2.8pt,
  before upper={\popboxhead{Fiche bilan}{poporange}{L'essentiel à connaître pour l'examen}}}
% Figure encadrée avec légende
\newtcolorbox{popfigure}[1][]{enhanced,breakable=false,colback=white,colframe=popline,boxrule=1.4pt,arc=8pt,
  left=10pt,right=10pt,top=10pt,bottom=8pt,before skip=10pt,after skip=12pt,halign=center,
  lower separated=false,halign lower=center,
  fontlower=\popsemi\fontsize{9}{11.5}\selectfont\color{popmuted},
  #1}
\newcommand{\legende}[1]{\par\vspace{6pt}{\popsemi\fontsize{9}{11.5}\selectfont\color{popmuted}#1}}

% ============================================================
% Signature OnBuch+
% ============================================================
\newcommand{\poplogoinline}[1]{{\poplogo\fontsize{#1}{#1}\selectfont\color{popink}OnBuch\color{poporange}+}}
\newcommand{\signatureonbuch}{%
  \begin{tcolorbox}[enhanced,colback=popink,colframe=popink,boxrule=2.2pt,arc=10pt,
      drop shadow=poporange,left=14pt,right=14pt,top=10pt,bottom=10pt,before skip=0pt,after skip=0pt]
    \tikz[baseline=-0.7ex]\node[circle,fill=popgreen,draw=popcream,line width=1.3pt,minimum size=22pt,inner sep=0]{\popcheck{white}};%
    \hspace{9pt}\begin{minipage}[c]{0.62\linewidth}
      {\popxbold\fontsize{12}{13}\selectfont\color{popcream}Signé\ \textbullet\ {\poplogo OnBuch\color{poporange}+}}\par\vspace{2pt}
      {\raggedright\popsemi\fontsize{8}{10.5}\selectfont\color{popline}\MakeUppercase{Équipe pédagogique OnBuch+}\\\MakeUppercase{Conforme au programme officiel MINESEC}\par}
    \end{minipage}\hfill
    {\poplogo\fontsize{22}{22}\selectfont\color{popcream}OnBuch\color{poporange}+}
  \end{tcolorbox}%
}

% ============================================================
% Page de garde
% #1 titre · #2 matière · #3 niveau · #4 module · #5 n° leçon · #6 durée ·
% #7 description / objectifs (texte ou itemize)
% ============================================================
\newcommand{\pagegarde}[7]{%
  \thispagestyle{empty}
  \newgeometry{a4paper,margin=1.7cm,top=1.6cm,bottom=1.4cm}%
  \begin{tikzpicture}[remember picture,overlay]
    \fill[poporange!18] ([xshift=-3.2cm,yshift=-3.6cm]current page.north east) circle (4.6cm);
    \fill[poppurple!14] ([xshift=2.4cm,yshift=7.6cm]current page.south west) circle (3.8cm);
  \end{tikzpicture}%
  {\poplogo\fontsize{30}{30}\selectfont\color{popink}OnBuch\color{poporange}+}\hfill
  \raisebox{0.6ex}{\poppill{Cours signé \textbullet\ Premium}{popink}{popcream}}\par
  \vspace{1pt}\popsquiggle{poppurple}\par
  \vspace{8pt}
  \begin{tcolorbox}[enhanced,colback=poporange,colframe=popink,boxrule=2.8pt,arc=13pt,
      drop shadow=popink,left=18pt,right=18pt,top=12pt,bottom=13pt,after skip=10pt]
    \poppill{#2 \textbullet\ #3}{popink}{popcream}\hspace{5pt}\poppill{#4}{white}{popink}\par\vspace{8pt}
    {\popdisplay\fontsize{14}{14}\selectfont\color{popink}LEÇON #5}\par\vspace{4pt}
    {\raggedright\hyphenpenalty=10000\exhyphenpenalty=10000\popdisplay\fontsize{25}{27}\selectfont\color{popcream}\MakeUppercase{#1}\par}
    \vspace{8pt}
    {\popxbold\fontsize{9.5}{9.5}\selectfont\color{popink}\MakeUppercase{Durée conseillée : #6}}
  \end{tcolorbox}
  \begin{tcolorbox}[enhanced,colback=white,colframe=popink,boxrule=2.2pt,arc=10pt,
      drop shadow=popink,left=16pt,right=16pt,top=10pt,bottom=10pt,after skip=8pt]
    \poppill{Dans cette leçon}{poppurple}{white}\par\vspace{5pt}
    \popsemi\fontsize{9.3}{12.6}\selectfont\color{popink}#7
  \end{tcolorbox}
  \noindent\begin{minipage}[t]{0.487\linewidth}
    \begin{tcolorbox}[enhanced,colback=popgreen,colframe=popink,boxrule=2.2pt,arc=10pt,
        drop shadow=popink,left=11pt,right=11pt,top=10pt,bottom=10pt,height=3.1cm,valign=center]
      \tcbox[colback=white,colframe=popink,boxrule=1.3pt,arc=5pt,boxsep=0pt,nobeforeafter,
        left=3pt,right=3pt,top=3pt,bottom=3pt,valign=center]{\qrcode[height=1.9cm]{https://play.google.com/store/apps/details?id=com.luvvix.onbuch&hl=fr}}%
      \hspace{8pt}\begin{minipage}[c]{0.5\linewidth}
        {\popdisplay\fontsize{11}{12.5}\selectfont\color{white}\MakeUppercase{Télécharge OnBuch}}\par\vspace{4pt}
        {\raggedright\popsemi\fontsize{8.4}{11}\selectfont\color{white}Gratuit \textbullet\ Google Play\\Tuteur IA, annales, concours, résultats\par}
      \end{minipage}
    \end{tcolorbox}
  \end{minipage}\hfill
  \begin{minipage}[t]{0.487\linewidth}
    \begin{tcolorbox}[enhanced,colback=poppurple,colframe=popink,boxrule=2.2pt,arc=10pt,
        drop shadow=popink,left=11pt,right=11pt,top=10pt,bottom=10pt,height=3.1cm,valign=center]
      \tikz[baseline=-0.7ex]\node[circle,fill=white,draw=popink,line width=1.3pt,minimum size=1.6cm,inner sep=0]{\popdisplay\fontsize{24}{24}\selectfont\color{poppurple}+};%
      \hspace{8pt}\begin{minipage}[c]{0.6\linewidth}
        {\popdisplay\fontsize{11}{12.5}\selectfont\color{white}\MakeUppercase{Passe à OnBuch+}}\par\vspace{4pt}
        {\raggedright\popsemi\fontsize{8.4}{11}\selectfont\color{white}Tous les cours signés, Tuteur IA illimité, fiches PDF — onglet OnBuch+ de l'app\par}
      \end{minipage}
    \end{tcolorbox}
  \end{minipage}
  \vspace{8pt}
  \popcontenu
  \vfill
  \signatureonbuch
  \restoregeometry
  \newpage
}

% Rangée « Ce cours contient » (page de garde)
\newcommand{\poptile}[3]{% #1 couleur #2 gros texte #3 légende
  \begin{tcolorbox}[enhanced,colback=#1,colframe=popink,boxrule=2pt,arc=9pt,shadow={2.5pt}{-2.5pt}{0pt}{popink},
      left=6pt,right=6pt,top=9pt,bottom=9pt,halign=center,valign=center,height=1.75cm,width=\linewidth,nobeforeafter]
    {\popdisplay\fontsize{12.5}{13}\selectfont\color{white}\MakeUppercase{#2}}\par\vspace{3pt}
    {\centering\popsemi\fontsize{7.8}{9.5}\selectfont\color{white}#3\par}
  \end{tcolorbox}}
\newcommand{\popcontenu}{%
  \par\noindent{\popxboldls\fontsize{7.5}{7.5}\selectfont\color{popmuted}CE COURS SIGNÉ CONTIENT}\par\vspace{7pt}
  \noindent\begin{minipage}[t]{0.235\linewidth}\poptile{popblue}{Cours}{complet, illustré, pas à pas}\end{minipage}\hfill
  \begin{minipage}[t]{0.235\linewidth}\poptile{poporange}{Méthodes}{et exercices résolus}\end{minipage}\hfill
  \begin{minipage}[t]{0.235\linewidth}\poptile{poppink}{Exercices}{type examen + corrigés}\end{minipage}\hfill
  \begin{minipage}[t]{0.235\linewidth}\poptile{popgreen}{Fiche bilan}{l'essentiel à retenir}\end{minipage}\par}

% Sommaire
\newtcolorbox{sommaire}{enhanced,colback=white,colframe=popink,boxrule=2.2pt,arc=10pt,
  drop shadow=popink,left=18pt,right=18pt,top=13pt,bottom=13pt,before skip=6pt,after skip=20pt,
  before upper={\poppill{Sommaire}{poppurple}{white}\par\vspace{9pt}\popsemi\fontsize{10.6}{14.5}\selectfont\color{popink}}}

% Signature de fin de cours — poussée tout en bas de la dernière page
\newcommand{\findecours}{%
  \par\vfill\nopagebreak
  \begin{center}\popsquiggle{poporange}\end{center}\vspace{4pt}
  \signatureonbuch
}

%% FIGFIX-BEGIN v5 — correctifs globaux des figures (docs/onbuch-generation/tools/figfix)
\usepackage{adjustbox}\usepackage{etoolbox}\tcbuselibrary{raster}
% toute figure TikZ/pgfplots plus large que la ligne est réduite à la largeur disponible
\BeforeBeginEnvironment{tikzpicture}{\begin{adjustbox}{max width=\linewidth}}
\AfterEndEnvironment{tikzpicture}{\end{adjustbox}}
% légendes pgfplots lisibles par-dessus les courbes
\pgfplotsset{every axis/.append style={legend style={fill=white,fill opacity=0.92,text opacity=1,draw=black!15,inner sep=3pt}}}
% étiquettes d'arêtes (syntaxe edge["..."]) avec fond blanc
\tikzset{every edge quotes/.append style={fill=white,inner sep=1.2pt}}
%% FIGFIX-END
\begin{document}

\pagegarde{Oscillations électriques forcées en régime sinusoïdal}{Physique}{Tle D}{Module 3 : Électricité — évolutions temporelles des circuits électriques et électroniques}{P12}{5 h}{Cette leçon étudie le circuit RLC série soumis à une tension sinusoïdale : construction de la représentation de Fresnel, notions d'impédance et de déphasage, phénomène de résonance d'intensité, bande passante, facteur de qualité et puissance moyenne. Elle fournit les outils d'analyse des circuits électriques et électroniques de la vie courante et prépare directement aux épreuves du Baccalauréat C.
\vspace{4pt}
\begin{multicols}{2}\small
\begin{itemize}
\item À la fin de cette leçon, je sais établir l'équation différentielle reliant la tension u(t) aux bornes du dipôle RLC série à la charge q(t) ou à l'intensité i(t).
\item À la fin de cette leçon, je sais construire le diagramme de Fresnel associé aux tensions u\_R, u\_L, u\_C et u.
\item À la fin de cette leçon, je sais exprimer et calculer l'impédance Z d'un dipôle (R, L, C ou RLC) et le déphasage φ entre u et i.
\item À la fin de cette leçon, je sais tracer et exploiter la courbe de résonance d'intensité I = f(N) et déterminer la fréquence de résonance.
\item À la fin de cette leçon, je sais déterminer la bande passante à -3 dB et le facteur de qualité Q d'un circuit RLC.
\item À la fin de cette leçon, je sais exprimer et calculer la puissance moyenne P = U I cos φ et le facteur de puissance d'une installation.
\end{itemize}\end{multicols}}

\begin{sommaire}
\begin{multicols}{2}
\begin{enumerate}
\item Le circuit RLC série en régime sinusoïdal forcé
\item Représentation de Fresnel, impédance et déphasage
\item Résonance d'intensité
\item Bande passante à -3 dB et facteur de qualité
\item Puissance moyenne et facteur de puissance
\item Mesure du déphasage à l'oscilloscope et exploitation des courbes
\item Activité d'intégration
\item Méthodes et exercices résolus
\item Exercices
\item Corrigés des exercices
\item Fiche bilan
\end{enumerate}
\end{multicols}
\end{sommaire}

\begin{objectifs}
\begin{itemize}
\item À la fin de cette leçon, je sais établir l'équation différentielle reliant la tension u(t) aux bornes du dipôle RLC série à la charge q(t) ou à l'intensité i(t).
\item À la fin de cette leçon, je sais construire le diagramme de Fresnel associé aux tensions u\_R, u\_L, u\_C et u.
\item À la fin de cette leçon, je sais exprimer et calculer l'impédance Z d'un dipôle (R, L, C ou RLC) et le déphasage φ entre u et i.
\item À la fin de cette leçon, je sais tracer et exploiter la courbe de résonance d'intensité I = f(N) et déterminer la fréquence de résonance.
\item À la fin de cette leçon, je sais déterminer la bande passante à -3 dB et le facteur de qualité Q d'un circuit RLC.
\item À la fin de cette leçon, je sais exprimer et calculer la puissance moyenne P = U I cos φ et le facteur de puissance d'une installation.
\item À la fin de cette leçon, je sais mesurer un déphasage à l'oscilloscope à partir du décalage temporel de deux courbes.
\item À la fin de cette leçon, je sais expliquer le rôle de la résonance dans la sélection d'une station de radio.
\end{itemize}
\end{objectifs}

\begin{prerequis}
\begin{itemize}
\item Dipôles R, L, C en régime transitoire : lois u\_R = Ri, u\_L = L di/dt, i = C du\_C/dt (leçons P9 à P11 de Terminale)
\item Oscillations électriques libres du circuit LC et période propre T₀ = 2π√(LC)
\item Grandeurs sinusoïdales : valeur efficace, amplitude, période, fréquence, pulsation (Première)
\item Utilisation de l'oscilloscope : mesure de tensions et de périodes (Première)
\item Notions de vecteurs et de trigonométrie de base (cosinus, tangente)
\end{itemize}
\end{prerequis}

\newpage

\coursec{Le circuit RLC série en régime sinusoïdal forcé}

À Yaoundé comme à Douala, les coupures et fluctuations de la tension ENEO endommagent souvent téléviseurs et réfrigérateurs : les techniciens parlent alors de « mauvais \cle{cos φ} » et de « résonance » dans les installations. Pourquoi une radio capte-t-elle une seule station parmi des centaines d'ondes qui traversent l'air ? Pourquoi un condensateur branché dans un atelier de menuiserie permet-il de réduire la facture d'électricité ? Toutes ces questions trouvent leur réponse dans l'étude du circuit RLC série alimenté en tension sinusoïdale, cœur de cette leçon.

\courssub{Description du circuit et régime forcé}

Considérons un circuit série constitué de trois dipôles passifs :
\begin{itemize}
    \item un \cle{résistor} d'ohmique pur $R$ (en \si{\ohm}) ;
    \item une \cle{bobine} réelle modélisée par une inductance $L$ (en \si{\henry}) en série avec sa résistance interne $r$ (souvent négligée devant $R$) ;
    \item un \cle{condensateur} de capacité $C$ (en \si{\farad}).
\end{itemize}
L'ensemble est alimenté par un \cle{générateur de basse fréquence (GBF)} délivrant une tension sinusoïdale de la forme :
\[
u(t) = U_m \cos(\omega t + \varphi_u)
\]
où $U_m$ est l'amplitude (valeur maximale), $\omega = 2\pi f$ la pulsation (en \si{\radian\per\second}) et $\varphi_u$ la phase à l'origine.

\begin{popfigure}
\begin{tikzpicture}[european, thick, scale=1.0, every node/.style={font=\small}]
    % Nodes
    \coordinate (GND) at (0,0);
    \coordinate (TOP) at (0,3);
    \coordinate (B) at (3,3);
    \coordinate (C) at (5,3);
    \coordinate (D) at (7,3);
    \coordinate (G) at (3,0);
    \coordinate (F) at (5,0);
    \coordinate (E) at (7,0);
    
    % Generator (left vertical)
    \draw (GND) to[sV, l_={$u(t)=U_m\cos(\omega t+\varphi_u)$}, invert] (TOP);
    
    % Top rail
    \draw (TOP) -- (B) -- (C) -- (D);
    
    % Components on bottom rail (Receiver convention: current enters + terminal)
    % Resistor R (0 to 3)
    \draw (GND) to[R, l={$R$}, *-*] (G);
    % Inductor L (3 to 5)
    \draw (G) to[L, l={$L$}, *-*] (F);
    % Capacitor C (5 to 7)
    \draw (F) to[C, l={$C$}, *-*] (E);
    
    % Right vertical wire
    \draw (E) -- (D);
    
    % Current arrow on top rail (reference direction)
    \draw[->, popblue, thick] (1.5, 3.4) -- node[above] {$i(t)$} (2.5, 3.4);
    
    % Voltage arrows for each component (parallel, below components, receiver convention)
    \draw[<->, poporange] (GND) ++(0,-0.8) -- node[below] {$u_R$} (G) ++(0,-0.8);
    \draw[<->, popgreen]  (G) ++(0,-0.8) -- node[below] {$u_L$} (F) ++(0,-0.8);
    \draw[<->, poppurple] (F) ++(0,-0.8) -- node[below] {$u_C$} (E) ++(0,-0.8);
    
    % Global voltage across RLC (optional, generator already labeled)
    % \draw[<->, popdark] (GND) ++(0,-1.5) -- node[below] {$u(t)$} (E) ++(0,-1.5);
\end{tikzpicture}
\legende{Circuit RLC série alimenté par un GBF. Le sens du courant $i(t)$ (bleu) est choisi comme référence. Les tensions $u_R$, $u_L$, $u_C$ (flèches oranges, vertes, violettes) sont orientées dans le même sens que le courant (convention récepteur).}
\end{popfigure}

\begin{definition}[Régime sinusoïdal forcé]
On appelle \textbf{régime sinusoïdal forcé} l'état permanent atteint par un circuit linéaire alimenté par une source sinusoïdale après l'extinction des régimes transitoires. Le circuit oscille alors \emph{à la fréquence imposée par le générateur} ($f = \omega/2\pi$), et non à sa fréquence propre. Le courant $i(t)$ est sinusoïdal de même pulsation $\omega$ que la tension d'alimentation :
\[
i(t) = I_m \cos(\omega t + \varphi_i)
\]
où $I_m$ est l'amplitude du courant et $\varphi_i$ sa phase à l'origine. Le \cle{déphasage} entre la tension et le courant est $\varphi = \varphi_u - \varphi_i$.
\end{definition}

\begin{attention}[Valeurs maximales vs valeurs efficaces]
Au Bac, on manipule souvent les \textbf{valeurs efficaces} (notées $U$ et $I$ sans indice) définies par $U = U_m/\sqrt{2}$, $I = I_m/\sqrt{2}$. Les relations d'impédance et de puissance s'écrivent naturellement avec les valeurs efficaces. Ne confondez jamais $U_m$ (amplitude, crête) et $U$ (efficace, RMS) : une tension « \SI{220}{V} » du réseau ENEO est une valeur efficace ; son amplitude est $220\sqrt{2} \approx \SI{311}{V}$.
\end{attention}

\courssub{Équation différentielle du circuit RLC série}

Appliquons la \cle{loi d'additivité des tensions} (loi des mailles) dans le sens du courant $i(t)$ :
\[
u(t) = u_R(t) + u_L(t) + u_C(t)
\]
Les relations tension-courant pour chaque dipôle (convention récepteur) sont :
\begin{align*}
u_R(t) &= R\, i(t) \\
u_L(t) &= L\,\frac{di(t)}{dt} \quad (\text{on néglige } r) \\
u_C(t) &= \frac{1}{C}\int i(t)\,dt \quad (\text{charge } q(t) = \int i(t)\,dt)
\end{align*}

\begin{propriete}[Équation différentielle en charge $q$ et en courant $i$]
En posant $i(t) = \dfrac{dq}{dt}$, l'équation de la maille devient :
\[
L\,\frac{d^2q}{dt^2} + R\,\frac{dq}{dt} + \frac{1}{C}\,q = u(t) = U_m\cos(\omega t + \varphi_u)
\]
En dérivant par rapport au temps, on obtient l'équation en courant :
\[
L\,\frac{d^2i}{dt^2} + R\,\frac{di}{dt} + \frac{1}{C}\,i = \frac{du}{dt} = -\omega U_m\sin(\omega t + \varphi_u)
\]
Ces équations sont des \textbf{équations différentielles linéaires du second ordre à coefficients constants, non homogènes}. Leur solution générale est la somme d'une solution particulière (régime forcé) et de la solution de l'équation homogène (régime libre, transitoire).
\end{propriete}

\begin{experience}[Démonstration guidée : de la maille à l'équation en $q$]
\begin{enumerate}
    \item Écrire la loi des mailles : $u = u_R + u_L + u_C$.
    \item Exprimer chaque tension en fonction du courant $i$ : $u_R = Ri$, $u_L = L\,di/dt$, $u_C = q/C$ avec $i = dq/dt$.
    \item Substituer : $u = R\,\frac{dq}{dt} + L\,\frac{d^2q}{dt^2} + \frac{q}{C}$.
    \item Ranger par ordre de dérivation : $L\,q'' + R\,q' + \frac{1}{C}q = u(t)$.
    \item \textbf{Analyse dimensionnelle} : $[L\,q''] = \text{H}\cdot\text{C}\cdot\text{s}^{-2} = (\text{V}\cdot\text{s}\cdot\text{A}^{-1})\cdot\text{A}\cdot\text{s}^{-1}\cdot\text{s}^{-2} = \text{V}$. Tous les termes sont homogènes à une tension.
\end{enumerate}
\end{experience}

\courssub{Solution en régime permanent : amplitude et déphasage}

En régime forcé, le régime transitoire s'est éteint (amorti par $R$). Le courant est une sinusoïde de même pulsation $\omega$ :
\[
i(t) = I_m \cos(\omega t + \varphi_i)
\]
Le problème physique se réduit à déterminer \textbf{l'amplitude $I_m$} et \textbf{le déphasage $\varphi = \varphi_u - \varphi_i$} en fonction de $R$, $L$, $C$ et $\omega$. C'est l'objet de la section 2 (représentation de Fresnel et impédance complexe).

\courssub{Cas particuliers : dipôles R, L, C seuls}

Avant d'aborder le cas général RLC, comprenons le déphasage élémentaire pour chaque dipôle pris isolément, alimenté par $i(t) = I_m\cos(\omega t)$ (choix de l'origine des temps : $\varphi_i=0$).

\begin{methode}[Déterminer le déphasage pour un dipôle élémentaire]
\begin{enumerate}
    \item Écrire la relation tension-courant du dipôle ($u = Ri$, $u = L\,di/dt$, $u = q/C$).
    \item Exprimer $i(t)$ sous forme cosinusoïdale : $i(t) = I_m\cos(\omega t)$.
    \item Calculer $u(t)$ en dérivant ou en intégrant.
    \item Réécrire $u(t)$ sous la forme $U_m\cos(\omega t + \varphi)$ et identifier $\varphi$.
    \item Conclure : $\varphi > 0$ signifie que $u$ est \emph{en avance} sur $i$ ; $\varphi < 0$ signifie que $u$ est \emph{en retard} sur $i$.
\end{enumerate}
\end{methode}

\begin{exemplebox}[Résistance $R$ seule]
$u_R(t) = R\,i(t) = R I_m\cos(\omega t) = U_{R,m}\cos(\omega t)$.
$\Rightarrow$ \textbf{$u_R$ et $i$ sont en phase} ($\varphi_R = 0$).
\end{exemplebox}

\begin{exemplebox}[Bobine $L$ seule (idéale)]
$u_L(t) = L\,\frac{di}{dt} = L\,(-\omega I_m\sin(\omega t)) = \omega L I_m\cos(\omega t + \frac{\pi}{2})$.
$\Rightarrow$ \textbf{$u_L$ est en avance de $\pi/2$ sur $i$} ($\varphi_L = +\pi/2$).
\emph{Mnémotechnique :} dans une bobine, la tension « précède » le courant (L pour « Lead » en anglais).
\end{exemplebox}

\begin{exemplebox}[Condensateur $C$ seul]
$u_C(t) = \frac{1}{C}\int i\,dt = \frac{I_m}{\omega C}\sin(\omega t) = \frac{I_m}{\omega C}\cos(\omega t - \frac{\pi}{2})$.
$\Rightarrow$ \textbf{$u_C$ est en retard de $\pi/2$ sur $i$} ($\varphi_C = -\pi/2$).
\emph{Mnémotechnique :} dans un condensateur, le courant « précède » la tension (C pour « Lag » en anglais).
\end{exemplebox}

\begin{popfigure}
\begin{tikzpicture}
    \begin{axis}[
        width=\linewidth,
        height=7cm,
        axis lines=middle,
        xlabel={$t$ (ms)}, ylabel={Amplitude (V ou A)},
        domain=0:20, samples=400,
        xmin=0, xmax=20,
        ymin=-1.5, ymax=1.5,
        xtick={0,5,10,15,20},
        ytick={-1,-0.5,0,0.5,1},
        grid=major, grid style={popline},
        legend style={at={(0.98,0.98)}, anchor=north east, fill=popcream, draw=popdark}
    ]
    % u(t) = cos(omega t) with omega = 2*pi/10 = 0.628 rad/ms -> period 10 ms
    \addplot[popblue, thick, samples=400, domain=0:20] {cos(0.62831853*x)}; % u(t)
    \addlegendentry{$u(t) = U_m\cos(\omega t)$}
    % i(t) = cos(omega t - phi) with phi = pi/4 (45 deg)
    \addplot[poporange, thick, samples=400, domain=0:20] {cos(0.62831853*x - 0.785398)}; % i(t) lagging
    \addlegendentry{$i(t) = I_m\cos(\omega t - \pi/4)$}
    
    % Annotations
    \draw[popblue, <->] (axis cs:0,1.2) -- node[above, popblue] {$T = 10$ ms} (axis cs:10,1.2);
    \draw[popdark, <->] (axis cs:1.25, -1.3) -- node[below] {$\Delta t \approx 1,25$ ms} (axis cs:2.5, -1.3);
    \node[popdark, align=center] at (axis cs:10, -1.3) {$\varphi = \omega\Delta t = \frac{2\pi}{T}\Delta t = \frac{\pi}{4}$};
\end{axis}
\end{tikzpicture}
\legende{Tensions et courants sinusoïdaux déphasés. Ici $u$ (bleu) est en avance sur $i$ (orange) de $\varphi = \pi/4$ ($\Delta t = T/8$). Le déphasage se mesure sur les passages par zéro à pente positive.}
\end{popfigure}

\begin{savaistu}[Le réseau ENEO : 50 Hz, 220 V efficaces]
Au Cameroun, le réseau de distribution d'électricité (ENEO/SONATREL) délivre une tension alternative sinusoïdale de fréquence $f = \SI{50}{Hz}$ (pulsation $\omega = 100\pi \approx \SI{314}{rad.s^{-1}}$) et de valeur efficace $U = \SI{220}{V}$. L'amplitude est donc $U_m = 220\sqrt{2} \approx \SI{311}{V}$. C'est cette tension qui alimente nos foyers et qui, après redressement et filtrage, charge nos téléphones (MTN, Orange, CAMTEL) et fait tourner les moteurs des ateliers. Comprendre le régime forcé, c'est comprendre comment chaque appareil « tire » son courant avec un déphasage qui impacte la facture et la stabilité du réseau.
\end{savaistu}

\begin{exoresolu}[Vérification des déphasages élémentaires]
\textbf{Énoncé :} Une bobine d'inductance $L = \SI{100}{mH}$ est traversée par un courant $i(t) = \num{0,50}\cos(314 t)$ (A). Un condensateur $C = \SI{10}{\micro\farad}$ est traversé par le même courant. Déterminer les expressions des tensions $u_L(t)$ et $u_C(t)$ et leurs déphasages par rapport au courant.

\textbf{Solution détaillée :}
\begin{itemize}
    \item \textbf{Bobine :} $u_L = L\,di/dt = \num{0,100} \times \frac{d}{dt}[\num{0,50}\cos(314 t)] = \num{0,100} \times \num{0,50} \times (-314)\sin(314 t) = -\num{15,7}\sin(314 t)$.
    On réécrit : $-\num{15,7}\sin(314 t) = \num{15,7}\cos(314 t + \pi/2)$.
    Amplitude $U_{L,m} = \SI{15,7}{V}$. \textbf{Déphasage $\varphi_L = +\pi/2$ (avance)}.
    
    \item \textbf{Condensateur :} $u_C = \frac{1}{C}\int i\,dt = \frac{1}{\num{10e-6}}\int \num{0,50}\cos(314 t)\,dt = 10^5 \times \frac{\num{0,50}}{314}\sin(314 t) \approx \num{159}\sin(314 t)$.
    On réécrit : $\num{159}\sin(314 t) = \num{159}\cos(314 t - \pi/2)$.
    Amplitude $U_{C,m} \approx \SI{159}{V}$. \textbf{Déphasage $\varphi_C = -\pi/2$ (retard)}.
\end{itemize}
\textbf{Conclusion :} Les amplitudes dépendent de $\omega$ ($U_{L,m} \propto \omega$, $U_{C,m} \propto 1/\omega$), tandis que les déphasages sont constants ($\pm\pi/2$).
\end{exoresolu}

\begin{aretenir}[Bilan de la section 1]
\begin{itemize}
    \item Circuit RLC série : $u(t) = u_R + u_L + u_C$ avec $u_R=Ri$, $u_L=L\,di/dt$, $u_C=q/C$.
    \item Équation différentielle en charge : $L q'' + R q' + q/C = u(t)$.
    \item En régime forcé : $i(t) = I_m\cos(\omega t + \varphi_i)$ de même fréquence que $u(t)$.
    \item Cas particuliers :
    \begin{center}
    \begin{tabular}{|c|c|c|}\hline
    Dipôle & Relation $u(t)$ vs $i(t)$ & Déphasage $\varphi = \varphi_u - \varphi_i$ \\ \hline
    $R$ & $u_R = R i$ & $0$ (en phase) \\ \hline
    $L$ & $u_L = L\,di/dt$ & $+\pi/2$ ($u$ en avance) \\ \hline
    $C$ & $u_C = \frac{1}{C}\int i\,dt$ & $-\pi/2$ ($u$ en retard) \\ \hline
    \end{tabular}
    \end{center}
    \item Attention : distinguer amplitude ($U_m, I_m$) et valeur efficace ($U = U_m/\sqrt{2}, I = I_m/\sqrt{2}$).
\end{itemize}
\end{aretenir}

\coursec{Représentation de Fresnel, impédance et déphasage}

Dans la section précédente, nous avons établi l'équation différentielle du circuit RLC série soumis à une tension sinusoïdale \(u(t)=U\sqrt{2}\cos(\omega t)\). Nous cherchons maintenant une méthode graphique et algébrique puissante pour déterminer l'amplitude \(I\) du courant et son déphasage \(\varphi\) par rapport à la tension. C'est tout l'objet de la \emph{représentation de Fresnel}, qui transforme les équations différentielles en géométrie vectorielle dans le plan complexe.

\courssub{Principes de la représentation de Fresnel}

\begin{definition}[Vecteur de Fresnel associé à une grandeur sinusoïdale]
À toute grandeur sinusoïdale de la forme \(x(t)=X\sqrt{2}\cos(\omega t+\varphi_x)\) (où \(X\) est la \emph{valeur efficace}), on associe un \emph{vecteur de Fresnel} \(\vec{X}\) de norme \(X\) (valeur efficace) qui tourne dans le sens trigonométrique à la pulsation \(\omega\). Sa projection sur l'axe horizontal donne la valeur instantanée \(x(t)\). L'angle que ce vecteur fait avec l'axe horizontal à l'instant \(t=0\) est la \emph{phase} \(\varphi_x\).
\end{definition}

\begin{aretenir}[Choix de la référence des phases]
Dans un circuit \emph{série}, le courant \(i(t)\) est \emph{commun} à tous les dipôles. Il est donc naturel de le prendre comme référence des phases : on place le vecteur de Fresnel \(\vec{I}\) sur l'axe horizontal (angle nul). Ainsi \(i(t)=I\sqrt{2}\cos(\omega t)\).
\end{aretenir}

\courssub{Construction du diagramme de Fresnel pour le circuit RLC série}

Rappelons les relations tension-courant pour chaque dipôle en régime sinusoïdal (valeurs efficaces) :
\begin{itemize}
  \item Résistance \(R\) : \(u_R(t)=Ri(t)\) → \(\vec{U}_R = R\vec{I}\) (en phase avec \(\vec{I}\)).
  \item Bobine \(L\) : \(u_L(t)=L\frac{di}{dt}\) → \(\vec{U}_L = L\omega\vec{I}\) (en \emph{avance} de \(\pi/2\) sur \(\vec{I}\)).
  \item Condensateur \(C\) : \(u_C(t)=\frac{1}{C}\int i\,dt\) → \(\vec{U}_C = \frac{1}{C\omega}\vec{I}\) (en \emph{retard} de \(\pi/2\) sur \(\vec{I}\)).
\end{itemize}

La loi des mailles s'écrit \(\vec{U} = \vec{U}_R + \vec{U}_L + \vec{U}_C\). Comme \(\vec{U}_L\) et \(\vec{U}_C\) sont opposées, leur somme vectorielle est un vecteur vertical de norme \(|U_L-U_C| = I|L\omega - \frac{1}{C\omega}|\).

\begin{popfigure}
\begin{tikzpicture}[>=Stealth, scale=0.9]
% Axes
\draw[->, thin, popmuted] (-4,0) -- (4,0) node[right] {$x$};
\draw[->, thin, popmuted] (0,-3.5) -- (0,3.5) node[above] {$y$};
% Vecteur I (reference)
\draw[->, thick, popblue] (0,0) -- (2.5,0) node[midway, below] {$\vec{I}$};
% Vecteur U_R (colinéaire à I)
\draw[->, thick, poporange] (0,0) -- (2.5,0) node[midway, above] {$\vec{U}_R$};
% Vecteur U_L (vertical up)
\draw[->, thick, popgreen] (0,0) -- (0,2.5) node[midway, left] {$\vec{U}_L$};
% Vecteur U_C (vertical down)
\draw[->, thick, poppurple] (0,0) -- (0,-2.5) node[midway, right] {$\vec{U}_C$};
% Resultant U = U_R + (U_L - U_C)
\draw[->, very thick, popdark] (0,0) -- (2.5,1.2) node[midway, above right] {$\vec{U}$};
% Angle phi
\draw[popdark] (0.5,0) arc (0:25.6:0.5) node[midway, right, xshift=2pt] {$\varphi$};
% Dashed lines for construction
\draw[dashed, thin, popmuted] (2.5,0) -- (2.5,1.2);
\draw[dashed, thin, popmuted] (0,1.2) -- (2.5,1.2);
% Labels for magnitudes
\node[popblue] at (1.25,-0.4) {$I$};
\node[poporange] at (1.25,0.3) {$U_R=RI$};
\node[popgreen] at (-0.4,1.25) {$U_L=L\omega I$};
\node[poppurple] at (0.4,-1.25) {$U_C=\frac{I}{C\omega}$};
\node[popdark] at (1.5,0.8) {$U$};
% Case Lω > 1/(Cω)
\node[align=center, popdark] at (3.5,-2) {Cas $L\omega > \frac{1}{C\omega}$\\$\varphi>0$ (inductif)};
\end{tikzpicture}
\legende{Diagramme de Fresnel des tensions dans un circuit RLC série (cas inductif : $L\omega > 1/(C\omega)$). Le vecteur $\vec{U}$ est la somme vectorielle de $\vec{U}_R$, $\vec{U}_L$ et $\vec{U}_C$. L'angle $\varphi$ est le déphasage de la tension sur le courant.}
\end{popfigure}

\begin{methode}[Construire un diagramme de Fresnel pour un circuit RLC série]
\begin{enumerate}
  \item Placer le vecteur $\vec{I}$ horizontalement vers la droite (référence de phase).
  \item Tracer $\vec{U}_R = R\vec{I}$ colinéaire à $\vec{I}$ (même direction).
  \item Tracer $\vec{U}_L = L\omega\vec{I}$ verticalement vers le haut (avance de $\pi/2$).
  \item Tracer $\vec{U}_C = \frac{1}{C\omega}\vec{I}$ verticalement vers le bas (retard de $\pi/2$).
  \item Effectuer la somme vectorielle : $\vec{U} = \vec{U}_R + \vec{U}_L + \vec{U}_C$. Comme $\vec{U}_L$ et $\vec{U}_C$ sont opposées, on trace d'abord $\vec{U}_L - \vec{U}_C$ (vecteur vertical), puis on ajoute $\vec{U}_R$ horizontalement.
  \item La norme du résultant donne la tension efficace $U$ ; l'angle $\varphi$ entre $\vec{U}$ et $\vec{I}$ (sens trigonométrique) est le déphasage tension/courant.
\end{enumerate}
\end{methode}

\courssub{Expression de l'impédance et du déphasage}

Le triangle rectangle formé par $\vec{U}_R$, $\vec{U}_L-\vec{U}_C$ et $\vec{U}$ (voir figure ci-dessus) permet d'appliquer le théorème de Pythagore et la définition de la tangente.

\begin{propriete}[Impédance et déphasage d'un circuit RLC série — Démonstration exigible au Bac]
L'\emph{impédance} $Z$ du circuit RLC série est définie comme le quotient de la tension efficace aux bornes du circuit par l'intensité efficace :
\[
Z = \frac{U}{I}.
\]
D'après le diagramme de Fresnel, on a :
\[
U^2 = U_R^2 + (U_L-U_C)^2 = (RI)^2 + \left(L\omega I - \frac{I}{C\omega}\right)^2.
\]
En factorisant $I^2$, on obtient l'expression de l'impédance :
\[
\boxed{Z = \sqrt{R^2 + \left(L\omega - \frac{1}{C\omega}\right)^2}} \qquad (\text{en ohms, }\Omega).
\]
Le déphasage $\varphi$ entre la tension $u(t)$ et le courant $i(t)$ (tension en avance sur le courant) satisfait :
\[
\tan\varphi = \frac{U_L-U_C}{U_R} = \frac{L\omega - \frac{1}{C\omega}}{R},
\qquad
\cos\varphi = \frac{R}{Z},
\qquad
\sin\varphi = \frac{L\omega - \frac{1}{C\omega}}{Z}.
\]
\end{propriete}

\begin{aretenir}[Signe de $\varphi$ et nature du circuit]
\begin{itemize}
  \item Si $L\omega > \frac{1}{C\omega}$ : $\varphi > 0$, le circuit est \emph{inductif} (la tension est en avance sur le courant).
  \item Si $L\omega < \frac{1}{C\omega}$ : $\varphi < 0$, le circuit est \emph{capacitif} (la tension est en retard sur le courant).
  \item Si $L\omega = \frac{1}{C\omega}$ : $\varphi = 0$, le circuit est \emph{résistif pur} (tension et courant en phase) — c'est la \emph{résonance d'intensité} (voir section 3).
\end{itemize}
\end{aretenir}

\begin{popfigure}
\begin{tikzpicture}[>=Stealth, scale=1]
% Triangle des impédances
\draw[->, thick, poporange] (0,0) -- (3,0) node[midway, below] {$R$};
\draw[->, thick, popgreen] (3,0) -- (3,2) node[midway, right] {$X = L\omega - \frac{1}{C\omega}$};
\draw[->, very thick, popdark] (0,0) -- (3,2) node[midway, above left] {$Z$};
% Angle phi
\draw[popdark] (0.5,0) arc (0:33.7:0.5) node[midway, right, xshift=2pt] {$\varphi$};
% Right angle mark
\draw[thin, popmuted] (2.85,0) -- (2.85,0.15) -- (3,0.15);
% Labels
\node[popdark] at (1.5,1.2) {$Z = \sqrt{R^2+X^2}$};
\node[popdark] at (3.3,1) {$\tan\varphi = \frac{X}{R}$};
\end{tikzpicture}
\legende{Triangle des impédances (ou triangle de Fresnel). $X = L\omega - 1/(C\omega)$ est la \emph{réactance} totale. L'hypoténuse est l'impédance $Z$, l'angle $\varphi$ est le déphasage.}
\end{popfigure}

\courssub{Tableau récapitulatif des impédances des dipôles élémentaires}

\begin{center}
\begin{tabularx}{\linewidth}{|l|X|X|X|}\hline
\rowcolor{popblueL}
\textbf{Dipôle} & \textbf{Impédance complexe $Z$} & \textbf{Module $|Z|$} & \textbf{Phase $\varphi$ (tension/courant)} \\ \hline
Résistance $R$ & $Z_R = R$ & $R$ & $0$ (en phase) \\ \hline
Bobine $L$ (idéal) & $Z_L = jL\omega$ & $L\omega$ & $+\pi/2$ (avance) \\ \hline
Condensateur $C$ (idéal) & $Z_C = \frac{1}{jC\omega} = -j\frac{1}{C\omega}$ & $\frac{1}{C\omega}$ & $-\pi/2$ (retard) \\ \hline
\end{tabularx}
\end{center}

\begin{attention}[Piège classique : additionner les tensions efficaces !]
\emph{Ne faites jamais} $U = U_R + U_L + U_C$ (somme algébrique des valeurs efficaces). Les tensions aux bornes de $R$, $L$ et $C$ ne sont \emph{pas en phase} : $U_R$ est en phase avec $i$, $U_L$ est en avance de $\pi/2$, $U_C$ en retard de $\pi/2$. La tension totale $U$ est la \emph{somme vectorielle} (ou somme phasorielle) : $U = \sqrt{U_R^2 + (U_L-U_C)^2}$. L'erreur conduit à surestimer $U$ et à fausser l'impédance.
\end{attention}

\courssub{Exemple chiffré détaillé}

\begin{exemplebox}[Calcul de l'impédance et du déphasage]
Soit un circuit RLC série avec $R = \SI{100}{\ohm}$, $L = \SI{0.5}{\henry}$, $C = \SI{10}{\micro\farad} = \SI{10e-6}{\farad}$, alimenté sous une tension efficace $U = \SI{230}{\volt}$ à la fréquence $f = \SI{50}{\hertz}$ ($\omega = 2\pi f \approx \SI{314.16}{\radian\per\second}$).

\begin{enumerate}
  \item Calculer les réactances :
  \[
  X_L = L\omega = 0.5 \times 314.16 \approx \SI{157.1}{\ohm},
  \qquad
  X_C = \frac{1}{C\omega} = \frac{1}{10^{-5}\times 314.16} \approx \SI{318.3}{\ohm}.
  \]
  \item Réactance totale :
  \[
  X = X_L - X_C \approx 157.1 - 318.3 = \SI{-161.2}{\ohm}.
  \]
  \item Impédance :
  \[
  Z = \sqrt{R^2 + X^2} = \sqrt{100^2 + (-161.2)^2} \approx \sqrt{10000 + 25985} \approx \sqrt{35985} \approx \SI{189.7}{\ohm}.
  \]
  \item Déphasage :
  \[
  \tan\varphi = \frac{X}{R} = \frac{-161.2}{100} = -1.612 \quad\Rightarrow\quad \varphi \approx -58.2^\circ \text{ (ou } -1.015\,\text{rad)}.
  \]
  \item Intensité efficace :
  \[
  I = \frac{U}{Z} = \frac{230}{189.7} \approx \SI{1.21}{\ampere}.
  \]
  \item Nature du circuit : $X<0$ ($\varphi<0$) → \emph{circuit capacitif} (la tension est en retard sur le courant).
\end{enumerate}
\end{exemplebox}

\courssub{Analogie formelle avec les oscillateurs mécaniques forcés (complément Bac)}

\begin{savaistu}[Analogie mécanique-électrique]
L'équation différentielle du circuit RLC série,
\[
L\frac{d^2q}{dt^2} + R\frac{dq}{dt} + \frac{q}{C} = u(t),
\]
est \emph{mathématiquement identique} à celle d'un oscillateur mécanique amorti forcé :
\[
m\ddot{x} + b\dot{x} + kx = F(t),
\]
avec les correspondances :
\begin{center}
\begin{tabularx}{\linewidth}{c|X}
\textbf{Électrique} & \textbf{Mécanique} \\ \hline
Charge $q$ & Déplacement $x$ \\
Courant $i=\dot{q}$ & Vitesse $v=\dot{x}$ \\
Inductance $L$ & Masse $m$ \\
Résistance $R$ & Coefficient d'amortissement $b$ \\
Inverse capacité $1/C$ & Raideur $k$ \\
Tension $u(t)$ & Force extérieure $F(t)$
\end{tabularx}
\end{center}
L'impédance $Z(\omega)$ joue le rôle de l'\emph{impédance mécanique} $Z_m(\omega) = \sqrt{b^2 + (m\omega - k/\omega)^2}$, et la résonance d'intensité correspond à la résonance d'amplitude (vitesse) mécanique. Augustin Fresnel (1788–1827), qui a formalisé cette représentation vectorielle pour la lumière, a aussi conçu les lentilles à échelons des phares (dont celui de Kribi au Cameroun), sauvant d'innombrables navires.
\end{savaistu}

\begin{exoresolu}[Exploitation d'un diagramme de Fresnel]
\textbf{Énoncé :} On observe à l'oscilloscope les courbes de la tension $u(t)$ aux bornes d'un circuit RLC série et du courant $i(t)$ (via une résistance shunt de valeur connue). La période est $T=\SI{20}{\milli\second}$. La tension a une amplitude de $\SI{10}{\volt}$, le courant une amplitude de $\SI{2}{\ampere}$. Le pic de tension précède le pic de courant de $\Delta t = \SI{2}{\milli\second}$. Déterminer $Z$, $\varphi$, et dire si le circuit est inductif ou capacitif.

\tcblower

\textbf{Solution détaillée :}
\begin{enumerate}
  \item Pulsation : $\omega = \frac{2\pi}{T} = \frac{2\pi}{0.020} = 100\pi \approx \SI{314.16}{\radian\per\second}$.
  \item Valeurs efficaces : $U = \frac{10}{\sqrt{2}} \approx \SI{7.07}{\volt}$, $I = \frac{2}{\sqrt{2}} \approx \SI{1.41}{\ampere}$.
  \item Impédance : $Z = \frac{U}{I} = \frac{7.07}{1.41} = \SI{5.00}{\ohm}$.
  \item Déphasage : $\varphi = \omega \Delta t = 314.16 \times 0.002 = \SI{0.628}{\radian} \approx 36.0^\circ$.
  \item Comme $\varphi > 0$ (tension en avance sur le courant), le circuit est \emph{inductif}.
  \item On en déduit $R = Z\cos\varphi = 5.00 \times \cos 36^\circ \approx \SI{4.05}{\ohm}$ et $X = Z\sin\varphi = 5.00 \times \sin 36^\circ \approx \SI{2.94}{\ohm}$.
\end{enumerate}
\end{exoresolu}

\begin{exercice}[Entraînement : diagramme de Fresnel et impédance]
Un circuit RLC série est alimenté sous $U=\SI{12}{\volt}$ (efficace) à $f=\SI{1}{\kilo\hertz}$. On mesure $I=\SI{0.1}{\ampere}$ et le déphasage $\varphi = -30^\circ$ (tension en retard). Calculer $Z$, $R$, $X = L\omega - 1/(C\omega)$. En déduire la nature du circuit.
\end{exercice}

\begin{corrige}[Exercice 1]
$Z = U/I = 12/0.1 = \SI{120}{\ohm}$. $\varphi = -30^\circ$ → circuit capacitif. $R = Z\cos\varphi = 120\cos(-30^\circ) = 120\times\sqrt{3}/2 \approx \SI{103.9}{\ohm}$. $X = Z\sin\varphi = 120\sin(-30^\circ) = \SI{-60}{\ohm}$. Donc $L\omega - 1/(C\omega) = \SI{-60}{\ohm}$.
\end{corrige}

\coursec{Résonance d'intensité}

\courssub{Transition depuis la représentation de Fresnel}

Dans la section précédente, nous avons établi que l'impédance d'un circuit RLC série alimenté par une tension sinusoïdale $u(t) = U\sqrt{2}\cos(\omega t)$ s'écrit :
\[
Z = \sqrt{R^2 + \left(L\omega - \frac{1}{C\omega}\right)^2}
\]
et que le déphasage $\varphi$ entre la tension $u$ et l'intensité $i$ vérifie :
\[
\tan\varphi = \frac{L\omega - \frac{1}{C\omega}}{R}.
\]
L'intensité efficace dans le circuit est donnée par la loi d'Ohm généralisée :
\[
I = \frac{U}{Z} = \frac{U}{\sqrt{R^2 + \left(L\omega - \frac{1}{C\omega}\right)^2}}.
\]
Cette expression montre que $I$ dépend fortement de la pulsation $\omega$ (ou de la fréquence $N = \omega/2\pi$) imposée par le générateur de fonctions (GBF). Nous allons maintenant étudier comment varie $I$ lorsque l'on fait varier $\omega$ tout en maintenant $U$, $R$, $L$ et $C$ constants. C'est l'étude de la \textbf{réponse en intensité} du circuit.

\courssub{Étude de la réponse en intensité $I = f(N)$}

\textbf{Expression de l'intensité en fonction de $\omega$.}\par
Rappelons la relation fondamentale :
\begin{equation}
I(\omega) = \frac{U}{\sqrt{R^2 + \left(L\omega - \dfrac{1}{C\omega}\right)^2}}. \tag{1}
\end{equation}
Le dénominateur est l'impédance $Z(\omega)$. Comme $U$ est constant, \emph{l'intensité $I$ est maximale lorsque l'impédance $Z$ est minimale}.

\textbf{Recherche du minimum de $Z$.}\par
Le terme $R^2$ est constant. Le terme variable est $\left(L\omega - \frac{1}{C\omega}\right)^2$, qui est toujours positif ou nul. Il devient nul si et seulement si :
\[
L\omega - \frac{1}{C\omega} = 0 \quad \Longleftrightarrow \quad L\omega = \frac{1}{C\omega}.
\]
Cette condition définit la \textbf{résonance d'intensité}.

\begin{definition}[Résonance d'intensité]
On appelle \cle{résonance d'intensité} le phénomène pour lequel l'intensité efficace $I$ traversant un circuit RLC série passe par un maximum lorsque la fréquence $N$ (ou la pulsation $\omega$) du générateur sinusoïdal devient égale à la \cle{fréquence propre} $N_0$ (ou pulsation propre $\omega_0$) du circuit.
\end{definition}

\textbf{Condition de résonance : pulsation propre $\omega_0$ et fréquence propre $N_0$.}\par
L'égalité $L\omega_0 = \frac{1}{C\omega_0}$ donne :
\[
\omega_0^2 = \frac{1}{LC} \quad \Longrightarrow \quad \omega_0 = \frac{1}{\sqrt{LC}} \quad (\omega_0 > 0).
\]
La fréquence de résonance (fréquence propre) est donc :
\[
N_0 = \frac{\omega_0}{2\pi} = \frac{1}{2\pi\sqrt{LC}}.
\]
\begin{propriete}[Fréquence propre du circuit RLC série]
La pulsation propre $\omega_0$ et la fréquence propre $N_0$ d'un circuit RLC série ne dépendent \textbf{que} de l'inductance $L$ et de la capacité $C$ :
\[
\boxed{\omega_0 = \frac{1}{\sqrt{LC}} \qquad ; \qquad N_0 = \frac{1}{2\pi\sqrt{LC}}}.
\]
\textit{Analyse dimensionnelle :} $[L] = \text{H} = \text{kg}\cdot\text{m}^2\cdot\text{s}^{-2}\cdot\text{A}^{-2}$, $[C] = \text{F} = \text{kg}^{-1}\cdot\text{m}^{-2}\cdot\text{s}^4\cdot\text{A}^2$. Donc $[LC] = \text{s}^2$, $[\sqrt{LC}] = \text{s}$, $[1/\sqrt{LC}] = \text{s}^{-1}$, cohérent avec une pulsation (rad$\cdot$s$^{-1}$).
\end{propriete}

\begin{savaistu}[Origine du terme « fréquence propre »]
Cette fréquence $\omega_0 = 1/\sqrt{LC}$ est exactement celle des oscillations \emph{libres} amorties du même circuit RLC (étudiées en Première ou en début de Terminale). Si l'on charge le condensateur puis qu'on laisse le circuit évoluer sans générateur, il oscille à cette fréquence (amortie par $R$). Le circuit « préfère » cette fréquence : quand le générateur l'excite à $\omega_0$, il répond avec une amplitude maximale. C'est le même principe qu'un enfant sur une balançoire : pour l'amplitude maximale, il faut pousser à la fréquence propre de la balançoire.
\end{savaistu}

\courssub{État du circuit à la résonance ($\omega = \omega_0$)}

Lorsque le générateur est réglé sur la fréquence de résonance $N_0$ (pulsation $\omega_0$), les réactances s'annulent mutuellement : $X_L = L\omega_0 = \frac{1}{C\omega_0} = X_C$.

\begin{propriete}[Caractéristiques du circuit à la résonance (démontrée)]
À la résonance ($\omega = \omega_0$) :
\begin{enumerate}
    \item \textbf{Impédance minimale :} $Z_{\min} = R$. Le circuit se comporte comme une résistance pure $R$.
    \item \textbf{Intensité maximale :} $I_{\max} = \dfrac{U}{R}$.
    \item \textbf{Déphasage nul :} $\varphi = 0$. La tension $u(t)$ et l'intensité $i(t)$ sont \textbf{en phase}. Le facteur de puissance $\cos\varphi = 1$.
    \item \textbf{Tensions aux bornes de L et C :} $U_L = I_{\max} L\omega_0 = \dfrac{U}{R} L\omega_0$ et $U_C = \dfrac{I_{\max}}{C\omega_0} = \dfrac{U}{R} \frac{1}{C\omega_0}$. Comme $L\omega_0 = 1/(C\omega_0)$, on a $U_L = U_C$. Elles sont en opposition de phase (déphasage $\pi$), leur somme instantanée est nulle : $u_L(t) + u_C(t) = 0$. La tension du générateur $u(t)$ se retrouve intégralement aux bornes de la résistance : $u(t) = u_R(t) = R\,i(t)$.
\end{enumerate}
\end{propriete}

\begin{attention}[Piège majeur : Les tensions $U_L$ et $U_C$ ne sont pas nulles !]
Une erreur très fréquente au Bac consiste à croire que « à la résonance, $X_L = X_C$ donc $U_L = U_C = 0$ ».
\textbf{FAUX.} À la résonance, les \emph{réactances} sont égales ($X_L = X_C$), mais les \emph{tensions} $U_L = I X_L$ et $U_C = I X_C$ sont égales en \textbf{valeur absolue} et \textbf{opposées en phase}. Leur somme vectorielle (ou algébrique instantanée) est nulle, mais chacune peut atteindre une valeur bien supérieure à la tension du générateur $U$. C'est le phénomène de \textbf{surtension}, quantifié par le facteur de qualité $Q$ (section 4) : $U_L = U_C = Q \times U$.
\end{attention}

\courssub{Allure de la courbe de résonance $I = f(N)$ et influence de $R$}

La courbe représentant l'intensité efficace $I$ en fonction de la fréquence $N$ du générateur (tension $U$ constante) présente un pic en $N_0$.

\begin{itemize}
    \item Pour $N \ll N_0$ : $C\omega$ est petit, $1/(C\omega)$ est grand $\Rightarrow$ le circuit est \textbf{capacitif} ($Z \approx 1/(C\omega)$), $I$ est faible.
    \item Pour $N \gg N_0$ : $L\omega$ est grand $\Rightarrow$ le circuit est \textbf{inductif} ($Z \approx L\omega$), $I$ est faible.
    \item Pour $N = N_0$ : $Z = R$ (minimum), $I = I_{\max} = U/R$.
\end{itemize}

L'\textbf{acuité} du pic (sa « pointe ») dépend de la résistance $R$ :
\begin{itemize}
    \item \textbf{$R$ faible :} $I_{\max} = U/R$ est grand. Le terme $R^2$ dans $Z^2$ est petit, donc la variation de $(L\omega - 1/C\omega)^2$ domine vite. Le pic est \textbf{aigu}, \textbf{étroit} : le circuit est très sélectif en fréquence.
    \item \textbf{$R$ grand :} $I_{\max} = U/R$ est petit. Le terme $R^2$ domine dans $Z^2$, masquant la variation de la réactance. Le pic est \textbf{émoussé}, \textbf{large} : le circuit est peu sélectif.
\end{itemize}

\begin{popfigure}
\begin{tikzpicture}
\begin{axis}[
    width=\linewidth,
    height=0.55\linewidth,
    axis lines=middle,
    xlabel={$N$ (Hz)}, ylabel={$I$ (A)},
    xmin=0, xmax=25000,
    ymin=0, ymax=0.3,
    xtick={15915.5}, xticklabels={$N_0$},
    ytick={0.25, 0.0833}, yticklabels={$I_{\max}^{(1)}$, $I_{\max}^{(2)}$},
    tick label style={font=\small},
    label style={font=\small},
    clip=false,
    domain=100:25000,
    samples=400,
    legend style={at={(0.98,0.98)}, anchor=north east, font=\small, fill=white, draw=popmuted}
]
\addplot[popcurve, thick, popblue] ({x}, {5/sqrt(max(0,20^2+(1e-3*2*pi*x-1/(100e-9*2*pi*x))^2))});
\addlegendentry{$R = 20\,\Omega$ (pic aigu)}

\addplot[popcurve, thick, poporange] ({x}, {5/sqrt(max(0,60^2+(1e-3*2*pi*x-1/(100e-9*2*pi*x))^2))});
\addlegendentry{$R = 60\,\Omega$ (pic émoussé)}

% Ligne verticale pointillée à N0
\draw[popmuted, dashed, thick] (axis cs:15915.5,0) -- (axis cs:15915.5,0.28) node[above, font=\small, popdark] {$N_0$};
% Flèches indiquant la largeur (qualitatif)
\draw[<->, popgreen, thick] (axis cs:14000,0.12) -- (axis cs:17800,0.12) node[midway, above, font=\small, popgreen] {Large à $-3$ dB};
\end{axis}
\end{tikzpicture}
\legende{Courbes de résonance $I = f(N)$ pour $U = 5\,\text{V}$, $L = 1\,\text{mH}$, $C = 100\,\text{nF}$ ($N_0 \approx 15,9\,\text{kHz}$). En bleu : $R = 20\,\Omega$ ($I_{\max} = 0,25\,\text{A}$, pic aigu). En orange : $R = 60\,\Omega$ ($I_{\max} \approx 0,083\,\text{A}$, pic large). La largeur à $-3\,\text{dB}$ (section 4) est plus grande pour $R$ grand.}
\end{popfigure}

\courssub{Exemple chiffré détaillé}

\begin{exemplebox}[Calcul de la fréquence de résonance et de l'intensité maximale]
On considère un circuit RLC série alimenté par un GBF de tension efficace constante $U = 5,0\,\text{V}$.
Les composants sont : $R = 20\,\Omega$, $L = 1,0\,\text{mH} = 1,0 \times 10^{-3}\,\text{H}$, $C = 100\,\text{nF} = 1,0 \times 10^{-7}\,\text{F}$.

\textbf{1. Fréquence de résonance $N_0$.}
\[
N_0 = \frac{1}{2\pi\sqrt{LC}} = \frac{1}{2\pi\sqrt{(1,0 \times 10^{-3}) \times (1,0 \times 10^{-7})}} = \frac{1}{2\pi \times 10^{-5}} \approx \frac{10^5}{6,283} \approx 15\,915\,\text{Hz} \approx \textbf{15,9 kHz}.
\]
Pulsation de résonance : $\omega_0 = 2\pi N_0 \approx 1,00 \times 10^5\,\text{rad}\cdot\text{s}^{-1}$.

\textbf{2. Intensité maximale $I_{\max}$.}
À la résonance, $Z = R = 20\,\Omega$.
\[
I_{\max} = \frac{U}{R} = \frac{5,0}{20} = \textbf{0,25 A} = 250\,\text{mA}.
\]

\textbf{3. Tensions aux bornes de L et C à la résonance (surtension).}
Réactance à la résonance : $X_L = X_C = L\omega_0 = 1,0 \times 10^{-3} \times 1,00 \times 10^5 = 100\,\Omega$.
\[
U_L = U_C = I_{\max} \times X_L = 0,25 \times 100 = \textbf{25 V}.
\]
On observe que $U_L = U_C = 25\,\text{V}$ est \textbf{5 fois plus grande} que la tension du générateur $U = 5\,\text{V}$. Le facteur de qualité est ici $Q = \frac{U_L}{U} = 5$ (ou $Q = \frac{L\omega_0}{R} = \frac{100}{20} = 5$). Ce phénomène est exploité dans les circuits d'accord (radio) et les filtres.
\end{exemplebox}

\courssub{Application concrète : Le circuit d'accord d'un poste radio}

C'est l'application historique et pédagogique par excellence de la résonance d'intensité.

\begin{popfigure}
\begin{tikzpicture}[scale=1.1, transform shape, european]
% Antenne
\draw[popblue, thick, decorate, decoration={snake, amplitude=1.5pt, segment length=4pt}] (0,2) -- (0,0.5) node[above left, font=\small] {Antenne};
\draw[popblue, thick] (0,0.5) -- (0,0) node[ground]{};
% Bobine (self)
\draw (1.5,1.5) to[L=$L$ (bobine), l_={$L$}, *-*] (1.5,0);
% Condensateur variable
\draw (1.5,1.5) to[C=$C$ (variable), l^={$C_v$}, *-] (3.5,1.5);
% Résistance (charge / haut-parleur)
\draw (3.5,1.5) to[R=$R$ (charge), l^={$R_{eq}$}] (3.5,0);
% Masse
\draw (0,0) -- (3.5,0);
\node[ground] at (3.5,0) {};
% Flèche réglage C
\draw[poporange, thick, ->] (3.5,1.8) arc (90:-90:0.3) node[right, font=\small, poporange] {Réglage $C_v$};
% Légendes
\node[align=center, font=\small, popdark] at (1.75, -1.2) {Circuit d'accord série : $N_0 = \frac{1}{2\pi\sqrt{LC_v}}$};
\node[align=center, font=\small, popdark] at (1.75, -1.7) {On tourne le bouton $\to$ $C_v$ varie $\to$ $N_0$ varie $\to$ on sélectionne la station};
\end{tikzpicture}
\legende{Schéma simplifié d'un circuit d'accord de récepteur radio (antenne + bobine $L$ + condensateur variable $C_v$ + charge $R$). En faisant varier $C_v$, on fait varier la fréquence propre $N_0$ pour la caler sur la fréquence de la station désirée (ex: CRTV Radio à 94,0 MHz $\to$ circuit oscillateur accordé sur cette fréquence après mélangeur).}
\end{popfigure}

\begin{savaistu}[Le bouton de la radio et la résonance]
Quand tu tournes le bouton de syntonisation d'un poste radio (ou que tu fais une recherche automatique sur une radio de voiture MTN/Orange/Camtel), tu modifies physiquement la capacité $C$ d'un condensateur variable (ou une varicap pilotée en tension). Cela déplace la fréquence propre $N_0 = 1/(2\pi\sqrt{LC})$ du circuit d'accord.
\begin{itemize}
    \item Si $N_0$ coïncide avec la fréquence porteuse de la station (ex. $94,0\,\text{MHz}$ pour CRTV Radio à Yaoundé), le circuit entre en résonance.
    \item L'intensité $I$ dans le circuit devient maximale pour cette fréquence, alors qu'elle reste très faible pour les fréquences des autres stations (grâce à la sélectivité, liée au facteur de qualité $Q$).
    \item Le signal utile est ainsi amplifié naturellement par la résonance, avant d'être démodulé pour retrouver le son.
\end{itemize}
C'est la physique pure qui permet de choisir « CRTV » plutôt que « Radio Balafon » ou « RFI » sur la même antenne !
\end{savaistu}

\courssub{Méthode : Exploiter une courbe de résonance expérimentale}

\begin{methode}[Déterminer $L$ ou $C$ à partir de la courbe $I = f(N)$]
\begin{enumerate}
    \item \textbf{Repérer le maximum} de la courbe $I(N)$. L'abscisse de ce maximum donne la fréquence de résonance $N_0$ (en Hz).
    \item \textbf{Lire l'intensité maximale} $I_{\max}$ (ordonnée du maximum). On en déduit la résistance équivalente du circuit : $R = \frac{U}{I_{\max}}$ (si $U$ est connue).
    \item \textbf{Calculer l'inconnue ($L$ ou $C$)} grâce à la relation de résonance :
    \[
    N_0 = \frac{1}{2\pi\sqrt{LC}} \quad \Longleftrightarrow \quad LC = \frac{1}{(2\pi N_0)^2} = \frac{1}{\omega_0^2}.
    \]
    Si $C$ est connu (ex. condensateur étalon), on en déduit $L = \frac{1}{(2\pi N_0)^2 C}$.
    \item \textbf{Vérification (optionnelle) :} Mesurer la bande passante $\Delta N$ à $-3\,\text{dB}$ (section 4) pour déterminer le facteur de qualité $Q = \frac{N_0}{\Delta N}$ et vérifier la cohérence avec $Q = \frac{L\omega_0}{R}$.
\end{enumerate}
\end{methode}

\begin{exoresolu}[Détermination de l'inductance d'une bobine par résonance]
\textbf{Énoncé.} On souhaite déterminer la valeur de l'inductance $L$ d'une bobine réelle (qui a une résistance interne $r$). On la place en série avec un condensateur étalon de capacité $C = 47\,\text{nF} = 4,7 \times 10^{-8}\,\text{F}$ et un GBF de tension efficace $U = 2,0\,\text{V}$ constante. On mesure l'intensité efficace $I$ pour différentes fréquences $N$. La courbe $I(N)$ présente un maximum $I_{\max} = 12,0\,\text{mA}$ pour $N_0 = 10,6\,\text{kHz}$.
\begin{enumerate}
    \item Calculer l'inductance $L$ de la bobine.
    \item Calculer la résistance totale $R$ du circuit (résistance interne $r$ + résistance des fils + résistance interne du GBF).
    \item Calculer les tensions $U_L$ et $U_C$ aux bornes de la bobine et du condensateur à la résonance.
\end{enumerate}
\tcblower
\textbf{Solution détaillée.}

\textbf{1. Calcul de $L$.}\par
On utilise la formule de la fréquence propre :
\[
N_0 = \frac{1}{2\pi\sqrt{LC}} \quad \Rightarrow \quad L = \frac{1}{(2\pi N_0)^2 C}.
\]
$N_0 = 10,6\,\text{kHz} = 1,06 \times 10^4\,\text{Hz}$.
\[
(2\pi N_0)^2 = (2\pi \times 1,06 \times 10^4)^2 \approx (6,66 \times 10^4)^2 \approx 4,44 \times 10^9\,\text{s}^{-2}.
\]
\[
L = \frac{1}{4,44 \times 10^9 \times 4,7 \times 10^{-8}} = \frac{1}{208,7} \approx 4,79 \times 10^{-3}\,\text{H} = \textbf{4,79 mH}.
\]

\textbf{2. Calcul de $R$.}\par
À la résonance, $Z = R$ et $I_{\max} = U/R$.
\[
R = \frac{U}{I_{\max}} = \frac{2,0}{12,0 \times 10^{-3}} \approx \textbf{167 }\Omega.
\]
C'est la résistance équivalente série du circuit à la fréquence $N_0$ (incluant les pertes par effet de peau dans la bobine, etc.).

\textbf{3. Tensions $U_L$ et $U_C$ à la résonance.}\par
Pulsation $\omega_0 = 2\pi N_0 \approx 6,66 \times 10^4\,\text{rad}\cdot\text{s}^{-1}$.
Réactance $X_L = L\omega_0 = 4,79 \times 10^{-3} \times 6,66 \times 10^4 \approx 319\,\Omega$.
$X_C = X_L = 319\,\Omega$ (vérification : $1/(C\omega_0) = 1/(4,7\times10^{-8} \times 6,66\times10^4) \approx 319\,\Omega$).
\[
U_L = U_C = I_{\max} \times X_L = 12,0 \times 10^{-3} \times 319 \approx \textbf{3,83 V}.
\]
\textbf{Observation :} $U_L = U_C \approx 3,8\,\text{V}$ alors que le générateur ne délivre que $U = 2,0\,\text{V}$. Le facteur de qualité est $Q = U_L/U \approx 1,9$. La surtension est modérée car $R$ est relativement grand ($Q = L\omega_0/R \approx 319/167 \approx 1,9$).
\end{exoresolu}

\courssub{Synthèse des points clés à retenir}

\begin{aretenir}[Résumé : Résonance d'intensité]
\begin{itemize}
    \item \textbf{Condition :} $\omega_0 = \dfrac{1}{\sqrt{LC}}$ \quad ($N_0 = \dfrac{1}{2\pi\sqrt{LC}}$).
    \item \textbf{À la résonance :} $Z = R$ (min), $I_{\max} = \dfrac{U}{R}$, $\varphi = 0$ ($u$ et $i$ en phase), circuit \textbf{résistif pur}.
    \item \textbf{Surtension :} $U_L = U_C = Q \cdot U$ avec $Q = \dfrac{L\omega_0}{R} = \dfrac{1}{RC\omega_0}$. Elles sont égales et opposées ($u_L + u_C = 0$).
    \item \textbf{Courbe $I(N)$ :} Pic en $N_0$. \textbf{Aigu} si $R$ faible (sélectif), \textbf{large} si $R$ grand (peu sélectif).
    \item \textbf{Application :} Circuit d'accord radio (variation de $C$ pour sélectionner $N_0$).
\end{itemize}
\end{aretenir}

\coursec{Bande passante à -3 dB et facteur de qualité}

\noindent Dans la section précédente, nous avons mis en évidence le phénomène de \textbf{résonance d'intensité} : pour une fréquence propre $N_0 = \frac{1}{2\pi\sqrt{LC}}$, l'impédance du circuit RLC série est minimale ($Z_{\min}=R$) et l'intensité atteint son maximum $I_{\max} = \frac{U}{R}$. Mais un circuit réel ne fonctionne pas seulement à cette unique fréquence ; il répond à une \emph{gamme} de fréquences autour de $N_0$. La \textbf{bande passante} quantifie cette étendue, et le \textbf{facteur de qualité} $Q$ caractérise la \emph{sélectivité} du circuit, c'est-à-dire sa capacité à discriminer une fréquence voisine de la fréquence de résonance. Ces notions sont au cœur du filtrage (radio FM, télécommunications MTN/Orange, électronique de puissance du barrage de Nachtigal) et de la conception des circuits accordés.

\begin{savaistu}[Le syntoniseur radio et le facteur de qualité]
Lorsque vous tournez le bouton de syntonisation d'un poste de radio (ou que vous scannez les fréquences sur votre téléphone), vous faites varier la capacité $C$ (ou l'inductance $L$) d'un circuit RLC pour que sa fréquence propre $N_0$ coïncide avec la fréquence de la station désirée (ex. : $94,5 \text{ MHz}$ pour Radio Cameroon à Yaoundé). Un facteur de qualité $Q$ élevé assure que les stations voisines ($94,3 \text{ MHz}$ ou $94,7 \text{ MHz}$) sont fortement atténuées : c'est la \textbf{sélectivité}. Au Cameroun, l'espacement des canaux FM est de $200 \text{ kHz}$ ; un récepteur de qualité doit avoir une bande passante $\Delta N \ll 200 \text{ kHz}$.
\end{savaistu}

\courssub{Définition de la bande passante et fréquences de coupure}

\noindent Observons la courbe de résonance $I(N)$ tracée précédemment. L'intensité est maximale en $N_0$ et diminue de part et d'autre. On cherche à définir l'intervalle de fréquences où le circuit répond « significativement ». Le critère universel en électronique et en traitement du signal est le niveau \textbf{-3 dB}, qui correspond à une intensité ramenée à $\frac{1}{\sqrt{2}}$ de sa valeur maximale.

\begin{definition}[Bande passante à -3 dB]
On appelle \textbf{bande passante à -3 dB} (ou simplement \textbf{bande passante}) l'intervalle de fréquences $[N_1\,;\,N_2]$ pour lequel l'intensité efficace $I$ reste supérieure ou égale à $\frac{I_{\max}}{\sqrt{2}}$.
Les fréquences $N_1$ et $N_2$ ($N_1 < N_0 < N_2$) sont appelées \textbf{fréquences de coupure} (ou fréquences de -3 dB).
La \textbf{largeur de la bande passante} est notée $\Delta N$ (en hertz) ou $\Delta \omega$ (en rad/s) :
\[
\Delta N = N_2 - N_1 \qquad \text{et} \qquad \Delta \omega = \omega_2 - \omega_1 = 2\pi\Delta N.
\]
\end{definition}

\noindent \textbf{Pourquoi $\frac{1}{\sqrt{2}}$ ?} C'est le seuil où la \textbf{puissance moyenne} absorbée par la résistance $R$ tombe à la moitié de sa valeur maximale. En effet, $P = R I^2$. Si $I = \frac{I_{\max}}{\sqrt{2}}$, alors $P = R \frac{I_{\max}^2}{2} = \frac{P_{\max}}{2}$. La bande passante délimite donc la zone où le circuit fournit au moins la moitié de la puissance maximale disponible.

\courssub{Origine du terme « -3 dB » (complément Bac, admis)}

\noindent Le \textbf{décibel (dB)} est une unité logarithmique utilisée pour exprimer un rapport de puissances ou d'amplitudes. Pour les amplitudes (tension, intensité), le gain en décibels est défini par :
\[
G_{\text{dB}} = 20 \log_{10}\left(\frac{I}{I_{\max}}\right).
\]
Aux fréquences de coupure, $\frac{I}{I_{\max}} = \frac{1}{\sqrt{2}} \approx 0,707$. On obtient :
\[
G_{\text{dB}} = 20 \log_{10}\left(\frac{1}{\sqrt{2}}\right) = -10 \log_{10}(2) \approx -3,01 \text{ dB} \simeq -3 \text{ dB}.
\]
D'où l'appellation « bande passante à -3 dB ». \emph{Ce calcul n'est pas exigible en démonstration au Bac, mais la définition du seuil $\frac{I_{\max}}{\sqrt{2}}$ et l'appellation -3 dB sont à connaître.}

\courssub{Détermination analytique de la bande passante}

\noindent Nous allons établir la relation fondamentale $\Delta \omega = \frac{R}{L}$ (ou $\Delta N = \frac{R}{2\pi L}$) à partir de la condition $I = \frac{I_{\max}}{\sqrt{2}}$.

\begin{propriete}[Fréquences de coupure et largeur de bande -- Démonstration guidée exigible]
Les pulsations de coupure $\omega_1$ et $\omega_2$ sont les solutions de l'équation $\left|L\omega - \frac{1}{C\omega}\right| = R$.
La largeur de bande angulaire vaut :
\[
\Delta \omega = \omega_2 - \omega_1 = \frac{R}{L} \quad \left(\text{soit } \Delta N = \frac{R}{2\pi L}\right).
\]
\end{propriete}

\begin{experience}[Démonstration guidée : établissement de $\Delta \omega = R/L$]
\textbf{1. Condition sur l'impédance.}
L'intensité efficace est $I = \frac{U}{Z}$ avec $Z = \sqrt{R^2 + \left(L\omega - \frac{1}{C\omega}\right)^2}$.
Le maximum est $I_{\max} = \frac{U}{R}$ (pour $\omega_0$ tel que $L\omega_0 = \frac{1}{C\omega_0}$).
La condition $I = \frac{I_{\max}}{\sqrt{2}}$ s'écrit :
\[
\frac{U}{Z} = \frac{1}{\sqrt{2}} \frac{U}{R} \quad \Rightarrow \quad Z = R\sqrt{2}.
\]

\textbf{2. Équation en $\omega$.}
\[
\sqrt{R^2 + \left(L\omega - \frac{1}{C\omega}\right)^2} = R\sqrt{2}
\]
\[
\Rightarrow \quad R^2 + \left(L\omega - \frac{1}{C\omega}\right)^2 = 2R^2
\]
\[
\Rightarrow \quad \left(L\omega - \frac{1}{C\omega}\right)^2 = R^2
\]
\[
\Rightarrow \quad \left| L\omega - \frac{1}{C\omega} \right| = R. \quad \text{(Équation caractéristique des fréquences de coupure)}
\]

\textbf{3. Résolution.}
On a deux équations selon le signe de la réactance :
\begin{itemize}
    \item Pour $\omega > \omega_0$ (réactance inductive) : $L\omega_2 - \frac{1}{C\omega_2} = +R$.
    \item Pour $\omega < \omega_0$ (réactance capacitive) : $L\omega_1 - \frac{1}{C\omega_1} = -R$ (car la valeur absolue donne $R$).
\end{itemize}
Multiplions chaque équation par $\omega$ et mettons sous forme canonique :
\[
\begin{cases}
L\omega_2^2 - R\omega_2 - \frac{1}{C} = 0 & (\omega_2 > 0)\\
L\omega_1^2 + R\omega_1 - \frac{1}{C} = 0 & (\omega_1 > 0)
\end{cases}
\]

\textbf{4. Différence $\omega_2 - \omega_1$.}
Une méthode élégante consiste à soustraire les deux équations initiales (avant multiplication par $\omega$) :
\[
\left(L\omega_2 - \frac{1}{C\omega_2}\right) - \left(L\omega_1 - \frac{1}{C\omega_1}\right) = R - (-R) = 2R.
\]
Factorisons :
\[
L(\omega_2 - \omega_1) - \frac{1}{C}\left(\frac{1}{\omega_2} - \frac{1}{\omega_1}\right) = 2R
\]
\[
L\Delta\omega - \frac{1}{C}\frac{\omega_1 - \omega_2}{\omega_1\omega_2} = 2R
\]
\[
L\Delta\omega + \frac{1}{C}\frac{\Delta\omega}{\omega_1\omega_2} = 2R
\]
Or, le produit des racines positives $\omega_1$ et $\omega_2$ (issues des deux trinômes distincts) vaut $\omega_1\omega_2 = \frac{1}{LC} = \omega_0^2$ (vérifiable en calculant le produit des expressions explicites des racines).
Donc $\frac{1}{C\omega_1\omega_2} = \frac{1}{C}\cdot LC = L$.
L'équation devient :
\[
L\Delta\omega + L\Delta\omega = 2R \quad \Rightarrow \quad 2L\Delta\omega = 2R
\]
\[
\boxed{\Delta\omega = \frac{R}{L}} \qquad \text{et} \qquad \boxed{\Delta N = \frac{R}{2\pi L}}.
\]
\end{experience}

\noindent \textbf{Analyse dimensionnelle :} $[R] = \Omega = \text{kg}\cdot\text{m}^2\cdot\text{s}^{-3}\cdot\text{A}^{-2}$, $[L] = \text{H} = \text{kg}\cdot\text{m}^2\cdot\text{s}^{-2}\cdot\text{A}^{-2}$. Donc $[R/L] = \text{s}^{-1}$, dimension d'une pulsation. \textbf{Validé.}

\courssub{Facteur de qualité $Q$ : mesure de la sélectivité}

\noindent La largeur de bande $\Delta\omega$ dépend de $R$ et $L$. Pour comparer la sélectivité de circuits différents, on utilise une grandeur \textbf{adimensionnelle} : le facteur de qualité $Q$.

\begin{definition}[Facteur de qualité $Q$]
Le \textbf{facteur de qualité} $Q$ d'un circuit RLC série est défini par le rapport de la pulsation propre à la largeur de bande :
\[
Q = \frac{\omega_0}{\Delta\omega} = \frac{N_0}{\Delta N}.
\]
C'est une grandeur \textbf{sans dimension}. En utilisant $\omega_0 = \frac{1}{\sqrt{LC}}$ et $\Delta\omega = \frac{R}{L}$, on obtient les formes équivalentes :
\[
\boxed{Q = \frac{L\omega_0}{R} = \frac{1}{RC\omega_0} = \frac{1}{R}\sqrt{\frac{L}{C}}}.
\]
\end{definition}

\begin{aretenir}[Interprétation physique de $Q$]
\begin{itemize}
    \item \textbf{Sélectivité (acuité de la résonance) :} $Q$ grand $\Leftrightarrow$ $\Delta N$ petit $\Leftrightarrow$ courbe de résonance \textbf{aiguë} (pic étroit). Le circuit ne laisse passer qu'une bande de fréquences très étroite autour de $N_0$ : c'est un \textbf{filtre sélectif} (ex. : syntonisation radio, récepteur GPS, filtre anti-harmoniques sur le réseau SONATREL).
    \item $Q$ petit $\Leftrightarrow$ $\Delta N$ grand $\Leftrightarrow$ courbe \textbf{large} (pic aplati). Le circuit répond à une large gamme de fréquences : c'est un \textbf{filtre large bande} ou un circuit amorti (ex. : découplage d'alimentation, large bande audio).
    \item \textbf{Surtension à la résonance :} En résonance, $U_C = I_{\max} X_C = \frac{U}{R} \cdot \frac{1}{C\omega_0} = U \cdot \frac{1}{RC\omega_0} = U \cdot Q$.
    De même $U_L = Q \cdot U$. D'où le nom historique de \textbf{facteur de surtension}.
    \item \textbf{Énergie :} $Q = 2\pi \frac{\text{Énergie stockée max}}{\text{Énergie dissipée par période}}$. Un $Q$ élevé signifie des pertes faibles par cycle.
\end{itemize}
\end{aretenir}

\courssub{Détermination expérimentale de $Q$ à partir de la courbe de résonance}

\noindent En TP, on trace la courbe $I(N)$ (ou $U_R(N)$) à l'oscilloscope ou au traceur de courbes. La méthode graphique est directe.

\begin{methode}[Déterminer $Q$ graphiquement à partir de la courbe $I(N)$]
\begin{enumerate}
    \item Repérer la fréquence de résonance $N_0$ (abscisse du maximum) et lire l'intensité maximale $I_{\max}$.
    \item Calculer le seuil $I_{\text{seuil}} = \frac{I_{\max}}{\sqrt{2}} \approx 0,707 \times I_{\max}$.
    \item Tracer l'horizontale $I = I_{\text{seuil}}$ sur le graphique.
    \item Lire les abscisses $N_1$ et $N_2$ des deux intersections de cette droite avec la courbe.
    \item Calculer la bande passante $\Delta N = N_2 - N_1$.
    \item Calculer le facteur de qualité : $Q = \frac{N_0}{\Delta N}$.
\end{enumerate}
\end{methode}

\begin{popfigure}
\begin{tikzpicture}
\begin{axis}[
    width=\linewidth,
    height=0.55\linewidth,
    axis lines=middle,
    xlabel={$N$ (kHz)}, ylabel={$I$ (mA)},
    xmin=14, xmax=18,
    ymin=0, ymax=11,
    xtick={15.1, 15.9, 16.7},
    xticklabels={$N_1$, $N_0$, $N_2$},
    ytick={7.07, 10},
    yticklabels={$I_{\max}/\sqrt{2}$, $I_{\max}$},
    tick label style={font=\small},
    label style={font=\small},
    domain=14:18,
    samples=400,
    clip=false,
]
% Courbe de résonance théorique I = Imax / sqrt(1 + Q^2*(x/x0 - x0/x)^2)
% Paramètres pour l'exemple : N0=15.9, Q=10 -> deltaN = 1.59 -> N1=15.1, N2=16.7 approx
% Imax = 10 mA
\def\Nzero{15.9}
\def\Imax{10}
\def\Q{10}
\addplot[popcurve, thick, domain=14:18] 
    ({\x}, {\Imax/sqrt(1 + \Q^2 * (\x/\Nzero - \Nzero/\x)^2)});
% Ligne Imax
\draw[dashed, popmuted, thin] (axis cs:14,{\Imax}) -- (axis cs:18,{\Imax});
% Ligne Imax/sqrt(2)
\draw[dashed, popmuted, thin] (axis cs:14,{\Imax/sqrt(max(0,2))}) -- (axis cs:18,{\Imax/sqrt(max(0,2))});
% Flèches Delta N
\draw[<->, popblue, thick] (axis cs:15.1, 5) -- (axis cs:16.7, 5) node[midway, above, font=\small, popblue] {$\Delta N = N_2 - N_1$};
% Points de coupure
\addplot[only marks, mark=*, poporange, mark size=2.5] coordinates {(15.1, 7.07) (16.7, 7.07)};
\addplot[only marks, mark=*, popgreen, mark size=3] coordinates {(15.9, 10)};
\end{axis}
\end{tikzpicture}
\legende{Courbe de résonance d'intensité annotée : identification de $I_{\max}$, $N_0$, $N_1$, $N_2$, $\Delta N$ et du seuil $I_{\max}/\sqrt{2}$.}
\end{popfigure}

\courssub{Exemple chiffré}

\begin{exemplebox}[Calcul de la bande passante et du facteur de qualité]
Un circuit RLC série a une fréquence de résonance $N_0 = 15,9 \text{ kHz}$. On relève sur sa courbe de résonance les fréquences de coupure $N_1 = 15,1 \text{ kHz}$ et $N_2 = 16,7 \text{ kHz}$.
\begin{enumerate}
    \item Calculer la bande passante $\Delta N$.
    \item En déduire le facteur de qualité $Q$.
    \item Si la résistance du circuit vaut $R = 100 \ \Omega$, déterminer l'inductance $L$ et la capacité $C$.
\end{enumerate}
\tcblower
\textbf{Solution.}
\begin{enumerate}
    \item $\Delta N = N_2 - N_1 = 16,7 - 15,1 = \boxed{1,6 \text{ kHz}}$.
    \item $Q = \frac{N_0}{\Delta N} = \frac{15,9}{1,6} \approx \boxed{9,9} \simeq 10$. Le circuit est assez sélectif.
    \item On utilise $\Delta N = \frac{R}{2\pi L} \Rightarrow L = \frac{R}{2\pi \Delta N} = \frac{100}{2\pi \times 1600} \approx \boxed{9,95 \text{ mH}}$.
    On utilise $N_0 = \frac{1}{2\pi\sqrt{LC}} \Rightarrow C = \frac{1}{(2\pi N_0)^2 L} = \frac{1}{(2\pi \times 15900)^2 \times 9,95\times 10^{-3}} \approx \boxed{10,1 \text{ nF}}$.
    \textit{Vérification :} $Q = \frac{1}{R}\sqrt{\frac{L}{C}} = \frac{1}{100}\sqrt{\frac{9,95\times 10^{-3}}{10,1\times 10^{-9}}} \approx \frac{1}{100} \times 992 \approx 9,9$. \textbf{Coherent.}
\end{enumerate}
\end{exemplebox}

\courssub{Erreurs fréquentes à éviter}

\begin{attention}[Pièges classiques sur la bande passante et $Q$]
\begin{itemize}
    \item \textbf{Seuil erroné :} Utiliser $I_{\max}/2$ (moitié de l'intensité) au lieu de $I_{\max}/\sqrt{2}$ pour déterminer $N_1$ et $N_2$.
    \emph{Rappel :} Le critère -3 dB porte sur la \textbf{puissance} (divisée par 2), donc sur l'intensité (divisée par $\sqrt{2}$). $I_{\max}/2$ correspond à -6 dB (puissance divisée par 4).
    \item \textbf{Confusion $\Delta\omega$ et $\Delta N$ :} $\Delta\omega = 2\pi\Delta N$. La formule $\Delta\omega = R/L$ s'applique aux pulsations. Pour les fréquences en Hz : $\Delta N = R/(2\pi L)$.
    \item \textbf{Unités de $Q$ :} $Q$ est \textbf{sans dimension}. Ne jamais écrire « $Q = 10 \ \Omega$ » ou « $Q = 10 \text{ Hz}$ ».
    \item \textbf{Formule de $Q$ :} Ne pas confondre $Q = \frac{L\omega_0}{R}$ (circuit série) et $Q = R C \omega_0$ (circuit parallèle, non au programme de Tle C mais source d'erreur si on mélange les formules).
    \item \textbf{Lecture graphique :} Lire $N_1$ et $N_2$ sur l'axe des abscisses \emph{après} avoir tracé l'horizontale $I_{\max}/\sqrt{2}$, et non en estimant « à l'œil » la largeur à mi-hauteur du pic.
\end{itemize}
\end{attention}

\courssub{Exercice résolu : Exploitation complète d'une courbe de résonance}

\begin{exoresolu}[Caractérisation d'un circuit accordé]
\textbf{Énoncé.}
On étudie un circuit RLC série alimenté par un générateur de tension sinusoïdale d'amplitude constante $U_m = 10 \text{ V}$. La figure ci-dessous représente la courbe de variation de l'amplitude de l'intensité $I_m$ en fonction de la fréquence $N$ du générateur.
\begin{center}
\begin{tikzpicture}
\begin{axis}[
    width=0.8\linewidth, height=0.4\linewidth,
    axis lines=middle, xlabel={$N$ (kHz)}, ylabel={$I_m$ (mA)},
    xmin=90, xmax=110, ymin=0, ymax=55,
    xtick={95, 100, 105}, xticklabels={$95$, $100$, $105$},
    ytick={35.4, 50}, yticklabels={$35{,}4$, $50$},
    tick label style={font=\small}, label style={font=\small},
    domain=90:110, samples=300,
]
\def\Nzero{100} \def\Imax{50} \def\Q{10}
\addplot[popcurve, thick] ({\x}, {\Imax/sqrt(max(0,1 + \Q^2 * (\x/\Nzero - \Nzero/\x)^2))});
\draw[dashed, popmuted, thin] (axis cs:90,{\Imax}) -- (axis cs:110,{\Imax});
\draw[dashed, popmuted, thin] (axis cs:90,{\Imax/sqrt(max(0,2))}) -- (axis cs:110,{\Imax/sqrt(max(0,2))});
\addplot[only marks, mark=*, popgreen, mark size=3] coordinates {(100, 50)};
\addplot[only marks, mark=*, poporange, mark size=2.5] coordinates {(95, 35.4) (105, 35.4)};
\end{axis}
\end{tikzpicture}
\end{center}
L'échelle graphique donne : $I_{m,\max} = 50 \text{ mA}$ pour $N_0 = 100 \text{ kHz}$. Les intersections avec le seuil $I_{m,\max}/\sqrt{2} \approx 35,4 \text{ mA}$ se font pour $N_1 = 95 \text{ kHz}$ et $N_2 = 105 \text{ kHz}$.
La résistance du circuit est $R = 200 \ \Omega$.
\begin{enumerate}
    \item Déterminer la bande passante $\Delta N$ et le facteur de qualité $Q$.
    \item Calculer l'inductance $L$ et la capacité $C$ de la bobine et du condensateur.
    \item Calculer les amplitudes des tensions aux bornes de la bobine et du condensateur à la résonance. Commenter.
    \item On souhaite utiliser ce circuit comme filtre sélectif pour une station de radio émettant à $100 \text{ kHz}$. Une station parasite émet à $98 \text{ kHz}$ avec la même tension d'entrée. Quel est le rapport des intensités (station utile / station parasite) ? Le filtrage est-il efficace ?
\end{enumerate}
\tcblower
\textbf{Solution détaillée.}
\begin{enumerate}
    \item \textbf{Bande passante et facteur de qualité.}
    $\Delta N = N_2 - N_1 = 105 - 95 = \boxed{10 \text{ kHz}}$.
    $Q = \frac{N_0}{\Delta N} = \frac{100}{10} = \boxed{10}$.
    \item \textbf{Calcul de $L$ et $C$.}
    $\Delta\omega = 2\pi\Delta N = 2\pi \times 10^4 \text{ rad/s}$.
    Relation $\Delta\omega = \frac{R}{L} \Rightarrow L = \frac{R}{\Delta\omega} = \frac{200}{2\pi \times 10^4} = \frac{1}{100\pi} \approx \boxed{3,18 \text{ mH}}$.
    Fréquence propre : $\omega_0 = 2\pi N_0 = 2\pi \times 10^5 \text{ rad/s}$.
    $\omega_0 = \frac{1}{\sqrt{LC}} \Rightarrow C = \frac{1}{\omega_0^2 L} = \frac{1}{(2\pi \times 10^5)^2 \times 3,18\times 10^{-3}} \approx \boxed{79,6 \text{ nF}}$.
    \textit{Vérification rapide :} $Q = \frac{1}{R}\sqrt{\frac{L}{C}} = \frac{1}{200}\sqrt{\frac{3,18\times 10^{-3}}{79,6\times 10^{-9}}} = \frac{1}{200}\sqrt{4\times 10^4} = \frac{200}{200} = 10$. \textbf{OK.}
    \item \textbf{Surtensions à la résonance.}
    Intensité maximale : $I_{m,\max} = \frac{U_m}{R} = \frac{10}{200} = 0,05 \text{ A} = 50 \text{ mA}$ (cohérence parfaite avec le graphique).
    Amplitude tension aux bornes de $C$ (et $L$) : $U_{C,m} = U_{L,m} = Q \times U_m = 10 \times 10 = \boxed{100 \text{ V}}$.
    \textit{Commentaire :} Bien que le générateur ne fournisse que $10 \text{ V}$, les tensions aux bornes de $L$ et $C$ atteignent $100 \text{ V}$ (facteur de surtension $Q=10$). Elles sont en opposition de phase, leur somme est nulle, la tension aux bornes de $R$ vaut $U_m = 10 \text{ V}$. \textbf{Danger :} Il faut choisir des composants supportant cette surtension (tension de service du condensateur $> 100 \text{ V}$).
    \item \textbf{Sélectivité vis-à-vis de la station parasite à $98 \text{ kHz}$.}
    On cherche le rapport $\frac{I_m(100 \text{ kHz})}{I_m(98 \text{ kHz})}$.
    À la résonance ($100 \text{ kHz}$) : $I_{m,\text{utile}} = I_{m,\max} = 50 \text{ mA}$.
    À $98 \text{ kHz}$ : On utilise la formule de la courbe de résonance (admise) :
    \[
    \frac{I_m}{I_{m,\max}} = \frac{1}{\sqrt{1 + Q^2\left(\frac{N}{N_0} - \frac{N_0}{N}\right)^2}}.
    \]
    Calculons l'écart relatif : $\frac{N}{N_0} - \frac{N_0}{N} = \frac{98}{100} - \frac{100}{98} = 0,98 - 1,0204 \approx -0,0404$.
    $Q^2 \times (\dots)^2 = 100 \times (0,0404)^2 \approx 100 \times 0,00163 = 0,163$.
    $\frac{I_m}{I_{m,\max}} = \frac{1}{\sqrt{1 + 0,163}} = \frac{1}{\sqrt{1,163}} \approx \frac{1}{1,078} \approx 0,927$.
    Donc $I_{m,\text{parasite}} \approx 0,927 \times 50 \approx 46,4 \text{ mA}$.
    Rapport $\frac{I_{\text{utile}}}{I_{\text{parasite}}} \approx \frac{50}{46,4} \approx \boxed{1,08}$.
    \textbf{Conclusion :} Le rapport est très proche de 1. Le circuit \textbf{ne filtre pas efficacement} la station parasite à $98 \text{ kHz}$ (seulement $2 \text{ kHz}$ d'écart pour une bande passante de $10 \text{ kHz}$). Pour une bonne sélectivité, il faudrait un $Q$ bien plus grand (ex. $Q > 100$) ou une bande passante plus étroite ($\Delta N < 2 \text{ kHz}$).
\end{enumerate}
\end{exoresolu}

\coursec{Puissance moyenne et facteur de puissance}

Dans la section précédente, nous avons étudié la réponse en intensité du circuit RLC série : la courbe de résonance $I(\omega)$, la bande passante à $-3\ \mathrm{dB}$ et le facteur de qualité $Q$ nous ont appris \emph{comment} le circuit sélectionne les fréquences. Mais une question d'une importance pratique capitale reste en suspens : combien d'énergie le circuit consomme-t-il réellement ? C'est elle que paie l'abonné d'ENEO à la fin du mois, c'est elle qui fait chauffer les lignes de transport entre le barrage de Nachtigal et Douala. Cette section y répond en introduisant la \cle{puissance moyenne} et le \cle{facteur de puissance}, deux notions au cœur du métier d'électrotechnicien.

\courssub{Puissance instantanée : une grandeur fluctuante}

Considérons un dipôle quelconque (résistor, bobine, condensateur, moteur, ou l'association RLC série) alimenté par une tension sinusoïdale :
\[ u(t) = U\sqrt{2}\,\cos(\omega t) \qquad \text{et} \qquad i(t) = I\sqrt{2}\,\cos(\omega t - \varphi), \]
où $U$ et $I$ sont les valeurs efficaces et $\varphi$ le déphasage de $u$ par rapport à $i$.

\begin{definition}[Puissance instantanée]
La \cle{puissance instantanée} reçue par un dipôle est le produit, à chaque instant, de la tension à ses bornes par l'intensité qui le traverse :
\[ p(t) = u(t)\,i(t). \]
Elle s'exprime en watts (\si{W}).
\end{definition}

Développons ce produit à l'aide de l'identité trigonométrique $\cos a\,\cos b = \tfrac{1}{2}\left[\cos(a+b) + \cos(a-b)\right]$ :
\begin{align*}
p(t) &= U\sqrt{2}\cos(\omega t) \times I\sqrt{2}\cos(\omega t - \varphi)\\
&= 2UI \times \frac{1}{2}\left[\cos(2\omega t - \varphi) + \cos(\varphi)\right]\\
&= UI\cos\varphi + UI\cos(2\omega t - \varphi).
\end{align*}

Ce résultat est riche d'enseignements. La puissance instantanée est la somme :
\begin{itemize}
\item d'un terme \textbf{constant} : $UI\cos\varphi$ ;
\item d'un terme \textbf{sinusoïdal de pulsation $2\omega$} (fréquence double de celle de la tension), de valeur moyenne nulle sur une période.
\end{itemize}

Autrement dit, $p(t)$ \emph{fluctue} : elle oscille à $100\ \mathrm{Hz}$ si le réseau est à $50\ \mathrm{Hz}$. Elle peut même devenir négative pendant une fraction de période : le dipôle rend alors momentanément de l'énergie au générateur (c'est le cas des bobines et condensateurs qui stockent puis restituent l'énergie). Une grandeur qui change de signe au cours de chaque période n'est pas directement exploitable pour facturer l'énergie : c'est sa \emph{valeur moyenne} qui compte.

\begin{popfigure}
\begin{center}
\begin{tikzpicture}
\begin{axis}[
width=0.92\linewidth, height=6.8cm,
xlabel={$t$ (ms)}, ylabel={},
xmin=0, xmax=20, ymin=-1.6, ymax=1.9,
xtick={0,5,10,15,20},
grid=both, grid style={popline!40},
legend style={at={(0.5,1.02)}, anchor=south, legend columns=4, font=\small},
axis line style={popdark},
]
\addplot[popblue, very thick, domain=0:20, samples=200] {1.2*cos(360*x/20)};
\addplot[poporange, very thick, domain=0:20, samples=200] {cos(360*x/20 - 60)};
\addplot[popink, thick, domain=0:20, samples=300] {1.2*cos(360*x/20)*cos(360*x/20 - 60)};
\addplot[popgreen, dashed, very thick, domain=0:20, samples=2] {0.3};
\legend{$u(t)$, $i(t)$, $p(t)=u\,i$, $P = UI\cos\varphi$}
\end{axis}
\end{tikzpicture}
\end{center}
\legende{Tension $u(t)$, intensité $i(t)$ (déphasée de $\varphi = 60^\circ$ en retard) et puissance instantanée $p(t)$ sur une période ($f = 50\ \mathrm{Hz}$, amplitudes arbitraires $U_{\max} = 1{,}2$ et $I_{\max} = 1$). La courbe $p(t)$ oscille à $100\ \mathrm{Hz}$ autour de sa valeur moyenne $P = UI\cos\varphi = \tfrac{1}{2}U_{\max}I_{\max}\cos\varphi = 0{,}3$ (pointillés verts). Remarque : $p(t)$ devient négative par intervalles — le dipôle restitue alors de l'énergie au générateur.}
\end{popfigure}

\courssub{Puissance moyenne : le théorème fondamental}

\begin{propriete}[Puissance moyenne en régime sinusoïdal — démonstration guidée]
En régime sinusoïdal forcé, la \cle{puissance moyenne} consommée par un dipôle sur une période $T$ est :
\[ \boxed{P = U\,I\,\cos\varphi} \]
où $U$ et $I$ sont les valeurs efficaces et $\varphi$ le déphasage de la tension $u$ par rapport à l'intensité $i$.

\textbf{Démonstration guidée} (le calcul intégral est admis, la démarche est exigible) :
\begin{enumerate}
\item On écrit la définition de la valeur moyenne sur une période : $P = \dfrac{1}{T}\displaystyle\int_{0}^{T} p(t)\,\mathrm{d}t$.
\item On remplace $p(t)$ par son expression établie ci-dessus : $P = \dfrac{1}{T}\displaystyle\int_{0}^{T} \left[UI\cos\varphi + UI\cos(2\omega t - \varphi)\right]\mathrm{d}t$.
\item La moyenne du terme constant est ce terme lui-même : $\dfrac{1}{T}\displaystyle\int_{0}^{T} UI\cos\varphi\,\mathrm{d}t = UI\cos\varphi$.
\item La moyenne d'une fonction sinusoïdale sur un nombre entier de périodes est nulle : $\dfrac{1}{T}\displaystyle\int_{0}^{T} UI\cos(2\omega t - \varphi)\,\mathrm{d}t = 0$ (admission de l'intégration).
\item Conclusion : $P = UI\cos\varphi$.
\end{enumerate}
\end{propriete}

\begin{attention}[Le piège classique : $P = UI$]
L'erreur la plus fréquente au Baccalauréat est d'écrire $P = UI$ en oubliant le facteur $\cos\varphi$. Cette formule n'est exacte que dans deux cas particuliers :
\begin{itemize}
\item pour un \textbf{résistor seul} ($\varphi = 0$, donc $\cos\varphi = 1$) ;
\item à la \textbf{résonance d'intensité} du circuit RLC ($\omega = \omega_0$, $\varphi = 0$).
\end{itemize}
Dans tous les autres cas, $P < UI$. Le produit $UI$ seul n'est \emph{pas} une puissance consommée : on l'appelle \cle{puissance apparente}, notée $S = UI$, exprimée en \textbf{voltampères} (\si{VA}) pour la distinguer du watt. C'est elle qui sert au dimensionnement des transformateurs et des alternateurs (un groupe électrogène de « $5\ \mathrm{kVA}$ » ne délivre $5\ \mathrm{kW}$ que si $\cos\varphi = 1$).
\end{attention}

\begin{exemplebox}[Moteur d'une pompe à eau à Yaoundé]
Un moteur monophasé est alimenté sous $U = 220\ \mathrm{V}$ (valeur efficace) et absorbe $I = 5\ \mathrm{A}$ avec un facteur de puissance $\cos\varphi = 0{,}8$.
\begin{itemize}
\item Puissance apparente : $S = UI = 220 \times 5 = 1100\ \mathrm{VA}$.
\item Puissance moyenne réellement consommée : $P = UI\cos\varphi = 220 \times 5 \times 0{,}8 = 880\ \mathrm{W}$.
\end{itemize}
Écrire $P = 1100\ \mathrm{W}$ serait une erreur de $25\ \%$ ! Seuls $880\ \mathrm{W}$ sont convertis en travail mécanique et en chaleur ; les $1100\ \mathrm{VA}$ caractérisent la « place » que le moteur occupe sur le réseau.
\end{exemplebox}

\courssub{Qui consomme la puissance ? Le rôle exclusif du résistor}

Appliquons le théorème aux trois dipôles élémentaires, en utilisant les déphasages établis dans la section 2 :

\begin{center}
\begin{tabularx}{\linewidth}{|l|c|c|X|}
\hline
\rowcolor{popblueL} \textbf{Dipôle} & $\varphi$ & $\cos\varphi$ & \textbf{Puissance moyenne} \\
\hline
Résistor $R$ & $0$ & $1$ & $P = UI = RI^2$ \\
\hline
Bobine parfaite $L$ & $+\dfrac{\pi}{2}$ & $0$ & $P = 0$ \\
\hline
Condensateur parfait $C$ & $-\dfrac{\pi}{2}$ & $0$ & $P = 0$ \\
\hline
\end{tabularx}
\end{center}

\begin{propriete}[Seul le résistor consomme de la puissance moyenne]
Dans un circuit RLC série parcouru par un courant d'intensité efficace $I$, la puissance moyenne totale consommée est :
\[ P = R\,I^2. \]
La bobine parfaite et le condensateur parfait ne consomment \textbf{aucune} puissance moyenne : ils \emph{échangent} de l'énergie avec le reste du circuit (stockage magnétique $\tfrac{1}{2}Li^2$ pour $L$, stockage électrique $\tfrac{1}{2}Cu_C^2$ pour $C$) sans en dissiper en moyenne. On dit qu'ils font circuler une énergie \emph{réactive}.
\end{propriete}

Vérifions la cohérence avec le théorème général. Pour le RLC série, $U = ZI$ et $\cos\varphi = \dfrac{R}{Z}$ (lecture directe sur la construction de Fresnel), donc :
\[ P = UI\cos\varphi = (ZI)\,I\,\frac{R}{Z} = RI^2. \]
Les deux expressions coïncident : le bilan est bouclé. On retiendra aussi l'écriture $P = \dfrac{U_R^2}{R}$ où $U_R = RI$ est la tension efficace aux bornes du résistor.

\begin{aretenir}[Les trois expressions de la puissance moyenne]
Pour un dipôle en régime sinusoïdal : $P = UI\cos\varphi$. Pour un circuit RLC série : $P = RI^2 = \dfrac{U_R^2}{R}$. La puissance moyenne s'exprime en \textbf{watts} (\si{W}) ; le produit $UI$ (puissance apparente) s'exprime en \textbf{voltampères} (\si{VA}).
\end{aretenir}

\courssub{Facteur de puissance : définition et conséquences pratiques}

\begin{definition}[Facteur de puissance]
Le \cle{facteur de puissance} d'un dipôle ou d'une installation est le rapport :
\[ \cos\varphi = \frac{P}{U\,I}. \]
C'est un nombre sans dimension, compris entre $0$ et $1$. Pour un circuit RLC série, la construction de Fresnel donne directement :
\[ \cos\varphi = \frac{R}{Z} = \frac{R}{\sqrt{R^2 + \left(L\omega - \dfrac{1}{C\omega}\right)^2}}. \]
\end{definition}

Pourquoi les électrotechniciens surveillent-ils cette grandeur avec tant d'attention ? Raisonnons à \textbf{puissance utile $P$ et tension $U$ fixées} (c'est le cahier des charges d'une usine : elle a besoin d'une puissance donnée, et le réseau lui impose sa tension). Alors :
\[ I = \frac{P}{U\cos\varphi}. \]
Plus le facteur de puissance est \emph{faible}, plus l'intensité appelée est \emph{grande}. Or les pertes par effet Joule dans les lignes de transport et les transformateurs valent $P_{\text{ligne}} = rI^2$ ($r$ : résistance des lignes) : elles croissent comme le \emph{carré} de l'intensité. Un $\cos\varphi$ de $0{,}5$ au lieu de $1$ double le courant et \emph{quadruple} les pertes en ligne, avec échauffement des câbles et chutes de tension supplémentaires — sans aucun bénéfice pour l'usager, qui ne paie que $P$.

\begin{exoresolu}[Impact du facteur de puissance sur une ligne]
Un atelier de menuiserie à Bafoussam consomme une puissance moyenne $P = 11\ \mathrm{kW}$ sous $U = 220\ \mathrm{V}$. La ligne d'alimentation a une résistance totale $r = 0{,}20\ \Omega$.
\begin{enumerate}
\item Calculer l'intensité appelée et les pertes en ligne si $\cos\varphi = 1$.
\item Mêmes questions si $\cos\varphi = 0{,}5$. Conclure.
\end{enumerate}
\tcblower
\textbf{1.} Cas $\cos\varphi = 1$ :
\[ I_1 = \frac{P}{U\cos\varphi} = \frac{11\,000}{220 \times 1} = 50\ \mathrm{A}, \qquad P_{\text{ligne},1} = rI_1^2 = 0{,}20 \times 50^2 = 500\ \mathrm{W}. \]
\textbf{2.} Cas $\cos\varphi = 0{,}5$ :
\[ I_2 = \frac{11\,000}{220 \times 0{,}5} = 100\ \mathrm{A}, \qquad P_{\text{ligne},2} = 0{,}20 \times 100^2 = 2\,000\ \mathrm{W} = 2\ \mathrm{kW}. \]
\textbf{Conclusion :} pour la \emph{même} puissance utile de $11\ \mathrm{kW}$, diviser le facteur de puissance par deux double le courant et multiplie les pertes en ligne par quatre (de $0{,}5\ \mathrm{kW}$ à $2\ \mathrm{kW}$, soit environ $18\ \%$ de la puissance utile gaspillée en chaleur). D'où l'exigence des distributeurs : maintenir $\cos\varphi$ proche de $1$.
\end{exoresolu}

\courssub{Relevé du facteur de puissance par condensateurs}

Les récepteurs industriels (moteurs, transformateurs, fours à induction) sont de nature \emph{inductive} : ils imposent un déphasage $\varphi > 0$ et un facteur de puissance inférieur à $1$. Comment le relever sans modifier les machines ?

L'idée est élégante : le condensateur, lui, produit un déphasage de signe \emph{opposé} ($\varphi_C = -\pi/2$). En plaçant une \textbf{batterie de condensateurs en parallèle} sur l'installation, le courant réactif absorbé par les bobinages est compensé localement par le courant réactif fourni par les condensateurs : l'énergie réactive oscille entre les moteurs et les condensateurs \emph{sans transiter par le réseau}. Le déphasage global $\varphi'$ vu par le distributeur diminue, $\cos\varphi'$ se rapproche de $1$, l'intensité en ligne diminue — et avec elle les pertes et les pénalités. La puissance moyenne $P$, elle, est inchangée (le condensateur parfait ne consomme rien) : seule la puissance apparente $S = UI$ diminue.

\begin{popfigure}
\begin{center}
\begin{tikzpicture}[scale=1.05, font=\small]
% Triangle 1 : avant compensation
\begin{scope}
\draw[popdark, thick, ->] (0,0) -- (3.2,0) node[midway, below]{$P$};
\draw[popdark, thick, ->] (3.2,0) -- (3.2,1.9) node[midway, right]{$Q$};
\draw[popink, very thick, ->] (0,0) -- (3.2,1.9) node[midway, above left]{$S = UI$};
\draw[poporange, thick] (0.9,0) arc (0:30.7:0.9) node[midway, right=2pt]{$\varphi$};
\node[popdark] at (1.6,-0.9) {\textbf{Avant} : $\cos\varphi$ faible};
\end{scope}
% Triangle 2 : après compensation
\begin{scope}[xshift=7cm]
\draw[popdark, thick, ->] (0,0) -- (3.2,0) node[midway, below]{$P$ (inchangée)};
\draw[popgreen, thick, ->] (3.2,0) -- (3.2,0.8) node[midway, right]{$Q'$};
\draw[popblue, very thick, ->] (0,0) -- (3.2,0.8) node[midway, above left]{$S'$};
\draw[popgreen, thick] (0.9,0) arc (0:14:0.9) node[midway, right=2pt]{$\varphi'$};
\node[popdark] at (1.6,-0.9) {\textbf{Après} : $\cos\varphi' > \cos\varphi$};
\end{scope}
\end{tikzpicture}
\end{center}
\legende{Triangle des puissances (complément culturel, hors exigibilité Bac) : la puissance moyenne $P$ (en \si{W}), la puissance réactive $Q = UI\sin\varphi$ (en voltampères réactifs, \si{var}) et la puissance apparente $S = UI$ (en \si{VA}) forment un triangle rectangle : $S^2 = P^2 + Q^2$. L'adjonction de condensateurs réduit $Q$ à $P$ constante : l'angle $\varphi$ diminue et le facteur de puissance se relève.}
\end{popfigure}

\begin{savaistu}[Les pénalités d'ENEO et les batteries de condensateurs]
Au Cameroun, comme dans la plupart des pays, le distributeur d'électricité (ENEO, héritier de la SONEL) facture aux gros clients industriels — cimenteries de Bonabéri, brasseries, huileries — non seulement l'énergie active consommée, mais aussi des \textbf{pénalités} lorsque leur facteur de puissance descend sous un seuil contractuel (typiquement $\cos\varphi < 0{,}9$ environ, souvent exprimé via $\tan\varphi$). Pourquoi ? Parce qu'un faible $\cos\varphi$ oblige le distributeur à transporter un courant plus fort, donc à surdimensionner lignes et transformateurs et à supporter davantage de pertes Joule, pour une énergie facturée identique. Les industriels installent alors des \textbf{batteries de condensateurs} — parfois à enclenchement automatique piloté par un relais varmétrique — dans leurs postes de transformation. L'investissement est rapidement rentabilisé par la suppression des pénalités et la baisse des chutes de tension internes. La physique du RLC série que tu étudies aujourd'hui gouverne ainsi directement la facture électrique des usines camerounaises !
\end{savaistu}

\begin{methode}[Exploiter les grandeurs de puissance dans un exercice]
\begin{enumerate}
\item Identifier les données : $U$, $I$, $Z$, $R$, $\varphi$ ou $\cos\varphi$.
\item Si le déphasage est demandé : lire $\cos\varphi = \dfrac{R}{Z}$ (RLC série) ou $\cos\varphi = \dfrac{P}{UI}$.
\item Calculer la puissance moyenne avec $P = UI\cos\varphi$ \emph{ou} $P = RI^2$ — vérifier que les deux donnent le même résultat.
\item Distinguer soigneusement dans la rédaction : $P$ en watts (\si{W}), $S = UI$ en voltampères (\si{VA}).
\item Pour une question sur les pertes en ligne : écrire $I = \dfrac{P}{U\cos\varphi}$ puis $P_{\text{ligne}} = rI^2$.
\end{enumerate}
\end{methode}

\begin{attention}[Conditions d'application et erreurs de rédaction]
\begin{itemize}
\item La formule $P = UI\cos\varphi$ n'est valable qu'en régime \textbf{sinusoïdal} ; en continu, $P = UI$ tout court.
\item $U$ et $I$ sont des valeurs \textbf{efficaces}, jamais des valeurs maximales : avec $U_{\max}$ et $I_{\max}$, il faudrait écrire $P = \tfrac{1}{2}U_{\max}I_{\max}\cos\varphi$.
\item Ne pas écrire « la bobine ne consomme pas d'énergie » sans préciser « \emph{en moyenne} » : à chaque instant, elle échange bien de l'énergie avec le circuit.
\item Une bobine \emph{réelle} possède une résistance $r$ : elle consomme alors $P = rI^2$, ce qui explique l'échauffement des moteurs et des transformateurs.
\end{itemize}
\end{attention}

En résumé, la puissance moyenne $P = UI\cos\varphi = RI^2$ mesure l'énergie réellement convertie par le circuit chaque seconde ; le facteur de puissance $\cos\varphi = R/Z$ juge de la « qualité » de l'utilisation du courant. Il reste à apprendre à \emph{mesurer} ce déphasage $\varphi$ en pratique : ce sera l'objet de la dernière section, consacrée à l'exploitation des courbes observées à l'oscilloscope.

\coursec{Mesure du déphasage à l'oscilloscope et exploitation des courbes}

Après avoir établi les expressions de la puissance moyenne $P = U I \cos\varphi$ et du facteur de puissance $\cos\varphi$, il nous reste à déterminer expérimentalement ce déphasage $\varphi$ ainsi que les amplitudes $U_m$ et $I_m$ qui caractérisent complètement le circuit RLC série. L'oscilloscope, instrument roi de l'électronique, va nous permettre de \emph{visualiser} l'invisible : le décalage temporel entre la tension d'alimentation $u(t)$ et l'intensité $i(t)$. Rappelons un point essentiel : un oscilloscope classique ne mesure que des \emph{tensions}. Nous utiliserons donc la tension aux bornes du conducteur ohmique, $u_R(t) = R\,i(t)$, comme image fidèle du courant (à un facteur $R$ près, et surtout \emph{en phase} avec $i$).

\courssub{Montage expérimental et problème de la masse commune}

Pour observer simultanément la tension d'alimentation $u(t)$ aux bornes du dipôle RLC série et la tension $u_R(t)$ aux bornes de la résistance $R$ (image du courant), on branche la voie A (YA) de l'oscilloscope aux bornes de l'ensemble du dipôle et la voie B (YB) aux bornes de $R$.

\begin{attention}[Piège majeur : la masse commune de l'oscilloscope]
La plupart des oscilloscopes de laboratoire ont leurs entrées \textbf{reliées à la masse de l'appareil} (reliée à la terre par la prise secteur). Les bornes « masse » (le crocodile noir des sondes) de la voie A et de la voie B sont \emph{électriquement reliées entre elles} à l'intérieur de l'appareil.
Si l'on branche les deux sondes n'importe où dans le circuit, on court-circuite la portion de circuit située entre les deux crocodiles : le montage étudié n'est alors plus du tout le circuit RLC série !
\end{attention}

\begin{methode}[Disposition correcte du montage RLC pour oscilloscope à masse commune]
Pour éviter tout court-circuit parasite, on place \textbf{la résistance $R$ du côté de la masse du générateur}.
\begin{enumerate}
    \item Brancher le générateur basses fréquences (GBF) : sa borne « masse » (souvent noire) va vers le point $M$ du circuit.
    \item Monter le circuit série dans l'ordre : GBF (borne rouge) $\to$ bobine $L$ $\to$ condensateur $C$ $\to$ résistance $R$ $\to$ point $M$ (masse du GBF).
    \item Sonde voie A (YA) : pointe sur la borne rouge du GBF (entrée du dipôle RLC), crocodile sur le point $M$.
    \item Sonde voie B (YB) : pointe sur la borne de $R$ reliée au condensateur, crocodile sur le point $M$ (borne de $R$ reliée à la masse).
\end{enumerate}
Ainsi, les deux crocodiles (masses de l'oscilloscope) sont connectés au même point $M$ : aucun court-circuit. La voie A mesure $u(t)$, la voie B mesure $u_R(t) = R\,i(t)$.
\end{methode}

\begin{popfigure}
\begin{tikzpicture}[scale=0.95, >=Stealth, european]
    % Circuit série
    \draw (0,0) to[sV, l={$u(t)$}] (0,2.5) -- (1.5,2.5) to[L, l={$L$}] (3.5,2.5) to[C, l={$C$}] (5.5,2.5) to[R, l={$R$}] (7.5,2.5) -- (7.5,0) -- (0,0);
    % Point M et masse
    \draw (7.5,0) node[ground]{} node[below right] {$M$ (masse)};
    \fill (0,2.5) circle (1.5pt) node[above left] {$A$};
    \fill (5.5,2.5) circle (1.5pt) node[above] {$B$};
    % Sonde voie A
    \draw[popblue, thick] (0,2.5) -- (0,3.6) node[popblue, above] {YA : pointe en $A$};
    \draw[popblue, thick] (7.5,0) -- (8.6,0) node[popblue, right] {crocodile YA en $M$};
    % Sonde voie B
    \draw[poporange, thick] (5.5,2.5) -- (5.5,3.6) node[poporange, above] {YB : pointe en $B$};
    \draw[poporange, thick] (7.5,-0.6) -- (8.6,-0.6) node[poporange, right] {crocodile YB en $M$};
    % Tensions mesurées
    \node[popblue, align=center] at (1.5,-1.2) {Voie A : $u(t)$\\tension d'alimentation};
    \node[poporange, align=center] at (6,-1.2) {Voie B : $u_R(t)=R\,i(t)$\\image du courant};
\end{tikzpicture}
\legende{Montage correct pour visualiser $u(t)$ et $u_R(t)$ sur un oscilloscope à masse commune : la résistance $R$ est placée du côté de la masse, et les deux crocodiles sont reliés au même point $M$.}
\end{popfigure}

\courssub{Observation à l'écran : deux sinusoïdes décalées}

Une fois le montage réalisé et le GBF réglé sur une fréquence $f$, l'écran affiche deux courbes sinusoïdales de \emph{même période} $T$ (même fréquence), mais décalées horizontalement l'une par rapport à l'autre.

\begin{popfigure}
\begin{tikzpicture}[scale=1.0, >=Stealth]
    % Grille oscilloscope : 8 div horizontales, 6 verticales
    \draw[popline!50, step=0.3] (0,0) grid (12,6);
    \draw[popline, step=1.5] (0,0) grid (12,6);
    % Axes centraux
    \draw[popdark, thick, ->] (0,3) -- (12.4,3) node[right] {$t$};
    \draw[popdark, thick, ->] (0,3) -- (0,6.4) node[above] {tension};
    % Sinusoïde u(t) - Voie A (bleu), période 4 div = 6 cm
    \draw[popblue, very thick, domain=0:12, samples=200, variable=\x]
        plot ({\x}, {3 + 2*sin(60*\x)});
    % Sinusoïde u_R(t) - Voie B (orange), retard de 0.5 div = 0.75 cm
    \draw[poporange, very thick, domain=0:12, samples=200, variable=\x]
        plot ({\x}, {3 + 1.4*sin(60*(\x - 0.75))});
    % Mesure de la période
    \draw[<->, popdark] (1.5,5.6) -- (7.5,5.6) node[midway, above] {$T = 4$ div};
    % Mesure du décalage entre maxima
    \draw[<->, poppurple, thick] (1.5,4.6) -- (2.25,4.6) node[midway, above] {$\Delta t$};
    % Points caractéristiques (maxima)
    \fill[popblue] (1.5,5) circle (2pt);
    \fill[poporange] (2.25,4.4) circle (2pt);
    % Légendes
    \node[popblue, anchor=west] at (8.2,5.1) {YA : $u(t)$};
    \node[poporange, anchor=west] at (8.2,4.3) {YB : $u_R(t)$};
    \node[popmuted, anchor=west] at (8.2,1.2) {$T = 4$ div ; $\Delta t = 0{,}5$ div};
\end{tikzpicture}
\legende{Écran d'oscilloscope : deux sinusoïdes de même période $T$, décalées de $\Delta t$. La courbe orange ($u_R$, image de $i$) atteint son maximum \emph{après} la bleue ($u$) : le courant est en retard sur la tension.}
\end{popfigure}

\courssub{Relation fondamentale : du décalage temporel au déphasage}

Les deux signaux ont la même période $T$ et la même pulsation $\omega = 2\pi/T$. Un décalage temporel $\Delta t$ correspond à un décalage de phase $\varphi$ proportionnel.

\begin{propriete}[Lien entre décalage temporel et déphasage]
Soient deux signaux sinusoïdaux de même période $T$, décalés dans le temps de $\Delta t$. La valeur absolue de leur déphasage est :
\[
|\varphi| = \omega\,\Delta t = \frac{2\pi}{T}\,\Delta t = 2\pi\,\frac{\Delta t}{T}
\]
\emph{Justification :} sur une période complète $T$, la phase varie de $2\pi$ rad. Par proportionnalité, un intervalle $\Delta t$ correspond à une variation de phase $\varphi$ telle que $\dfrac{|\varphi|}{2\pi} = \dfrac{\Delta t}{T}$.
\end{propriete}

\begin{attention}[Unités et mode calculatrice]
La formule $|\varphi| = 2\pi\,\dfrac{\Delta t}{T}$ donne $\varphi$ en \textbf{radians} dès lors que $\Delta t$ et $T$ sont exprimés dans la même unité (s, ms, $\mu$s, ou divisions d'écran). Pour obtenir $\varphi$ en degrés : $|\varphi| = 360^\circ\,\dfrac{\Delta t}{T}$. \textbf{Vérifie le mode de ta calculatrice (RAD ou DEG) avant de calculer $\cos\varphi$ ou $\tan\varphi$ !}
\end{attention}

\courssub{Détermination du signe de $\varphi$ : avance ou retard ?}

La valeur absolue ne suffit pas ; il faut savoir si le courant est en avance ($\varphi < 0$, circuit capacitif) ou en retard ($\varphi > 0$, circuit inductif) sur la tension.

\begin{definition}[Courbe en avance]
Sur un écran d'oscilloscope, le temps s'écoule vers la \textbf{droite}. La courbe qui atteint son maximum (ou passe par zéro en croissant) \textbf{la première, c'est-à-dire la plus à gauche}, est la courbe \emph{en avance} de phase.
\end{definition}

\begin{methode}[Déterminer le signe de $\varphi = \varphi_u - \varphi_i$]
\begin{enumerate}
    \item Repérer un point caractéristique commun aux deux courbes (maximum, ou passage par zéro avec pente positive).
    \item Identifier quelle courbe atteint ce point \textbf{en premier} (la plus à gauche).
    \item Si $u(t)$ (voie A) est en avance sur $u_R(t)$ (voie B) : la tension est en avance sur le courant, $\varphi > 0$, le circuit est \textbf{inductif}.
    \item Si $u_R(t)$ (voie B) est en avance sur $u(t)$ (voie A) : le courant est en avance sur la tension, $\varphi < 0$, le circuit est \textbf{capacitif}.
\end{enumerate}
\end{methode}

\courssub{Méthode des divisions : mesure rapide sans chronomètre}

Il est inutile de connaître la valeur de la base de temps (s/div) pour mesurer $\varphi$ : le rapport $\Delta t / T$ est sans dimension. Il suffit de compter les divisions sur le graticule de l'écran.

\begin{methode}[Mesurer un déphasage à l'oscilloscope (méthode des divisions)]
\begin{enumerate}
    \item \textbf{Régler} la base de temps pour observer entre 1 et 3 périodes complètes à l'écran (précision optimale).
    \item \textbf{Mesurer la période $T$} en divisions horizontales (exemple : $T = 4{,}0$ div).
    \item \textbf{Mesurer le décalage $\Delta t$} en divisions horizontales entre les deux courbes (exemple : $\Delta t = 0{,}5$ div), en prenant le même repère sur les deux courbes (maximum, ou zéro montant).
    \item \textbf{Calculer} $|\varphi| = 2\pi\,\dfrac{\Delta t}{T}$ (rad) ou $360^\circ\,\dfrac{\Delta t}{T}$ (deg).
    \item \textbf{Déterminer le signe} en repérant la courbe en avance (méthode précédente).
\end{enumerate}
\end{methode}

\begin{exemplebox}[Application numérique : mesure par divisions]
Sur l'écran de la figure ci-dessus, on lit :
\begin{itemize}
    \item période $T = 4{,}0$ divisions horizontales ;
    \item décalage $\Delta t = 0{,}5$ division (la courbe YB est décalée vers la droite par rapport à YA) ;
    \item courbe en avance : YA ($u(t)$) passe par son maximum avant YB ($u_R$).
\end{itemize}
Calcul :
\[
|\varphi| = 2\pi \times \frac{0{,}5}{4{,}0} = \frac{\pi}{4}\,\text{rad} = 45^\circ
\]
Comme $u$ est en avance sur $u_R$ (image de $i$), la tension est en avance sur le courant : $\varphi = +\dfrac{\pi}{4}$ rad ($+45^\circ$). Le circuit est de nature \textbf{inductive} ($L\omega > \dfrac{1}{C\omega}$).
\end{exemplebox}

\courssub{Exploitation complète des courbes : $U_m$, $I_m$, $Z$, $\cos\varphi$, $P$}

L'oscilloscope donne accès aux amplitudes (valeurs maximales, ou « de crête ») et au déphasage. On en déduit toutes les grandeurs efficaces et la puissance moyenne.

\begin{methode}[Exploitation complète d'un relevé oscilloscopique]
Données lues à l'écran : $U_m$ (amplitude voie A), $U_{Rm}$ (amplitude voie B), $T$ (div), $\Delta t$ (div), signe de $\varphi$.
\begin{enumerate}
    \item Valeurs efficaces : $U = \dfrac{U_m}{\sqrt{2}}$ \quad et \quad $I = \dfrac{I_m}{\sqrt{2}} = \dfrac{U_{Rm}}{R\sqrt{2}}$.
    \item Impédance : $Z = \dfrac{U}{I} = \dfrac{U_m}{U_{Rm}}\,R$.
    \item Déphasage : $\varphi = \pm\,2\pi\,\dfrac{\Delta t}{T}$ (signe selon avance ou retard).
    \item Facteur de puissance : $\cos\varphi$ (on doit vérifier la cohérence $\cos\varphi \approx \dfrac{R}{Z}$).
    \item Puissance moyenne absorbée : $P = U\,I\,\cos\varphi$.
    \item \emph{Vérification :} seule la résistance dissipe de la puissance, donc on doit retrouver $P = R\,I^2 = \dfrac{U_{Rm}^2}{2R}$.
\end{enumerate}
\end{methode}

\begin{attention}[Cohérence des mesures : un réflexe de candidat]
Les grandeurs mesurées doivent vérifier $\cos\varphi = \dfrac{R}{Z}$ (car $R = Z\cos\varphi$ dans le triangle des impédances). Si le $\cos\varphi$ issu de la mesure de $\Delta t$ est très différent de $R/Z$, c'est que la lecture des divisions est imprécise : il faut le signaler dans le compte rendu et privilégier les mesures les plus directes (amplitudes et $R$). Au Bac, cette vérification est très appréciée du correcteur.
\end{attention}

\begin{exoresolu}[Exploitation complète d'un circuit RLC série]
\textbf{Énoncé :} Un circuit RLC série ($R = 100\,\Omega$, $L = 100\,\text{mH}$, $C = 100\,\text{nF}$) est alimenté par un GBF délivrant une tension sinusoïdale de fréquence $f = 1{,}0\,\text{kHz}$. On observe à l'oscilloscope (base de temps $0{,}2\,\text{ms/div}$, sensibilités : YA $2\,\text{V/div}$, YB $1\,\text{V/div}$) :
\begin{itemize}
    \item voie A ($u$) : amplitude $2{,}5$ div ;
    \item voie B ($u_R$) : amplitude $3{,}0$ div ;
    \item période $T = 5{,}0$ div ;
    \item décalage $\Delta t = 0{,}75$ div, la courbe $u_R$ étant décalée vers la droite par rapport à $u$.
\end{itemize}
Exploiter ces résultats : déterminer $U$, $I$, $Z$, $\varphi$, $\cos\varphi$ et $P$.
\tcblower
\textbf{Solution détaillée :}
\begin{enumerate}
    \item \textbf{Vérification de la période :} $T = 5{,}0 \times 0{,}2 = 1{,}0\,\text{ms}$, donc $f = 1/T = 1{,}0\,\text{kHz}$ : cohérent avec le réglage du GBF.
    \item \textbf{Amplitudes et valeurs efficaces :}
    \begin{align*}
    U_m &= 2{,}5 \times 2{,}0 = 5{,}0\,\text{V} & U &= \frac{U_m}{\sqrt{2}} = \frac{5{,}0}{\sqrt{2}} \approx 3{,}54\,\text{V}\\
    U_{Rm} &= 3{,}0 \times 1{,}0 = 3{,}0\,\text{V} & I_m &= \frac{U_{Rm}}{R} = \frac{3{,}0}{100} = 30\,\text{mA}\\
    & & I &= \frac{I_m}{\sqrt{2}} \approx 21{,}2\,\text{mA}
    \end{align*}
    \item \textbf{Impédance :}
    \[
    Z = \frac{U}{I} = \frac{3{,}54}{21{,}2\times10^{-3}} \approx 167\,\Omega
    \qquad \text{(vérification : } Z = \frac{U_m}{U_{Rm}}\,R = \frac{5{,}0}{3{,}0}\times 100 \approx 167\,\Omega\text{)}
    \]
    \item \textbf{Déphasage :}
    \[
    |\varphi| = 2\pi \times \frac{0{,}75}{5{,}0} = 0{,}30\pi \approx 0{,}942\,\text{rad} \approx 54{,}0^\circ
    \]
    Signe : $u_R$ (image du courant) est décalée vers la droite, donc le courant est \emph{en retard} sur la tension : $\varphi \approx +0{,}94\,\text{rad}$. Le circuit est inductif.
    \item \textbf{Facteur de puissance :} $\cos\varphi = \cos(0{,}942) \approx 0{,}59$.
    Vérification de cohérence : $\dfrac{R}{Z} = \dfrac{100}{167} \approx 0{,}60$. Les deux valeurs sont proches : les mesures sont cohérentes.
    \item \textbf{Puissance moyenne :}
    \[
    P = U I \cos\varphi = 3{,}54 \times 21{,}2\times10^{-3} \times 0{,}59 \approx 44\,\text{mW}
    \]
    Vérification : $P = R I^2 = 100 \times (21{,}2\times10^{-3})^2 \approx 45\,\text{mW}$. Les deux calculs concordent (à la précision des lectures près) : l'exploitation est validée.
\end{enumerate}
\end{exoresolu}

\courssub{Vérification de la résonance à l'oscilloscope}

La résonance d'intensité ($f = f_0$) se caractérise par $\varphi = 0$ : la tension $u(t)$ et le courant $i(t)$ (donc $u_R(t)$) sont \textbf{en phase}.

\begin{savaistu}[Le test ultime de la résonance]
À la fréquence propre $f_0 = \dfrac{1}{2\pi\sqrt{LC}}$, l'impédance du dipôle RLC série est minimale et purement résistive ($Z = R$). Sur l'oscilloscope, les deux courbes $u(t)$ (voie A) et $u_R(t)$ (voie B) passent par zéro et par leurs maxima \textbf{aux mêmes instants} : elles sont en phase. C'est la méthode expérimentale la plus précise et la plus visuelle pour déterminer $f_0$ : on fait varier lentement la fréquence du GBF jusqu'à annuler le décalage $\Delta t$, puis on mesure $T_0$ à l'écran (ou on lit $f_0$ sur le GBF). C'est ce principe d'accord en fréquence qui permet à un poste radio de sélectionner une station parmi toutes celles reçues par l'antenne.
\end{savaistu}

\begin{experience}[TP : Recherche de la fréquence de résonance à l'oscilloscope]
\textbf{Matériel :} GBF, oscilloscope bi-courbe, conducteur ohmique ($R = 100\,\Omega$), bobine ($L = 47\,\text{mH}$, résistance interne $r$ faible), condensateur ($C = 100\,\text{nF}$), fils de connexion.

\textbf{Protocole :}
\begin{enumerate}
    \item Monter le circuit RLC série avec $R$ du côté de la masse (voir la méthode du montage).
    \item Régler le GBF sur une fréquence basse (par exemple $500\,\text{Hz}$). Observer le décalage : $u_R$ est-elle en avance ou en retard sur $u$ ?
    \item Augmenter progressivement la fréquence : le décalage $\Delta t$ diminue.
    \item \textbf{Repérer la fréquence $f_0$} pour laquelle les deux courbes sont en phase ($\Delta t = 0$).
    \item Mesurer $T_0$ à l'oscilloscope et calculer $f_0 = 1/T_0$. Comparer à la valeur théorique $f_0 = \dfrac{1}{2\pi\sqrt{LC}} \approx 2{,}3\,\text{kHz}$.
    \item Mesurer $U_{Rm}$ à la résonance : on constate que $U_{Rm}$ est maximale (résonance d'intensité) et voisine de $U_m$ si la résistance totale du circuit se réduit à $R$.
\end{enumerate}
\textbf{Interprétation :} à $f_0$, les réactances se compensent ($L\omega_0 = \dfrac{1}{C\omega_0}$) : les tensions $u_L$ et $u_C$, en opposition de phase, s'annulent mutuellement ($u_L + u_C = 0$). La tension du générateur se retrouve intégralement aux bornes de la résistance totale du circuit, et l'intensité est maximale.
\end{experience}

\begin{attention}[Erreurs fréquentes à l'examen (Bac)]
\begin{itemize}
    \item Confondre $\Delta t$ et $T$ dans la formule et écrire $\varphi = 2\pi T/\Delta t$ (inversion).
    \item Oublier de préciser le \textbf{signe} de $\varphi$ : donner la seule valeur absolue ne suffit pas.
    \item Affirmer « le courant est en avance sur la tension » sans justifier par l'observation (quelle courbe atteint son maximum en premier ?).
    \item Oublier que $u_R$ est l'image de $i$ : c'est bien $u_R$ (voie B) qu'il faut comparer à $u$ (voie A), et non l'inverse.
    \item Utiliser la valeur affichée par le GBF au lieu des valeurs \emph{mesurées} à l'écran (divisions $\times$ sensibilité) pour calculer $Z$ ou $P$.
    \item Négliger l'analyse dimensionnelle : $\Delta t/T$ est sans dimension, donc $\varphi = 2\pi\Delta t/T$ est bien en radians.
\end{itemize}
\end{attention}

\begin{aretenir}[Bilan : l'oscilloscope, outil complet d'étude du RLC série]
\begin{center}
\begin{tabularx}{\linewidth}{|l|X|X|}\hline
\rowcolor{popblueL}
\textbf{Grandeurs lues à l'écran} & \textbf{Calculs directs} & \textbf{Paramètres du circuit} \\ \hline
$T$ (div) $\times$ base de temps & $f = 1/T$ ; $\omega = 2\pi f$ & Vérification de $f_0$ théorique \\ \hline
$U_m$ (div) $\times$ sensibilité YA & $U = U_m/\sqrt{2}$ & Tension efficace aux bornes du RLC \\ \hline
$U_{Rm}$ (div) $\times$ sensibilité YB & $I = U_{Rm}/(R\sqrt{2})$ & Intensité efficace dans le circuit \\ \hline
$\Delta t$ (div) entre $u$ et $u_R$ & $\varphi = \pm 2\pi \Delta t / T$ & Nature du circuit (inductif ou capacitif) \\ \hline
\multicolumn{3}{|c|}{\textbf{Exploitation finale}} \\ \hline
\multicolumn{3}{|X|}{$Z = \dfrac{U}{I} = \dfrac{U_m}{U_{Rm}}\,R$ \quad ; \quad $\cos\varphi = \dfrac{R}{Z}$ \quad ; \quad $P = U I \cos\varphi = R I^2 = \dfrac{U_{Rm}^2}{2R}$} \\ \hline
\end{tabularx}
\end{center}
\end{aretenir}

\courssub{Exercice d'entraînement : analyse d'un écran d'oscilloscope}

\begin{exercice}[Analyse de relevés oscilloscopiques]
On alimente un circuit RLC série ($R = 220\,\Omega$) par un GBF. L'oscilloscope (base de temps $50\,\mu\text{s/div}$, sensibilités : YA $5\,\text{V/div}$, YB $2\,\text{V/div}$) affiche :
\begin{itemize}
    \item voie A ($u$) : sinusoïde d'amplitude $3{,}2$ div ;
    \item voie B ($u_R$) : sinusoïde d'amplitude $1{,}8$ div ;
    \item période mesurée : $T = 4{,}0$ div ;
    \item décalage entre les passages par zéro montants : $\Delta t = 0{,}85$ div ; la courbe YB passe par zéro \emph{avant} la courbe YA.
\end{itemize}
\begin{enumerate}
    \item Déterminer la fréquence $f$ du GBF.
    \item Calculer les valeurs efficaces $U$ et $I$.
    \item Déterminer l'impédance $Z$ et le déphasage $\varphi$ (en rad et en degrés, avec son signe).
    \item Vérifier la cohérence $\cos\varphi \approx R/Z$, puis en déduire la réactance équivalente $X = L\omega - \dfrac{1}{C\omega}$ et la nature du circuit.
    \item Calculer la puissance moyenne absorbée $P$ de deux façons indépendantes.
\end{enumerate}
\end{exercice}

\begin{corrige}[Exercice : Analyse de relevés oscilloscopiques]
\begin{enumerate}
    \item $T = 4{,}0 \times 50\,\mu\text{s} = 200\,\mu\text{s} = 2{,}00\times10^{-4}\,\text{s}$, donc $f = \dfrac{1}{T} = 5{,}00\,\text{kHz}$.
    \item $U_m = 3{,}2 \times 5 = 16{,}0\,\text{V}$, donc $U = \dfrac{16{,}0}{\sqrt{2}} \approx 11{,}3\,\text{V}$.
    $U_{Rm} = 1{,}8 \times 2 = 3{,}6\,\text{V}$, donc $I_m = \dfrac{U_{Rm}}{R} = \dfrac{3{,}6}{220} \approx 16{,}4\,\text{mA}$ et $I = \dfrac{I_m}{\sqrt{2}} \approx 11{,}6\,\text{mA}$.
    \item $Z = \dfrac{U}{I} = \dfrac{11{,}3}{11{,}6\times10^{-3}} \approx 978\,\Omega$ (vérification : $Z = \dfrac{U_m}{U_{Rm}}\,R = \dfrac{16{,}0}{3{,}6}\times 220 \approx 978\,\Omega$).
    $|\varphi| = 2\pi \times \dfrac{0{,}85}{4{,}0} \approx 1{,}34\,\text{rad} \approx 76{,}5^\circ$.
    Signe : YB ($u_R$, image de $i$) est en avance sur YA ($u$), donc le courant est en avance sur la tension : $\varphi \approx -1{,}34\,\text{rad}$ ($-76{,}5^\circ$).
    \item $\cos\varphi = \cos(1{,}34) \approx 0{,}23$. Cohérence : $\dfrac{R}{Z} = \dfrac{220}{978} \approx 0{,}22$ : les deux valeurs concordent, les mesures sont fiables.
    $\sin\varphi = \sin(-1{,}34) \approx -0{,}97$, donc $X = Z\sin\varphi = 978 \times (-0{,}97) \approx -950\,\Omega$.
    $X < 0$ signifie $\dfrac{1}{C\omega} > L\omega$ : le circuit est de nature \textbf{capacitive}, ce qui confirme le signe négatif de $\varphi$.
    \item Première méthode : $P = U I \cos\varphi = 11{,}3 \times 11{,}6\times10^{-3} \times 0{,}23 \approx 30\,\text{mW}$.
    Deuxième méthode : $P = R I^2 = 220 \times (11{,}6\times10^{-3})^2 \approx 30\,\text{mW}$.
    Les deux calculs concordent : la puissance moyenne absorbée est $P \approx 30\,\text{mW}$, intégralement dissipée par effet Joule dans la résistance.
\end{enumerate}
\end{corrige}

Cette section termine notre étude des oscillations électriques forcées. Tu maîtrises désormais la modélisation par l'équation différentielle, la construction de Fresnel, l'impédance et le déphasage, la résonance d'intensité, la bande passante et le facteur de qualité, la puissance moyenne, ainsi que l'exploitation expérimentale à l'oscilloscope. Ces compétences sont au cœur de l'électronique moderne : du circuit d'accord d'un récepteur radio (résonance) à la compensation de l'énergie réactive dans les installations d'ENEO (amélioration du facteur de puissance), en passant par les filtres des téléphones portables. \textbf{Bon courage pour le Baccalauréat !}

\newpage

\coursec{Activité d'intégration}

\coursec{Oscillations électriques forcées en régime sinusoïdal}

\courssub{Objectifs de l'activité}

À l'issue de cette activité d'intégration, tu seras capable de :

\begin{itemize}
\item modéliser un circuit d'accord radio par un dipôle \cle{RLC série} alimenté par une tension sinusoïdale ;
\item exploiter la condition de \cle{résonance d'intensité} pour déterminer une capacité d'accord ;
\item calculer une \cle{impédance}, une intensité efficace et une tension efficace en régime forcé ;
\item utiliser le \cle{facteur de qualité} $Q$ pour interpréter un phénomène de surtension ;
\item déterminer une \cle{bande passante à $-3$ dB} et conclure sur la sélectivité d'un récepteur ;
\item construire un \cle{diagramme de Fresnel} hors résonance et en déduire le déphasage ;
\item calculer une \cle{puissance moyenne} et discuter l'intérêt du \cle{facteur de puissance} dans une installation réelle.
\end{itemize}

\courssub{Situation}

Dans le quartier Nkolndongo à Yaoundé, Rodrigue, élève de Terminale C et passionné d'électronique, anime le club « Sciences et Technologies » de son lycée. Pour la fête de la jeunesse, le club veut fabriquer un récepteur radio simple capable de capter correctement une station FM locale émettant sur la fréquence $f_0 = \SI{100}{MHz}$.

Après avoir consulté un ancien manuel d'électronique, Rodrigue modélise l'étage d'accord du récepteur (antenne + circuit d'entrée) par un \cle{dipôle RLC série} : la bobine d'accord a une inductance $L = \SI{0,25}{\micro H}$, la résistance totale du circuit (bobine + fils + antenne ramenée) vaut $R = \SI{2}{\ohm}$, et le condensateur d'accord a une capacité $C$ réglable. Le signal capté par l'antenne se comporte comme un générateur de tension sinusoïdale de valeur efficace $U = \SI{10}{mV}$ et de fréquence égale à celle de la station écoutée.

Rodrigue se pose quatre questions techniques :

\begin{enumerate}
\item Quelle valeur donner à $C$ pour que le circuit soit \emph{accordé} sur la station à $\SI{100}{MHz}$ ?
\item Quelle sera alors l'intensité efficace dans le circuit, et quelle tension efficace obtiendra-t-on aux bornes du condensateur, tension que l'on prélèvera pour l'amplifier ?
\item Le récepteur sera-t-il capable de séparer la station visée d'une station concurrente émettant à seulement $\SI{0,5}{MHz}$ d'écart ?
\item Que devient le circuit si l'on tourne le bouton d'accord sans le régler exactement, par exemple en écoutant une station à $\SI{90}{MHz}$ avec le condensateur resté réglé sur $\SI{100}{MHz}$ ?
\end{enumerate}

Par ailleurs, l'atelier du club est alimenté par le réseau ENEO sous une tension efficace $U' = \SI{220}{V}$. L'ensemble des appareils (fer à souder, perceuse, éclairage) absorbe un courant efficace $I' = \SI{10}{A}$ avec un facteur de puissance $\cos\varphi' = 0,7$. Le proviseur s'étonne de la facture d'électricité : Rodrigue veut calculer la puissance réellement consommée et comprendre pourquoi ENEO pénalise les installations à faible facteur de puissance.

\emph{Données :} on rappelle que $\omega = 2\pi f$ ; on prendra $\pi^2 \approx 9,87$.

\courssub{Consignes}

\begin{enumerate}
\item \textbf{Accord du circuit.} Rappeler la condition de résonance d'intensité d'un circuit RLC série. En déduire l'expression littérale de la capacité $C$ en fonction de $L$ et $f_0$, puis calculer sa valeur numérique. Vérifier l'homogénéité du résultat.
\item \textbf{Fonctionnement à la résonance.} Calculer l'impédance du circuit à la résonance, puis l'intensité efficace $I_0$. Calculer le facteur de qualité $Q = \dfrac{L\omega_0}{R}$ et en déduire la tension efficace $U_C$ aux bornes du condensateur. Comparer $U_C$ à $U$ et commenter en deux lignes l'intérêt pratique de ce résultat pour la réception radio.
\item \textbf{Sélectivité.} Donner l'expression de la bande passante $\Delta f$ à $-3$ dB en fonction de $f_0$ et $Q$. Calculer $\Delta f$. Deux stations sont distantes de $\SI{0,5}{MHz}$ : le récepteur de Rodrigue peut-il les séparer ? Justifier et proposer une modification du circuit pour améliorer la sélectivité.
\item \textbf{Désaccord à \SI{90}{MHz}.} Le condensateur garde la valeur trouvée à la question 1. Pour $f = \SI{90}{MHz}$, calculer les réactances $X_L = L\omega$ et $X_C = \dfrac{1}{C\omega}$, puis l'impédance $Z$ et le déphasage $\varphi$ de $u$ par rapport à $i$. Construire le diagramme de Fresnel correspondant (échelle : $\SI{1}{cm}$ pour $\SI{20}{\ohm}$) et préciser si le circuit est capacitif ou inductif. Calculer enfin l'intensité efficace reçue et conclure.
\item \textbf{Facture de l'atelier.} Calculer la puissance moyenne $P$ réellement consommée par l'atelier. Calculer le courant qu'absorberait une installation consommant la même puissance avec $\cos\varphi' = 1$. Expliquer en quelques lignes pourquoi le distributeur d'électricité exige des clients industriels un facteur de puissance élevé.
\end{enumerate}

\courssub{Ressources à mobiliser}

\begin{aretenir}[Boîte à outils de la leçon]
\begin{itemize}
\item Équation du circuit RLC série en régime forcé : $u = Ri + L\dfrac{di}{dt} + \dfrac{q}{C}$, avec $i = \dfrac{dq}{dt}$.
\item Impédance : $Z = \sqrt{R^2 + \left(L\omega - \dfrac{1}{C\omega}\right)^2}$ ; loi d'Ohm en régime sinusoïdal : $U = Z\,I$.
\item Déphasage de $u$ par rapport à $i$ : $\tan\varphi = \dfrac{L\omega - \dfrac{1}{C\omega}}{R}$.
\item Résonance d'intensité : $L\omega_0 = \dfrac{1}{C\omega_0}$, soit $f_0 = \dfrac{1}{2\pi\sqrt{LC}}$ ; alors $Z = R$ (minimale), $I_0 = \dfrac{U}{R}$ (maximale), $\varphi = 0$.
\item Facteur de qualité : $Q = \dfrac{L\omega_0}{R} = \dfrac{1}{RC\omega_0} = \dfrac{f_0}{\Delta f}$ ; surtension : $U_C = Q\,U$.
\item Bande passante à $-3$ dB : $\Delta f = \dfrac{f_0}{Q}$, limitée par les fréquences où $I = \dfrac{I_0}{\sqrt{2}}$.
\item Puissance moyenne : $P = UI\cos\varphi$ ; seule la résistance consomme de la puissance moyenne.
\end{itemize}
\end{aretenir}

\courssub{Démarche et corrigé commenté}

\begin{exoresolu}[Dimensionnement du circuit d'accord et diagnostic de l'atelier]
Énoncé : reprendre les cinq consignes ci-dessus avec les données de la situation ($f_0 = \SI{100}{MHz}$, $L = \SI{0,25}{\micro H}$, $R = \SI{2}{\ohm}$, $U = \SI{10}{mV}$, station perturbatrice à $\pm \SI{0,5}{MHz}$, atelier : $\SI{220}{V}$, $\SI{10}{A}$, $\cos\varphi' = 0,7$).

\tcblower

\textbf{1. Capacité d'accord.} La résonance d'intensité a lieu lorsque les réactances se compensent :
\[
L\omega_0 = \frac{1}{C\omega_0} \quad\Longrightarrow\quad C = \frac{1}{L\omega_0^2} = \frac{1}{4\pi^2 f_0^2 L}.
\]
Analyse dimensionnelle : $[L\omega^2] = \mathrm{H \cdot s^{-2}} = \mathrm{\Omega \cdot s^{-1}}$, donc $[C] = \mathrm{\Omega^{-1} \cdot s} = \mathrm{F}$ : la formule est homogène. Application numérique :
\[
C = \frac{1}{4\pi^2 \times (10^8)^2 \times 0{,}25\times 10^{-6}} = \frac{1}{9{,}87\times 10^{10}} \approx \SI{1,01e-11}{F} \approx \SI{10,1}{pF}.
\]
C'est une capacité très faible, typique des circuits FM : Rodrigue utilisera un condensateur variable de quelques picofarads à quelques dizaines de picofarads.

\textbf{2. Fonctionnement à la résonance.} À la résonance, $Z = R = \SI{2}{\ohm}$, d'où :
\[
I_0 = \frac{U}{R} = \frac{10\times 10^{-3}}{2} = \SI{5,0}{mA}.
\]
Le facteur de qualité vaut :
\[
Q = \frac{L\omega_0}{R} = \frac{0{,}25\times 10^{-6} \times 2\pi \times 10^8}{2} = \frac{157}{2} \approx 78{,}5.
\]
La tension aux bornes du condensateur est alors :
\[
U_C = Q\,U = 78{,}5 \times 10\ \mathrm{mV} \approx \SI{0,79}{V}.
\]
\textit{Commentaire :} grâce au phénomène de \cle{surtension} à la résonance, le circuit fournit au premier étage amplificateur une tension environ $78$ fois plus grande que le signal capté par l'antenne, et ce \emph{uniquement pour la fréquence accordée}. C'est précisément ce double effet — amplification sélective — qui rend la réception radio possible avec un circuit passif.

\textbf{3. Sélectivité.} La bande passante à $-3$ dB est :
\[
\Delta f = \frac{f_0}{Q} = \frac{100\ \mathrm{MHz}}{78{,}5} \approx \SI{1,27}{MHz}.
\]
Or les deux stations ne sont distantes que de $\SI{0,5}{MHz} < \Delta f$ : la station concurrente tombe \emph{à l'intérieur} de la bande passante et sera reçue presque aussi bien que la station visée. \textbf{Le récepteur ne peut pas les séparer.} Pour améliorer la sélectivité, il faut augmenter $Q$, donc \emph{diminuer la résistance} $R$ du circuit (bobine à fil plus gros, soudures soignées) : avec $R = \SI{0,5}{\ohm}$, on aurait $Q \approx 314$ et $\Delta f \approx \SI{0,32}{MHz}$, ce qui permettrait la séparation.

\textbf{4. Désaccord sur \SI{90}{MHz}.} La pulsation vaut $\omega = 2\pi \times 90\times 10^6 \approx \SI{5,65e8}{rad/s}$. Les réactances :
\begin{align*}
X_L &= L\omega = 0{,}25\times 10^{-6} \times 5{,}65\times 10^8 \approx \SI{141}{\ohm},\\
X_C &= \frac{1}{C\omega} = \frac{1}{1{,}01\times 10^{-11} \times 5{,}65\times 10^8} \approx \SI{175}{\ohm}.
\end{align*}
Impédance et déphasage :
\[
Z = \sqrt{R^2 + (X_L - X_C)^2} = \sqrt{2^2 + (141-175)^2} \approx \SI{34}{\ohm},
\qquad
\tan\varphi = \frac{X_L - X_C}{R} = \frac{-34}{2} = -17 \quad\Longrightarrow\quad \varphi \approx -86{,}6^\circ.
\]
Comme $X_C > X_L$, le circuit est \cle{capacitif} : la tension $u$ est en retard sur $i$ (ou $i$ en avance sur $u$). L'intensité efficace reçue n'est plus que :
\[
I = \frac{U}{Z} = \frac{10\ \mathrm{mV}}{34\ \mathrm{\Omega}} \approx \SI{0,29}{mA},
\]
soit $17$ fois moins qu'à la résonance : la station désaccordée est pratiquement inaudible, ce qui confirme le rôle de filtre sélectif du circuit RLC.

\begin{popfigure}
\begin{tikzpicture}[scale=0.8, >=Stealth]
% Échelle : 1 cm = 20 Ω
\def\scale{0.05} % 1/20 = 0.05 cm per ohm
% Vecteurs
\coordinate (O) at (0,0);
\coordinate (R) at (2*\scale, 0); % R = 2
\coordinate (XL) at (2*\scale, 141*\scale); % XL = 141
\coordinate (XC) at (2*\scale, -175*\scale); % XC = -175
\coordinate (X) at (2*\scale, -34*\scale); % X = XL + XC = -34
\coordinate (Z) at (34*\scale, 0); % Z magnitude ~34, but vector is R + jX
% Actually Z vector is (R, X) = (2, -34)
\coordinate (Zvec) at (2*\scale, -34*\scale);

% Axes
\draw[->, thin, popmuted] (0, -9.5) -- (0, 7.5) node[left] {$X$};
\draw[->, thin, popmuted] (-1, 0) -- (2.5, 0) node[below] {$R$};

% Vecteurs Impédance
\draw[popblue, very thick, ->] (O) -- (R) node[midway, below=2pt] {$R = 2\,\Omega$};
\draw[popgreen, very thick, ->] (R) -- (XL) node[midway, right=2pt] {$X_L = 141\,\Omega$};
\draw[poporange, very thick, ->] (R) -- (XC) node[midway, right=2pt] {$X_C = 175\,\Omega$};
\draw[poppurple, very thick, ->] (O) -- (Zvec) node[midway, below right=2pt] {$Z \approx 34\,\Omega$};
\draw[poppurple, dashed, thin] (R) -- (Zvec);
\draw[poppurple, dashed, thin] (XC) -- (Zvec);

% Angle phi
\draw[poppink, thick] (0.3,0) arc (0:-86.6:0.3) node[midway, right=4pt, poppink] {$\varphi \approx -87^\circ$};

% Légendes
\node[popblue] at (1.5*\scale, -1) {$R$};
\node[popgreen] at (1.5*\scale, 3) {$X_L$};
\node[poporange] at (1.5*\scale, -5) {$X_C$};
\node[poppurple] at (1.5*\scale, -2) {$Z$};
\end{tikzpicture}
\legende{Diagramme de Fresnel des impédances à $f = \SI{90}{MHz}$ (échelle : 1 cm pour 20 $\Omega$). Le circuit est capacitif ($X_C > X_L$), $\varphi < 0$.}
\end{popfigure}

\textbf{5. Facture de l'atelier.} La puissance moyenne réellement consommée :
\[
P = U'I'\cos\varphi' = 220 \times 10 \times 0{,}7 = \SI{1540}{W} \approx \SI{1,54}{kW}.
\]
Avec un facteur de puissance égal à $1$, la même puissance ne demanderait que :
\[
I'' = \frac{P}{U'} = \frac{1540}{220} = \SI{7,0}{A} \quad \text{au lieu de } \SI{10}{A}.
\]
\textit{Explication :} à puissance utile égale, un mauvais $\cos\varphi$ impose un courant en ligne plus fort, donc des pertes par effet Joule $RI^2$ plus importantes dans les câbles et transformateurs, et un calibre de compteur plus élevé. C'est pourquoi le distributeur facture des pénalités aux clients dont le facteur de puissance est trop faible ; on le relève en branchant des condensateurs en parallèle sur l'installation.
\end{exoresolu}

\begin{attention}[Pièges rencontrés dans cette activité]
\begin{itemize}
\item Ne pas confondre $\omega$ (en $\mathrm{rad/s}$) et $f$ (en Hz) : la condition de résonance s'écrit $LC\omega_0^2 = 1$, pas $LCf_0^2 = 1$.
\item Les préfixes : $\SI{0,25}{\micro H} = \SI{2,5e-7}{H}$ et $\SI{100}{MHz} = \SI{1e8}{Hz}$ ; une erreur de préfixe fausse $C$ d'un facteur $10^6$.
\item À la résonance, $U_C = QU$ peut dépasser largement $U$ : ce n'est pas une violation de la loi des mailles, car $U_L$ et $U_C$ sont en opposition de phase et se compensent.
\item Dans $P = UI\cos\varphi$, $U$ et $I$ sont des valeurs \emph{efficaces}.
\end{itemize}
\end{attention}

\courssub{Bilan de l'activité}

Cette activité a permis de mobiliser l'intégralité des savoir-faire de la leçon dans un contexte technologique authentique. Sur le plan des calculs, nous avons vu que la condition de résonance $LC\omega_0^2 = 1$ permet de dimensionner un composant réel ($C \approx \SI{10}{pF}$), que le facteur de qualité $Q \approx 78$ quantifie à la fois la surtension exploitable ($U_C \approx \SI{0,79}{V}$) et la sélectivité ($\Delta f \approx \SI{1,3}{MHz}$), et que le diagramme de Fresnel hors résonance rend compte du déphasage ($\varphi \approx -87^\circ$, circuit capacitif) et de l'effondrement de l'intensité reçue. Sur le plan du jugement technique, l'activité a montré qu'un même paramètre — la résistance $R$ — gouverne un compromis fondamental : la diminuer améliore la sélectivité mais exige des composants de meilleure qualité. Enfin, l'étude de l'atelier a relié la leçon à la vie quotidienne camerounaise : le facteur de puissance n'est pas une notion abstraite, c'est un paramètre économique qu'ENEO surveille et que tout technicien doit savoir diagnostiquer et corriger. Tu es désormais prêt à aborder les exercices de type Baccalauréat sur les oscillations forcées, où ces cinq outils — impédance, résonance, bande passante, Fresnel, puissance — reviennent systématiquement.

\newpage

\coursec{Méthodes et exercices résolus}

\courssub{Fiches méthodes et exercices résolus}

\begin{methode}[Écrire l'équation de la tension aux bornes du dipôle RLC série]
Pour établir l'équation différentielle ou l'expression de la tension $u(t)$ aux bornes d'un dipôle RLC série alimenté par un GBF délivrant $u(t) = U_m\cos(\omega t)$ :
\begin{enumerate}
\item \textbf{Flécher le circuit} : choisir un sens positif pour $i$, flécher $u_R$, $u_L$, $u_C$ en convention récepteur et $u$ en convention générateur.
\item \textbf{Écrire la loi d'additivité des tensions} (loi des mailles) :
\[ u(t) = u_R + u_L + u_C = Ri + L\dfrac{di}{dt} + \dfrac{q}{C}. \]
\item \textbf{Choisir la variable} : avec $i = \dfrac{dq}{dt}$, on obtient l'équation en $q$ :
\[ L\ddot{q} + R\dot{q} + \dfrac{q}{C} = U_m\cos(\omega t). \]
\item En régime sinusoïdal forcé permanent, la solution est $i(t) = I_m\cos(\omega t - \varphi)$ ; on cherche alors $I_m$ et $\varphi$ par la \cle{représentation de Fresnel} (jamais en résolvant l'équation différentielle au Bac).
\item \textbf{Vérifier l'homogénéité} : $[L\ddot{q}] = \mathrm{H\cdot A\cdot s^{-1}} = \mathrm{V}$, $[q/C] = \mathrm{C/F} = \mathrm{V}$.
\end{enumerate}
\end{methode}

\begin{exoresolu}[Équation électrique d'un circuit RLC série]
Un dipôle série est constitué d'un résistor $R = \SI{20}{\ohm}$, d'une bobine d'inductance $L = \SI{0,10}{\henry}$ et d'un condensateur $C = \SI{10}{\micro\farad}$. Il est alimenté par un GBF délivrant $u(t) = 12\sqrt{2}\cos(100\pi t)$ (en volts).
\begin{enumerate}
\item Établir l'équation différentielle vérifiée par la charge $q(t)$.
\item Donner l'équation différentielle vérifiée par l'intensité $i(t)$.
\end{enumerate}
\tcblower
\textbf{1.} On flèche les trois dipôles en convention récepteur. La loi d'additivité des tensions s'écrit :
\[ u(t) = u_R + u_L + u_C \quad\text{avec}\quad u_R = Ri,\quad u_L = L\dfrac{di}{dt},\quad u_C = \dfrac{q}{C}. \]
Comme $i = \dot{q}$, on obtient :
\[ L\ddot{q} + R\dot{q} + \dfrac{q}{C} = U_m\cos(\omega t). \]
Application numérique ($\omega = 100\pi \approx \SI{314}{rad.s^{-1}}$, $U_m = 12\sqrt{2} \approx \SI{17,0}{\volt}$) :
\[ 0{,}10\,\ddot{q} + 20\,\dot{q} + 10^{5}\,q = 12\sqrt{2}\cos(100\pi t). \]
Vérification dimensionnelle : $[L\ddot{q}] = \mathrm{H\cdot A/s} = \mathrm{V}$ ; $[q/C] = \mathrm{C/F} = \mathrm{V}$. L'équation est homogène.

\textbf{2.} On dérive membre à membre par rapport au temps, en utilisant $\dot{q} = i$ :
\[ L\ddot{i} + R\dot{i} + \dfrac{i}{C} = -\omega U_m\sin(\omega t), \]
soit numériquement :
\[ 0{,}10\,\ddot{i} + 20\,\dot{i} + 10^{5}\,i = -1200\pi\sqrt{2}\sin(100\pi t). \]
\textbf{Conclusion} : en régime forcé, la charge et l'intensité oscillent à la pulsation $\omega$ imposée par le générateur ; l'amplitude et le déphasage se déterminent par Fresnel.
\end{exoresolu}

\begin{methode}[Construire la représentation de Fresnel ; calculer $Z$ et $\varphi$]
\begin{enumerate}
\item \textbf{Choisir l'axe des phases sur l'intensité} : $i$ est commune à tous les dipôles série, on prend $\vec{I}$ horizontal.
\item \textbf{Placer les vecteurs tensions} :
\begin{itemize}
\item $\vec{U}_R$ : en phase avec $i$ (même direction que $\vec{I}$), norme $RI_m$ ;
\item $\vec{U}_L$ : en \emph{avance} de $\pi/2$ sur $i$ (vers le haut), norme $L\omega I_m$ ;
\item $\vec{U}_C$ : en \emph{retard} de $\pi/2$ sur $i$ (vers le bas), norme $\dfrac{I_m}{C\omega}$.
\end{itemize}
\item \textbf{Somme vectorielle} : $\vec{U} = \vec{U}_R + \vec{U}_L + \vec{U}_C$. Le théorème de Pythagore donne :
\[ U_m = I_m\sqrt{R^2 + \left(L\omega - \dfrac{1}{C\omega}\right)^2} \quad\Longrightarrow\quad Z = \dfrac{U_m}{I_m} = \dfrac{U}{I} = \sqrt{R^2 + \left(L\omega - \dfrac{1}{C\omega}\right)^2}. \]
\item \textbf{Déphasage} de $u$ sur $i$ :
\[ \tan\varphi = \dfrac{L\omega - \dfrac{1}{C\omega}}{R}, \qquad \cos\varphi = \dfrac{R}{Z}. \]
Le signe de $\cos\varphi$ ne suffit pas : utiliser $\tan\varphi$ (ou le signe de $L\omega - \tfrac{1}{C\omega}$) pour lever l'ambiguïté sur le signe de $\varphi$.
\item \textbf{Cas particuliers à connaître} : $R$ seul : $Z = R$, $\varphi = 0$ ; $L$ seule : $Z = L\omega$, $\varphi = +\pi/2$ ; $C$ seul : $Z = \dfrac{1}{C\omega}$, $\varphi = -\pi/2$.
\end{enumerate}
\end{methode}

\begin{exoresolu}[Impédance et déphasage d'un dipôle RLC]
On reprend le circuit de l'exercice précédent : $R = \SI{20}{\ohm}$, $L = \SI{0,10}{\henry}$, $C = \SI{10}{\micro\farad}$, alimenté par $u(t) = 12\sqrt{2}\cos(100\pi t)$.
\begin{enumerate}
\item Calculer l'impédance $Z$ du dipôle.
\item Calculer le déphasage $\varphi$ de $u$ sur $i$ et préciser le caractère du circuit.
\item Écrire l'expression de $i(t)$.
\end{enumerate}
\tcblower
\textbf{1.} Pulsation : $\omega = 100\pi \approx \SI{314,2}{rad.s^{-1}}$.
\[ L\omega = 0{,}10 \times 314{,}2 = \SI{31,4}{\ohm} \; ; \qquad \dfrac{1}{C\omega} = \dfrac{1}{10^{-5}\times 314{,}2} = \SI{318,3}{\ohm}. \]
\[ Z = \sqrt{R^2 + \left(L\omega - \tfrac{1}{C\omega}\right)^2} = \sqrt{20^2 + (31{,}4 - 318{,}3)^2} = \sqrt{400 + 82\,300} \approx \SI{287,6}{\ohm}. \]

\textbf{2.} $\tan\varphi = \dfrac{31{,}4 - 318{,}3}{20} = -14{,}35$, d'où $\varphi \approx -86{,}0^{\circ} \approx \SI{-1,50}{rad}$.
Comme $\varphi < 0$ (car $L\omega < \tfrac{1}{C\omega}$), la tension est en \textbf{retard} sur l'intensité : le circuit est \cle{capacitif}.

\textbf{3.} Intensité maximale : $I_m = \dfrac{U_m}{Z} = \dfrac{12\sqrt{2}}{287{,}6} \approx \SI{5,90e-2}{\ampere}$, soit $I = \dfrac{U}{Z} = \dfrac{12}{287{,}6} \approx \SI{41,7}{\milli\ampere}$ (valeur efficace).
L'intensité est en avance de $|\varphi|$ sur $u$ :
\[ i(t) = 5{,}9\times 10^{-2}\cos\left(100\pi t + 1{,}50\right) \quad (i \text{ en A}). \]
\textbf{Conclusion} : le dipôle est capacitif à cette fréquence ; $i$ est en avance de $86^{\circ}$ sur $u$.
\end{exoresolu}

\begin{methode}[Étudier la résonance d'intensité et exploiter la courbe $I = f(N)$]
\begin{enumerate}
\item \textbf{Condition de résonance} : l'intensité efficace $I = \dfrac{U}{Z}$ est maximale quand $Z$ est minimale, c'est-à-dire quand $L\omega_0 = \dfrac{1}{C\omega_0}$, soit :
\[ \omega_0 = \dfrac{1}{\sqrt{LC}}, \qquad N_0 = \dfrac{1}{2\pi\sqrt{LC}}. \]
\item \textbf{À la résonance} : $Z = R$ (minimum), $I_0 = \dfrac{U}{R}$ (maximum), $\varphi = 0$ ($u$ et $i$ en phase), le circuit est résistif. Il y a aussi résonance de tension aux bornes de $C$.
\item \textbf{Tracer la courbe} : placer le maximum $I_0$ en $N_0$ ; la courbe est d'autant plus aiguë que $R$ est petit.
\item \textbf{Exploitation graphique} : lire $N_0$ au sommet, $I_0$ ; en déduire $R = U/I_0$ si besoin, ou $L$ et $C$ via $N_0$.
\item \textbf{Réflexe Bac} : toujours préciser que $U$ (tension efficace du GBF) est maintenue \emph{constante} pendant qu'on fait varier $N$.
\end{enumerate}
\end{methode}

\begin{exoresolu}[Résonance d'un circuit de réception radio]
Un élève de Yaoundé réalise un circuit série $R = \SI{10}{\ohm}$, $L = \SI{50}{\milli\henry}$, $C$ variable, alimenté par un GBF de tension efficace $U = \SI{6,0}{\volt}$ constante.
\begin{enumerate}
\item Calculer la capacité $C_0$ qui donne la résonance d'intensité à la fréquence $N_0 = \SI{3,18}{\kilo\hertz}$.
\item Calculer l'intensité efficace $I_0$ à la résonance et la tension efficace $U_C$ aux bornes du condensateur. Commenter.
\end{enumerate}
\tcblower
\textbf{1.} À la résonance : $\omega_0 = \dfrac{1}{\sqrt{LC_0}}$ avec $\omega_0 = 2\pi N_0 = 2\pi \times 3{,}18\times 10^{3} \approx \SI{2,00e4}{rad.s^{-1}}$.
\[ C_0 = \dfrac{1}{L\omega_0^{2}} = \dfrac{1}{0{,}050 \times (2{,}00\times 10^{4})^{2}} = \dfrac{1}{2{,}0\times 10^{7}} = \SI{5,0e-8}{\farad} = \SI{50}{\nano\farad}. \]

\textbf{2.} À la résonance, $Z = R$ donc :
\[ I_0 = \dfrac{U}{R} = \dfrac{6{,}0}{10} = \SI{0,60}{\ampere}. \]
Tension efficace aux bornes du condensateur :
\[ U_C = \dfrac{I_0}{C_0\omega_0} = \dfrac{0{,}60}{5{,}0\times 10^{-8}\times 2{,}0\times 10^{4}} = \dfrac{0{,}60}{1{,}0\times 10^{-3}} = \SI{600}{\volt}. \]
\textbf{Commentaire} : $U_C = 100\,U$ ! C'est le phénomène de \cle{surtension} à la résonance : les tensions aux bornes de $L$ et de $C$ peuvent largement dépasser la tension du générateur (elles sont en opposition de phase et se compensent). C'est ce principe qui permet à un poste radio d'amplifier la station sélectionnée, mais il peut aussi détruire un condensateur mal dimensionné.
\textbf{Conclusion} : $C_0 = \SI{50}{\nano\farad}$, $I_0 = \SI{0,60}{\ampere}$, $U_C = \SI{600}{\volt}$ : facteur de surtension égal à 100.
\end{exoresolu}

\begin{methode}[Déterminer la bande passante à $-3$ dB et le facteur de qualité]
\begin{enumerate}
\item \textbf{Définition} : la bande passante est l'intervalle de fréquences $[N_1\,;N_2]$ pour lesquelles $I \geq \dfrac{I_0}{\sqrt{2}}$, soit un gain $G = 20\log\dfrac{I}{I_0} \geq -3$ dB.
\item \textbf{Sur le graphe} : tracer l'horizontale à $I_0/\sqrt{2} \approx 0{,}707\,I_0$ ; lire les fréquences de coupure $N_1$ et $N_2$ aux intersections avec la courbe ; alors :
\[ \Delta N = N_2 - N_1. \]
\item \textbf{Facteur de qualité} :
\[ Q = \dfrac{N_0}{\Delta N} = \dfrac{\omega_0}{\Delta\omega} = \dfrac{L\omega_0}{R} = \dfrac{1}{R}\sqrt{\dfrac{L}{C}} = \dfrac{1}{RC\omega_0}. \]
\item \textbf{Interprétation} : plus $Q$ est grand (donc $R$ petit), plus la résonance est aiguë et le circuit sélectif ; $Q$ mesure aussi la surtension : $Q = \dfrac{U_C}{U}$ à la résonance.
\item \textbf{Vérification} : $Q$ est sans dimension ; vérifier que $N_1 < N_0 < N_2$ et que $Q > 1$ pour un circuit sélectif.
\end{enumerate}
\end{methode}

\begin{exoresolu}[Bande passante et sélectivité]
Pour le circuit précédent ($R = \SI{10}{\ohm}$, $L = \SI{50}{\milli\henry}$, $C_0 = \SI{50}{\nano\farad}$, $N_0 = \SI{3,18}{\kilo\hertz}$, $I_0 = \SI{0,60}{\ampere}$), la courbe de résonance coupe la droite $I = \dfrac{I_0}{\sqrt{2}}$ en $N_1 = \SI{3,02}{\kilo\hertz}$ et $N_2 = \SI{3,34}{\kilo\hertz}$.
\begin{enumerate}
\item Déterminer la bande passante $\Delta N$.
\item Calculer le facteur de qualité $Q$ de deux manières différentes et vérifier la cohérence.
\item On remplace $R$ par $R' = \SI{40}{\ohm}$. Que deviennent $Q$ et la sélectivité ?
\end{enumerate}
\tcblower
\textbf{1.} Bande passante : $\Delta N = N_2 - N_1 = 3{,}34 - 3{,}02 = \SI{0,32}{\kilo\hertz} = \SI{320}{\hertz}$.

\textbf{2.} Première méthode (graphique) :
\[ Q = \dfrac{N_0}{\Delta N} = \dfrac{3180}{320} \approx 9{,}9. \]
Deuxième méthode (calcul avec la résistance nominale $R = \SI{10}{\ohm}$) :
\[ Q = \dfrac{L\omega_0}{R} = \dfrac{0{,}050 \times 2{,}0\times 10^{4}}{10} = \dfrac{1000}{10} = 100. \]
Cette valeur contredit le résultat graphique. La relation exacte est $\Delta\omega = \dfrac{R_{\text{eq}}}{L}$, soit $\Delta N = \dfrac{R_{\text{eq}}}{2\pi L}$. Avec $\Delta N = \SI{320}{\hertz}$, on trouve une résistance équivalente série $R_{\text{eq}} = 2\pi L \Delta N \approx \SI{100}{\ohm}$ (résistance de la bobine incluse). Retenons la valeur cohérente avec les mesures : $Q = \dfrac{N_0}{\Delta N} \approx 9{,}9 \approx 10$.
Vérification par surtension : $Q = \dfrac{U_C}{U}$ à la résonance donnerait $U_C = QU = 10 \times 6{,}0 = \SI{60}{\volt}$ (contre $\SI{600}{\volt}$ pour le circuit idéal).

\textbf{3.} Avec $R' = 4R = \SI{40}{\ohm}$ (résistance externe augmentée) : $Q' = \dfrac{Q}{4} \approx 2{,}5$ et $\Delta N' = 4\,\Delta N \approx \SI{1,28}{\kilo\hertz}$.
\textbf{Conclusion} : augmenter $R$ élargit la bande passante et diminue $Q$ : le circuit devient moins sélectif (résonance plus « floue »), ce qui est mauvais pour sélectionner une station radio mais parfois recherché en électronique audio.
\end{exoresolu}

\begin{methode}[Calculer la puissance moyenne et le facteur de puissance]
\begin{enumerate}
\item \textbf{Formule fondamentale} (valeurs efficaces) :
\[ P = UI\cos\varphi = RI^{2} = \dfrac{U^{2}R}{Z^{2}}. \]
Seul le résistor consomme de la puissance moyenne ; $L$ et $C$ parfaits n'en consomment pas.
\item \textbf{Facteur de puissance} : $k = \cos\varphi = \dfrac{R}{Z} = \dfrac{P}{UI}$, compris entre 0 et 1.
\item \textbf{À la résonance} : $\varphi = 0$, $\cos\varphi = 1$, la puissance est maximale : $P_{\max} = UI = \dfrac{U^{2}}{R}$.
\item \textbf{Piège} : ne jamais écrire $P = UI$ sans le $\cos\varphi$ hors résonance ; ne pas utiliser les valeurs maximales dans $UI\cos\varphi$ (ou alors diviser par 2 : $P = \tfrac{1}{2}U_m I_m \cos\varphi$).
\item \textbf{Enjeu pratique} : un faible $\cos\varphi$ impose des courants plus forts pour la même puissance utile, d'où des pertes en ligne : c'est pourquoi ENEO facture les mauvais facteurs de puissance aux industriels.
\end{enumerate}
\end{methode}

\begin{exoresolu}[Facture d'énergie d'un atelier de Douala]
Un moteur d'atelier est modélisé par un dipôle série $R = \SI{40}{\ohm}$, $L = \SI{0,20}{\henry}$, alimenté sous $U = \SI{220}{\volt}$, $N = \SI{50}{\hertz}$.
\begin{enumerate}
\item Calculer l'impédance $Z$, l'intensité efficace $I$ et le facteur de puissance.
\item Calculer la puissance moyenne $P$ consommée.
\item On ajoute un condensateur en série pour ramener $\cos\varphi$ à 1 (résonance). Calculer $C$ et la nouvelle puissance consommée.
\end{enumerate}
\tcblower
\textbf{1.} $\omega = 2\pi N = 100\pi \approx \SI{314,2}{rad.s^{-1}}$ ; $L\omega = 0{,}20 \times 314{,}2 = \SI{62,8}{\ohm}$.
\[ Z = \sqrt{R^{2} + (L\omega)^{2}} = \sqrt{40^{2} + 62{,}8^{2}} = \sqrt{1600 + 3944} = \sqrt{5544} \approx \SI{74,5}{\ohm}. \]
\[ I = \dfrac{U}{Z} = \dfrac{220}{74{,}5} \approx \SI{2,95}{\ampere} \; ; \qquad \cos\varphi = \dfrac{R}{Z} = \dfrac{40}{74{,}5} \approx 0{,}537. \]

\textbf{2.} Puissance moyenne :
\[ P = UI\cos\varphi = 220 \times 2{,}95 \times 0{,}537 \approx \SI{349}{\watt}. \]
Vérification : $P = RI^{2} = 40 \times 2{,}95^{2} \approx \SI{348}{\watt}$. Cohérent (écarts dus aux arrondis).

\textbf{3.} Résonance : $C = \dfrac{1}{L\omega^{2}} = \dfrac{1}{0{,}20 \times (314{,}2)^{2}} = \dfrac{1}{19\,739} \approx \SI{5,07e-5}{\farad} \approx \SI{51}{\micro\farad}$.
Alors $Z = R$, $I' = \dfrac{U}{R} = \dfrac{220}{40} = \SI{5,5}{\ampere}$, et :
\[ P' = \dfrac{U^{2}}{R} = \dfrac{220^{2}}{40} = \SI{1210}{\watt}. \]
\textbf{Conclusion} : le condensateur relève le facteur de puissance à 1 ; la puissance utile passe de $\SI{349}{\watt}$ à $\SI{1210}{\watt}$. En pratique industrielle, on place le condensateur en \emph{parallèle} pour ne pas augmenter le courant dans le moteur, mais l'effet sur $\cos\varphi$ est le même.
\end{exoresolu}

\begin{methode}[Mesurer un déphasage à l'oscilloscope]
\begin{enumerate}
\item \textbf{Branchement} : visualiser $u(t)$ sur la voie A ; pour visualiser l'image de $i(t)$, afficher $u_R(t) = Ri(t)$ sur la voie B (en phase avec $i$). Attention à la masse commune : placer $R$ au point du circuit relié à la masse.
\item \textbf{Réglages} : mêmes sensibilités verticales si possible, balayage permettant d'observer une à deux périodes.
\item \textbf{Mesure} : repérer le décalage horizontal $\Delta t$ entre deux passages par zéro dans le même sens (ou deux maxima), et la période $T$ :
\[ |\varphi| = 2\pi\,\dfrac{\Delta t}{T} = \omega\,\Delta t \quad (\text{en rad}), \qquad \text{ou } |\varphi| = 360^{\circ}\times\dfrac{\Delta t}{T}. \]
\item \textbf{Sens du déphasage} : la courbe dont le maximum apparaît \emph{en premier} (à gauche) est en \emph{avance} sur l'autre.
\item \textbf{Méthode de Lissajous} (balayage supprimé, X-Y) : ellipse dont l'axe permet de calculer $\sin\varphi = \dfrac{b}{B}$ ; droite $\Rightarrow \varphi = 0$ ou $\pi$ ; cercle $\Rightarrow |\varphi| = \pi/2$ avec amplitudes égales.
\end{enumerate}
\end{methode}

\begin{exoresolu}[Exploitation d'un écran d'oscilloscope]
Sur l'écran d'un oscilloscope, la voie A affiche $u(t)$ aux bornes d'un dipôle RLC et la voie B affiche $u_R(t)$. Balayage : $\SI{1,0}{ms/div}$. La période occupe $8{,}0$ divisions ; le maximum de $u_R$ apparaît $1{,}5$ division après celui de $u$.
\begin{enumerate}
\item Calculer la fréquence $N$ de la tension.
\item Calculer le déphasage $\varphi$ de $u$ sur $i$ et préciser le caractère du circuit.
\item Le circuit est-il en résonance ? Justifier.
\end{enumerate}
\tcblower
\textbf{1.} Période : $T = 8{,}0 \times 1{,}0 = \SI{8,0}{ms} = \SI{8,0e-3}{s}$, d'où :
\[ N = \dfrac{1}{T} = \dfrac{1}{8{,}0\times 10^{-3}} = \SI{125}{\hertz}. \]

\textbf{2.} Décalage temporel : $\Delta t = 1{,}5 \times 1{,}0 = \SI{1,5}{ms}$.
\[ |\varphi| = 2\pi\dfrac{\Delta t}{T} = 2\pi \times \dfrac{1{,}5}{8{,}0} = \dfrac{3\pi}{8} \approx \SI{1,18}{rad} \quad (\approx 67{,}5^{\circ}). \]
Le maximum de $u_R$ (donc de $i$) apparaît \emph{après} celui de $u$ : $i$ est en \textbf{retard} sur $u$, donc $u$ est en avance sur $i$ : $\varphi = +\dfrac{3\pi}{8} \approx +\SI{1,18}{rad}$. Le circuit est \cle{inductif}.

\textbf{3.} À la résonance, $\varphi = 0$ : les deux courbes seraient en phase. Ici $\varphi \neq 0$, donc le circuit n'est \textbf{pas} en résonance ; de plus $\varphi > 0$ indique $L\omega > \dfrac{1}{C\omega}$, c'est-à-dire $N > N_0$.
\textbf{Conclusion} : $N = \SI{125}{\hertz}$, $\varphi = +67{,}5^{\circ}$, circuit inductif fonctionnant au-dessus de sa fréquence de résonance.
\end{exoresolu}

\begin{attention}[Les 5 erreurs qui coûtent des points au Bac]
\begin{enumerate}
\item \textbf{Oublier le $\cos\varphi$ dans la puissance.} Écrire $P = UI$ est faux en dehors de la résonance : la seule formule générale est $P = UI\cos\varphi = RI^{2}$. Et attention : $U$ et $I$ sont des valeurs \emph{efficaces}.
\item \textbf{Conclure sur le signe de $\varphi$ avec $\cos\varphi$ seul.} $\cos\varphi = R/Z$ est toujours positif : il ne permet pas de savoir si le circuit est inductif ou capacitif. Il faut utiliser $\tan\varphi$ ou comparer $L\omega$ et $\dfrac{1}{C\omega}$.
\item \textbf{Confondre valeurs maximales et efficaces.} $Z = \dfrac{U_m}{I_m} = \dfrac{U}{I}$ fonctionne avec les deux, mais mélanger $U_m$ et $I$ efficace dans un même calcul est une erreur. Rappel : $U_m = U\sqrt{2}$.
\item \textbf{Mal définir la bande passante.} Les fréquences de coupure correspondent à $I = \dfrac{I_0}{\sqrt{2}}$ (et non $I_0/2$) ; le « $-3$ dB » vient de $20\log(1/\sqrt{2}) \approx -3$ dB. Et $Q = \dfrac{N_0}{\Delta N}$, sans dimension.
\item \textbf{Se tromper de sens dans le déphasage à l'oscilloscope.} La courbe dont le maximum apparaît en premier (à gauche) est en avance. Dire « $u$ est en retard » alors que le circuit est inductif ($\varphi > 0$, $u$ en avance) est la faute classique ; vérifier la cohérence avec le signe de $L\omega - \tfrac{1}{C\omega}$.
\end{enumerate}
\end{attention}

\newpage

\coursec{Exercices}

\courssub{Je vérifie mes connaissances}

\begin{exercice}[QCM : une seule bonne réponse]
Pour chaque question, choisir la bonne réponse et justifier brièvement.
\begin{enumerate}
\item Dans un circuit RLC série alimenté par une tension sinusoïdale, l'impédance s'écrit :
\begin{enumerate}
\item $Z = R + L\omega + \dfrac{1}{C\omega}$ ;
\item $Z = \sqrt{R^2 + \left(L\omega - \dfrac{1}{C\omega}\right)^2}$ ;
\item $Z = \sqrt{R^2 + \left(L\omega + \dfrac{1}{C\omega}\right)^2}$ ;
\item $Z = \dfrac{1}{\sqrt{R^2 + \left(L\omega - \dfrac{1}{C\omega}\right)^2}}$.
\end{enumerate}
\item À la résonance d'intensité :
\begin{enumerate}
\item l'intensité efficace est minimale ;
\item le déphasage $\varphi$ entre $u$ et $i$ vaut $\dfrac{\pi}{2}$ ;
\item l'impédance du circuit est égale à $R$ ;
\item la pulsation vaut $\omega_0 = \sqrt{LC}$.
\end{enumerate}
\item La puissance moyenne consommée par un dipôle en régime sinusoïdal est :
\begin{enumerate}
\item $P = UI$ ;
\item $P = UI\sin\varphi$ ;
\item $P = UI\cos\varphi$ ;
\item $P = \dfrac{UI}{2}$.
\end{enumerate}
\item Le facteur de qualité d'un circuit RLC série est :
\begin{enumerate}
\item $Q = \dfrac{R}{L\omega_0}$ ;
\item $Q = \dfrac{L\omega_0}{R}$ ;
\item $Q = \dfrac{R}{L}\omega_0$ ;
\item $Q = RC\omega_0$.
\end{enumerate}
\end{enumerate}
\end{exercice}

\begin{exercice}[Vrai ou faux ?]
Répondre par vrai ou faux en justifiant chaque réponse par une phrase ou un calcul.
\begin{enumerate}
\item Un condensateur seul, alimenté en sinusoïdal, consomme une puissance moyenne nulle.
\item Si $L\omega > \dfrac{1}{C\omega}$, le circuit RLC série est capacitif.
\item La bande passante à $-3$ dB est l'intervalle de fréquences pour lesquelles $I \geq \dfrac{I_{\max}}{\sqrt{2}}$.
\item Plus la résistance $R$ du circuit est grande, plus la résonance est aiguë.
\item À l'oscilloscope, si la tension $u$ est en avance sur l'intensité $i$, le circuit est inductif.
\end{enumerate}
\end{exercice}

\begin{exercice}[Définitions et expressions]
\begin{enumerate}
\item Définir l'\cle{impédance} d'un dipôle et donner son unité.
\item Définir le \cle{déphasage} $\varphi$ entre la tension $u$ et l'intensité $i$. Que signifie $\varphi > 0$ ?
\item Donner l'expression de la pulsation propre $\omega_0$ d'un circuit RLC série, puis celle de la fréquence propre $N_0$. Vérifier par analyse dimensionnelle que $\dfrac{1}{\sqrt{LC}}$ est homogène à l'inverse d'un temps.
\item Rappeler la relation entre le facteur de qualité $Q$, la fréquence propre $N_0$ et la bande passante $\Delta N$.
\end{enumerate}
\end{exercice}

\begin{exercice}[Calculs directs]
Un dipôle RLC série est alimenté par une tension sinusoïdale de pulsation $\omega = 500$ rad$\cdot$s$^{-1}$. On donne $R = 40\ \Omega$, $L = 0{,}1$ H et $C = 10\ \mu$F.
\begin{enumerate}
\item Calculer les réactances $X_L = L\omega$ et $X_C = \dfrac{1}{C\omega}$.
\item En déduire l'impédance $Z$ du circuit.
\item Calculer $\tan\varphi$ puis le déphasage $\varphi$ en degrés. Le circuit est-il inductif ou capacitif ?
\item La tension efficace aux bornes du dipôle vaut $U = 24$ V. Calculer l'intensité efficace $I$ et la puissance moyenne $P$ consommée.
\end{enumerate}
\end{exercice}

\courssub{Je m'entraîne}

\begin{exercice}[Dipôle RLC série complet]
Un dipôle RLC série ($R = 50\ \Omega$, $L = 0{,}2$ H, $C = 5\ \mu$F) est alimenté par la tension $u(t) = 20\sqrt{2}\cos(100\pi t)$ (en volts).
\begin{enumerate}
\item Déterminer la tension efficace $U$, la pulsation $\omega$ et la fréquence $N$ de la tension d'alimentation.
\item Calculer l'impédance $Z$ du dipôle et le déphasage $\varphi$ de $u$ par rapport à $i$.
\item Calculer l'intensité efficace $I$, puis écrire l'expression de $i(t)$.
\item Calculer les tensions efficaces $U_R$, $U_L$ et $U_C$. Vérifier que $U \neq U_R + U_L + U_C$ et commenter.
\item Construire le diagramme de Fresnel des tensions (échelle : 1 cm pour 5 V).
\item Calculer la puissance moyenne $P$ consommée par le dipôle de deux manières différentes.
\end{enumerate}
\end{exercice}

\begin{exercice}[Surtension aux bornes du condensateur]
\begin{enumerate}
\item Montrer qu'à la résonance d'intensité d'un circuit RLC série, la tension efficace aux bornes du condensateur vaut $U_C = Q\,U$, où $U$ est la tension efficace d'alimentation et $Q$ le facteur de qualité.
\item Un circuit résonne sous $U = 12$ V avec un facteur de qualité $Q = 25$. Calculer $U_C$.
\item Le condensateur utilisé supporte au maximum une tension efficace de $250$ V. Y a-t-il danger ? Commenter en une phrase.
\end{enumerate}
\end{exercice}

\begin{exercice}[Bande passante et facteur de qualité]
On considère un circuit RLC série alimenté par un GBF délivrant une tension de valeur efficace constante $U = 6$ V.
\begin{enumerate}
\item En partant de la condition $I = \dfrac{I_{\max}}{\sqrt{2}}$, démontrer que la bande passante en pulsation vérifie $\Delta\omega = \dfrac{R}{L}$.
\item En déduire l'expression du facteur de qualité $Q = \dfrac{\omega_0}{\Delta\omega}$ et montrer que $Q = \dfrac{1}{R}\sqrt{\dfrac{L}{C}}$.
\item Application : avec $R = 20\ \Omega$, $L = 50$ mH et $C = 100$ nF, calculer $N_0$, $\Delta N$ et $Q$.
\end{enumerate}
\end{exercice}

\begin{exercice}[Relevage du facteur de puissance à Douala]
Un petit atelier de menuiserie à Douala (moteur électrique + lampes) est alimenté sous $U = 220$ V, $N = 50$ Hz. Il absorbe une intensité efficace $I = 8$ A avec un facteur de puissance $\cos\varphi = 0{,}75$.
\begin{enumerate}
\item Calculer la puissance moyenne $P$ consommée par l'atelier.
\item Calculer la puissance apparente $S = UI$ et la puissance réactive $Q_r = UI\sin\varphi$.
\item Après installation d'un condensateur de compensation, le facteur de puissance passe à $\cos\varphi' = 0{,}95$ pour la même puissance utile $P$. Calculer la nouvelle intensité efficace $I'$ absorbée.
\item Calculer la réduction relative de l'intensité en pourcentage. Conclure sur l'intérêt, pour ENEO et pour l'abonné, du relevage du facteur de puissance (échauffement des lignes, pertes par effet Joule).
\end{enumerate}
\end{exercice}

\courssub{Je me prépare au Bac}

\begin{exercice}[Type Bac : étude complète d'un dipôle RLC série]
\textbf{Données :} $R = 100\ \Omega$ ; $L = 0{,}5$ H ; $C = 5\ \mu$F ; la tension d'alimentation est $u(t) = 30\sqrt{2}\cos(1000\,t)$ (en volts) ; $\pi \approx 3{,}14$.

Un technicien du lycée technique de Bafoussam réalise un dipôle série constitué d'un résistor de résistance $R$, d'une bobine d'inductance $L$ (de résistance négligeable) et d'un condensateur de capacité $C$. Il alimente l'ensemble par un GBF délivrant la tension $u(t)$ ci-dessus.
\begin{enumerate}
\item \textbf{Équation électrique.} Établir, à partir de la loi d'additivité des tensions, l'équation différentielle vérifiée par l'intensité $i(t)$ dans le circuit.
\item \textbf{Impédance et déphasage.}
\begin{enumerate}
\item Calculer les réactances $L\omega$ et $\dfrac{1}{C\omega}$, puis l'impédance $Z$ du dipôle.
\item Calculer le déphasage $\varphi$ de $u$ par rapport à $i$ (valeur en radians puis en degrés). Préciser la nature (inductive ou capacitive) du circuit.
\end{enumerate}
\item \textbf{Intensité.} Calculer l'intensité efficace $I$ et écrire l'expression de $i(t)$.
\item \textbf{Diagramme de Fresnel.} Calculer $U_R$, $U_L$ et $U_C$, puis construire le diagramme de Fresnel des tensions (échelle : 1 cm pour 5 V). Retrouver graphiquement la valeur de $Z$.
\item \textbf{Puissance.} Calculer la puissance moyenne $P$ consommée par le dipôle. Quelle fraction de cette puissance est dissipée dans la bobine ? Justifier.
\item \textbf{Résonance.} On fait maintenant varier la fréquence du GBF, la tension efficace restant égale à $30$ V.
\begin{enumerate}
\item Calculer la fréquence de résonance $N_0$ du circuit.
\item Que valent alors l'intensité efficace $I_{\max}$ et le déphasage $\varphi$ ?
\item Calculer le facteur de qualité $Q$ et la tension efficace $U_C$ aux bornes du condensateur à la résonance. Commenter.
\end{enumerate}
\end{enumerate}
\end{exercice}

\begin{exercice}[Type Bac : exploitation d'une courbe de résonance]
\textbf{Données :} $C = 100$ nF ; tension efficace maintenue constante $U = 6$ V ; $\sqrt{2} \approx 1{,}41$ ; $\pi \approx 3{,}14$.

Dans le cadre d'un TP au lycée de Ngaoundéré, un groupe d'élèves alimente un circuit RLC série par un GBF et relève l'intensité efficace $I$ en fonction de la fréquence $N$. Ils obtiennent la courbe ci-dessous.

\begin{popfigure}
\begin{center}
\begin{tikzpicture}
\begin{axis}[width=0.85\linewidth,height=6cm,
xlabel={$N$ (kHz)},ylabel={$I$ (mA)},
xmin=0,xmax=6,ymin=0,ymax=70,
grid=major,grid style={popline!60},
xtick={0,1,2,3,4,5,6},ytick={0,10,20,30,40,50,60},
label style={font=\small},tick label style={font=\small}]
\addplot[popcurve,domain=0.9:3.6,samples=300]{6000/sqrt(max(0,100^2+(2*pi*x*1000*0.05-1/(2*pi*x*1000*1e-7))^2))};
\addplot[poporange,dashed,domain=0:6]{42.43};
\addplot[popmuted,dashed] coordinates {(2.25,0) (2.25,60)};
\node[font=\small,popdark] at (axis cs:4.6,46) {$\dfrac{I_{\max}}{\sqrt{2}}$};
\node[font=\small,popdark] at (axis cs:2.95,64) {$I_{\max}$};
\end{axis}
\end{tikzpicture}
\end{center}
\legende{Courbe de résonance $I = f(N)$ obtenue par le groupe d'élèves.}
\end{popfigure}

\begin{enumerate}
\item \textbf{Lecture graphique.}
\begin{enumerate}
\item Déterminer graphiquement la fréquence de résonance $N_0$ et l'intensité maximale $I_{\max}$.
\item Calculer $\dfrac{I_{\max}}{\sqrt{2}}$, puis déterminer graphiquement les fréquences $N_1$ et $N_2$ limitant la bande passante à $-3$ dB.
\item En déduire la bande passante $\Delta N$ et le facteur de qualité $Q$ du circuit.
\end{enumerate}
\item \textbf{Exploitation.}
\begin{enumerate}
\item Montrer qu'à la résonance, $I_{\max} = \dfrac{U}{R}$. En déduire la résistance $R$ du circuit.
\item À partir de l'expression de $N_0$, calculer l'inductance $L$ de la bobine.
\item Retrouver la valeur de $Q$ à l'aide de l'expression $Q = \dfrac{1}{R}\sqrt{\dfrac{L}{C}}$. Comparer avec la valeur graphique.
\end{enumerate}
\item \textbf{Interprétation.} Le groupe remplace le résistor par un autre de résistance deux fois plus faible. Décrire qualitativement (sans calcul) les modifications de la courbe de résonance : hauteur du pic, largeur de la bande passante, sélectivité du circuit. Citer une application pratique de cette sélectivité dans un récepteur radio (par exemple pour capter CRTV ou une radio communautaire).
\end{enumerate}
\end{exercice}

\begin{exercice}[Type Bac : mesures à l'oscilloscope]
\textbf{Données :} voie A (tension $u$ aux bornes du dipôle RLC) : sensibilité $5$ V/div ; voie B (tension aux bornes de $R$, image de $i$) : sensibilité $5$ V/div ; balayage : $0{,}5$ ms/div ; $R = 50\ \Omega$ ; $\pi \approx 3{,}14$.

Un élève de Terminale C du lycée de Maroua branche un oscilloscope sur un circuit RLC série alimenté par un GBF : la voie A visualise la tension totale $u$, la voie B visualise la tension $u_R$ aux bornes du résistor. Il obtient l'oscillogramme ci-dessous.

\begin{popfigure}
\begin{center}
\begin{tikzpicture}[scale=0.9]
\draw[step=0.8,popline!50,very thin] (-4,-2.4) grid (4,2.4);
\draw[popdark,thick,->] (-4.2,0) -- (4.3,0);
\draw[popdark,thick,->] (0,-2.6) -- (0,2.7);
\draw[popblue,thick,domain=-4:4,samples=200] plot (\x,{1.6*sin(90*\x)});
\draw[poporange,thick,domain=-4:4,samples=200] plot (\x,{1.29*sin(90*\x - 36)});
\node[popblue,font=\small] at (-2.9,2.1) {voie A : $u$};
\node[poporange,font=\small] at (-2.9,-1.9) {voie B : $u_R$};
\foreach \x in {-4,-3,-2,-1,1,2,3,4} \draw[popdark] (\x,0.06) -- (\x,-0.06);
\foreach \y in {-2,-1,1,2} \draw[popdark] (0.06,\y*0.8) -- (-0.06,\y*0.8);
\draw[popmuted,dashed] (1,0) -- (1,1.6);
\draw[popmuted,dashed] (1.4,0) -- (1.4,1.29);
\draw[popdark,<->] (1,-2.55) -- (1.4,-2.55) node[midway,below,font=\scriptsize] {$\Delta t$};
\end{tikzpicture}
\end{center}
\legende{Oscillogramme : voie A (bleu, 5 V/div), voie B (orange, 5 V/div), balayage 0,5 ms/div.}
\end{popfigure}

\begin{enumerate}
\item \textbf{Mesures.}
\begin{enumerate}
\item Mesurer sur l'écran l'amplitude $U_m$ de la tension $u$ et l'amplitude $U_{Rm}$ de la tension $u_R$. En déduire les valeurs efficaces $U$ et $U_R$.
\item Mesurer la période $T$ et en déduire la fréquence $N$ et la pulsation $\omega$ du GBF.
\item Mesurer le décalage temporel $\Delta t$ entre les deux courbes. En déduire la valeur absolue du déphasage $|\varphi|$ entre $u$ et $i$ (en radians puis en degrés).
\end{enumerate}
\item \textbf{Exploitation.}
\begin{enumerate}
\item Justifier que la voie B donne l'image de l'intensité $i(t)$. Calculer l'intensité efficace $I$ dans le circuit.
\item Calculer l'impédance $Z$ du dipôle RLC.
\item La courbe de la voie A est-elle en avance ou en retard sur celle de la voie B ? En déduire si le circuit est inductif ou capacitif, et le signe de $\varphi$.
\item Calculer la puissance moyenne $P$ consommée par le dipôle.
\end{enumerate}
\item \textbf{Pour aller plus loin.} Sachant que $L = 0{,}1$ H, déterminer la capacité $C$ du condensateur (on utilisera l'expression de $\tan\varphi$). Que faudrait-il faire pour observer la résonance d'intensité ?
\end{enumerate}
\end{exercice}

\newpage

\coursec{Corrigés des exercices}

\coursec{Exercices corrigés et problèmes types Bac — Oscillations électriques forcées en régime sinusoïdal}

\begin{corrige}[Exercice 1 — QCM : une seule bonne réponse]
\begin{enumerate}
\item \textbf{Réponse b.} Dans la construction de Fresnel, la tension aux bornes du résistor ($RI$, en phase avec $i$) est portée horizontalement, tandis que $U_L = L\omega I$ (en avance de $\dfrac{\pi}{2}$) et $U_C = \dfrac{I}{C\omega}$ (en retard de $\dfrac{\pi}{2}$) sont verticales et de sens opposés. Le théorème de Pythagore donne $U = I\sqrt{R^2 + \left(L\omega - \dfrac{1}{C\omega}\right)^2}$, d'où $Z = \sqrt{R^2 + \left(L\omega - \dfrac{1}{C\omega}\right)^2}$. La réponse a additionne algébriquement des grandeurs en quadrature ; la réponse d est l'admittance $Y = \dfrac{1}{Z}$.
\item \textbf{Réponse c.} À la résonance, $L\omega_0 = \dfrac{1}{C\omega_0}$, donc le terme réactif s'annule et $Z = R$ : l'impédance est minimale, l'intensité efficace est maximale et $\varphi = 0$. La pulsation propre est $\omega_0 = \dfrac{1}{\sqrt{LC}}$ et non $\sqrt{LC}$.
\item \textbf{Réponse c.} $P = UI\cos\varphi$ : seule la composante de la tension en phase avec l'intensité transporte de la puissance moyenne ; $UI$ est la puissance apparente.
\item \textbf{Réponse b.} $Q = \dfrac{L\omega_0}{R} = \dfrac{1}{RC\omega_0} = \dfrac{1}{R}\sqrt{\dfrac{L}{C}}$. Plus $R$ est petite, plus $Q$ est grand et plus la résonance est aiguë.
\end{enumerate}
\end{corrige}

\begin{corrige}[Exercice 2 — Vrai ou faux ?]
\begin{enumerate}
\item \textbf{Vrai.} Pour un condensateur parfait, l'intensité est en avance de $\dfrac{\pi}{2}$ sur la tension : $\varphi = -\dfrac{\pi}{2}$, donc $P = UI\cos\varphi = 0$. Le condensateur emmagasine puis restitue l'énergie électrostatique sans en consommer en moyenne.
\item \textbf{Faux.} Si $L\omega > \dfrac{1}{C\omega}$, alors $\tan\varphi = \dfrac{L\omega - \frac{1}{C\omega}}{R} > 0$, donc $\varphi > 0$ : la tension est en avance sur l'intensité, le circuit est \cle{inductif} (et non capacitif).
\item \textbf{Vrai.} Aux fréquences de coupure $N_1$ et $N_2$, $I = \dfrac{I_{\max}}{\sqrt{2}}$, ce qui correspond à une atténuation de $20\log\sqrt{2} \approx 3$ dB ; entre $N_1$ et $N_2$, $I \geq \dfrac{I_{\max}}{\sqrt{2}}$.
\item \textbf{Faux.} $Q = \dfrac{L\omega_0}{R}$ : si $R$ augmente, $Q$ diminue et la bande passante $\Delta N = \dfrac{N_0}{Q}$ s'élargit : la résonance devient \emph{floue}. C'est une petite résistance qui donne une résonance aiguë.
\item \textbf{Vrai.} $u$ en avance sur $i$ signifie $\varphi > 0$, donc $L\omega > \dfrac{1}{C\omega}$ : le circuit a un comportement inductif.
\end{enumerate}
\end{corrige}

\begin{corrige}[Exercice 3 — Définitions et expressions]
\begin{enumerate}
\item L'\cle{impédance} $Z$ d'un dipôle est le quotient de la tension efficace à ses bornes par l'intensité efficace qui le traverse : $Z = \dfrac{U}{I}$. Elle s'exprime en ohms ($\Omega$).
\item Le \cle{déphasage} $\varphi$ est la phase de $u(t)$ diminuée de la phase de $i(t)$ : $\varphi = \varphi_u - \varphi_i$. Si $\varphi > 0$, la tension est en avance sur l'intensité : le dipôle a un comportement inductif.
\item $\omega_0 = \dfrac{1}{\sqrt{LC}}$ et $N_0 = \dfrac{\omega_0}{2\pi} = \dfrac{1}{2\pi\sqrt{LC}}$. Analyse dimensionnelle : d'après $u_L = L\dfrac{di}{dt}$, $[L] = \dfrac{[U]\,[T]}{[I]}$ ; d'après $i = C\dfrac{du_C}{dt}$, $[C] = \dfrac{[I]\,[T]}{[U]}$. Donc $[LC] = [T]^2$ et $\left[\dfrac{1}{\sqrt{LC}}\right] = [T]^{-1}$ : c'est bien l'inverse d'un temps, c'est-à-dire une pulsation (rad$\cdot$s$^{-1}$).
\item $Q = \dfrac{N_0}{\Delta N} = \dfrac{\omega_0}{\Delta\omega}$ : le facteur de qualité mesure l'acuité de la résonance.
\end{enumerate}
\end{corrige}

\begin{corrige}[Exercice 4 — Calculs directs]
\begin{enumerate}
\item Réactances :
\[ X_L = L\omega = \SI{0,1}{\henry} \times \SI{500}{s^{-1}} = \SI{50}{\ohm} \qquad ; \qquad X_C = \dfrac{1}{C\omega} = \dfrac{1}{\SI{10^{-5}}{\farad} \times \SI{500}{s^{-1}}} = \SI{200}{\ohm}. \]
\item Impédance :
\[ Z = \sqrt{R^2 + (X_L - X_C)^2} = \sqrt{\SI{40}{\ohm}^2 + (\SI{50}{\ohm} - \SI{200}{\ohm})^2} = \sqrt{1600 + 22500} = \sqrt{24100} \approx \SI{155}{\ohm}. \]
\item Déphasage :
\[ \tan\varphi = \dfrac{X_L - X_C}{R} = \dfrac{50 - 200}{40} = -3{,}75 \quad \Rightarrow \quad \varphi \approx -75^{\circ}. \]
Comme $\varphi < 0$ (car $X_C > X_L$), le circuit est \cle{capacitif} : l'intensité est en avance sur la tension.
\item Intensité efficace : $I = \dfrac{U}{Z} = \dfrac{\SI{24}{\volt}}{\SI{155}{\ohm}} \approx \SI{0,155}{\ampere}$. Puissance moyenne :
\[ P = UI\cos\varphi = \SI{24}{\volt} \times \SI{0,155}{\ampere} \times \cos 75^{\circ} \approx \SI{0,96}{\watt}. \]
Vérification : $P = RI^2 = \SI{40}{\ohm} \times (\SI{0,155}{\ampere})^2 \approx \SI{0,96}{\watt}$. Les deux calculs concordent : seul le résistor dissipe de la puissance.
\end{enumerate}
\end{corrige}

\begin{corrige}[Exercice 5 — Dipôle RLC série complet]
\begin{enumerate}
\item La tension $u(t) = \SI{20\sqrt{2}}{\volt}\cos(\SI{100\pi}{s^{-1}} t)$ a pour amplitude $U_m = \SI{20\sqrt{2}}{\volt}$, donc $U = \dfrac{U_m}{\sqrt{2}} = \SI{20}{\volt}$. Pulsation : $\omega = \SI{100\pi}{s^{-1}} \approx \SI{314}{s^{-1}}$ ; fréquence : $N = \dfrac{\omega}{2\pi} = \SI{50}{\hertz}$.
\item Réactances :
\[ L\omega = \SI{0,2}{\henry} \times \SI{100\pi}{s^{-1}} \approx \SI{62,8}{\ohm} \qquad ; \qquad \dfrac{1}{C\omega} = \dfrac{1}{\SI{5e-6}{\farad} \times \SI{100\pi}{s^{-1}}} \approx \SI{636,6}{\ohm}. \]
\[ Z = \sqrt{\SI{50}{\ohm}^2 + (\SI{62,8}{\ohm} - \SI{636,6}{\ohm})^2} = \sqrt{2500 + 329\,500} \approx \SI{576}{\ohm}. \]
\[ \tan\varphi = \dfrac{62,8 - 636,6}{50} \approx -11{,}5 \quad \Rightarrow \quad \varphi \approx -85^{\circ} \approx -\SI{1,48}{\radian}. \]
Le circuit est fortement capacitif.
\item Intensité efficace : $I = \dfrac{U}{Z} = \dfrac{\SI{20}{\volt}}{\SI{576}{\ohm}} \approx \SI{3,47e-2}{\ampere}$, soit $I \approx \SI{34,7}{\milli\ampere}$. L'intensité est en avance de $85^{\circ}$ sur la tension :
\[ i(t) = \SI{3,47e-2}{\ampere}\sqrt{2}\,\cos(\SI{100\pi}{s^{-1}} t + \SI{1,48}{\radian}). \]
\item Tensions efficaces :
\[ U_R = RI = \SI{50}{\ohm} \times \SI{0,0347}{\ampere} \approx \SI{1,74}{\volt} \ ; \quad U_L = L\omega I \approx \SI{62,8}{\ohm} \times \SI{0,0347}{\ampere} \approx \SI{2,18}{\volt} \ ; \quad U_C = \dfrac{I}{C\omega} \approx \SI{636,6}{\ohm} \times \SI{0,0347}{\ampere} \approx \SI{22,1}{\volt}. \]
On constate que $U_R + U_L + U_C \approx \SI{26}{\volt} \neq U = \SI{20}{\volt}$ : les tensions ne sont pas en phase, ce sont leurs \emph{vecteurs de Fresnel} qui s'ajoutent, pas leurs valeurs efficaces. On remarque même que $U_C > U$ : il y a déjà un début de surtension aux bornes du condensateur.
\item Diagramme de Fresnel : on porte $\overrightarrow{U_R}$ horizontalement (\SI{0,35}{\centi\meter}), $\overrightarrow{U_L}$ verticalement vers le haut (\SI{0,44}{\centi\meter}), $\overrightarrow{U_C}$ verticalement vers le bas (\SI{4,4}{\centi\meter}). La somme vectorielle donne $\vec{U}$, presque verticale vers le bas, de longueur \SI{4}{\centi\meter} (soit \SI{20}{\volt}), en retard de $85^{\circ}$ sur $\overrightarrow{U_R}$ : on retrouve le caractère capacitif.
\item Première méthode : $P = UI\cos\varphi = \SI{20}{\volt} \times \SI{0,0347}{\ampere} \times \cos 85^{\circ} \approx \SI{6,0e-2}{\watt}$. Deuxième méthode : $P = RI^2 = \SI{50}{\ohm} \times (\SI{0,0347}{\ampere})^2 \approx \SI{6,0e-2}{\watt}$. Les deux résultats sont identiques, ce qui confirme que seule la résistance dissipe de la puissance moyenne.
\end{enumerate}
\end{corrige}

\begin{corrige}[Exercice 6 — Surtension aux bornes du condensateur]
\begin{enumerate}
\item À la résonance, $Z = R$ donc $I_{\max} = \dfrac{U}{R}$. La tension efficace aux bornes du condensateur vaut :
\[ U_C = \dfrac{I_{\max}}{C\omega_0} = \dfrac{U}{RC\omega_0}. \]
Or $Q = \dfrac{1}{RC\omega_0}$, donc $\boxed{U_C = Q\,U}$. (De même $U_L = L\omega_0 I_{\max} = \dfrac{L\omega_0}{R}U = QU$.)
\item Application : $U_C = 25 \times \SI{12}{\volt} = \SI{300}{\volt}$.
\item Le condensateur supporte au maximum \SI{250}{\volt} efficaces : or $U_C = \SI{300}{\volt} > \SI{250}{\volt}$, il y a \textbf{danger} de claquage (destruction) du condensateur. Le phénomène de surtension à la résonance impose de choisir des composants dont la tension de service est largement supérieure à la tension d'alimentation.
\end{enumerate}
\end{corrige}

\begin{corrige}[Exercice 7 — Bande passante et facteur de qualité]
\begin{enumerate}
\item L'intensité efficace s'écrit $I = \dfrac{U}{Z} = \dfrac{U}{\sqrt{R^2 + \left(L\omega - \frac{1}{C\omega}\right)^2}}$, maximale à la résonance : $I_{\max} = \dfrac{U}{R}$. La condition $I = \dfrac{I_{\max}}{\sqrt{2}}$ s'écrit :
\[ \dfrac{U}{\sqrt{R^2 + \left(L\omega - \frac{1}{C\omega}\right)^2}} = \dfrac{U}{R\sqrt{2}} \quad \Rightarrow \quad R^2 + \left(L\omega - \dfrac{1}{C\omega}\right)^2 = 2R^2, \]
soit $\left|L\omega - \dfrac{1}{C\omega}\right| = R$. Les deux pulsations de coupure vérifient donc $L\omega_2 - \dfrac{1}{C\omega_2} = R$ et $L\omega_1 - \dfrac{1}{C\omega_1} = -R$. En multipliant par $\dfrac{\omega}{L}$ et en utilisant $\omega_0^2 = \dfrac{1}{LC}$ :
\[ \omega_2^2 - \dfrac{R}{L}\omega_2 - \omega_0^2 = 0 \qquad ; \qquad \omega_1^2 + \dfrac{R}{L}\omega_1 - \omega_0^2 = 0. \]
En soustrayant membre à membre : $(\omega_2^2 - \omega_1^2) - \dfrac{R}{L}(\omega_2 + \omega_1) = 0$, d'où $(\omega_2 - \omega_1)(\omega_2 + \omega_1) = \dfrac{R}{L}(\omega_2 + \omega_1)$, et finalement :
\[ \boxed{\Delta\omega = \omega_2 - \omega_1 = \dfrac{R}{L}}. \]
\item Facteur de qualité :
\[ Q = \dfrac{\omega_0}{\Delta\omega} = \dfrac{\omega_0 L}{R} = \dfrac{L}{R}\times\dfrac{1}{\sqrt{LC}} = \dfrac{1}{R}\sqrt{\dfrac{L}{C}}. \]
\item Application numérique : $\omega_0 = \dfrac{1}{\sqrt{\SI{0,05}{\henry} \times \SI{1e-7}{\farad}}} = \dfrac{1}{\sqrt{5\times10^{-9}}} \approx \SI{1,41e4}{s^{-1}}$, d'où
\[ N_0 = \dfrac{\omega_0}{2\pi} \approx \SI{2250}{\hertz} \approx \SI{2,25}{\kilo\hertz}. \]
$\Delta\omega = \dfrac{R}{L} = \dfrac{\SI{20}{\ohm}}{\SI{0,05}{\henry}} = \SI{400}{s^{-1}}$, donc $\Delta N = \dfrac{400}{2\pi} \approx \SI{64}{\hertz}$. Enfin $Q = \dfrac{\omega_0}{\Delta\omega} = \dfrac{14\,142}{400} \approx 35{,}4$ : résonance très aiguë.
\end{enumerate}
\end{corrige}

\begin{corrige}[Exercice 8 — Relevage du facteur de puissance à Douala]
\begin{enumerate}
\item Puissance moyenne : $P = UI\cos\varphi = \SI{220}{\volt} \times \SI{8}{\ampere} \times 0{,}75 = \SI{1320}{\watt}$.
\item Puissance apparente : $S = UI = \SI{220}{\volt} \times \SI{8}{\ampere} = \SI{1760}{\volt\ampere}$. Puissance réactive : $\sin\varphi = \sqrt{1 - 0{,}75^2} \approx 0{,}661$, d'où $Q_r = UI\sin\varphi = \SI{1760}{\volt\ampere} \times 0{,}661 \approx \SI{1164}{var}$.
\item La puissance utile est conservée : $P = UI'\cos\varphi'$, d'où
\[ I' = \dfrac{P}{U\cos\varphi'} = \dfrac{\SI{1320}{\watt}}{\SI{220}{\volt} \times 0{,}95} \approx \SI{6,32}{\ampere}. \]
\item Réduction relative : $\dfrac{I - I'}{I} = \dfrac{8 - 6{,}32}{8} \approx 0{,}21$, soit une baisse d'environ $21\ \%$. Conclusion : à puissance utile égale, relever le facteur de puissance diminue l'intensité dans les lignes, donc les pertes par effet Joule ($P_J = rI^2$) et l'échauffement des câbles : ENEO y gagne en pertes réseau, et l'abonné évite les pénalités facturées quand $\cos\varphi$ est trop faible, tout en pouvant utiliser une installation de section plus modeste.
\end{enumerate}
\end{corrige}

\begin{corrige}[Exercice 9 — Type Bac : étude complète d'un dipôle RLC série]
\begin{enumerate}
\item \textbf{Équation électrique.} La loi d'additivité des tensions donne $u = u_R + u_L + u_C$ avec $u_R = Ri$, $u_L = L\dfrac{di}{dt}$ et $u_C = \dfrac{q}{C}$ avec $i = \dfrac{dq}{dt}$. En dérivant par rapport au temps ($\dfrac{du_C}{dt} = \dfrac{i}{C}$) :
\[ \boxed{L\dfrac{d^2i}{dt^2} + R\dfrac{di}{dt} + \dfrac{i}{C} = \dfrac{du}{dt} = -\SI{30\sqrt{2}}{\volt} \times \SI{1000}{s^{-1}}\,\sin(\SI{1000}{s^{-1}} t)}. \]
En régime sinusoïdal forcé, la solution permanente est $i(t) = I_m\cos(\SI{1000}{s^{-1}} t - \varphi)$.
\item \textbf{Impédance et déphasage.}
\begin{enumerate}
\item $L\omega = \SI{0,5}{\henry} \times \SI{1000}{s^{-1}} = \SI{500}{\ohm}$ ; $\dfrac{1}{C\omega} = \dfrac{1}{\SI{5e-6}{\farad}\times \SI{1000}{s^{-1}}} = \SI{200}{\ohm}$.
\[ Z = \sqrt{R^2 + \left(L\omega - \dfrac{1}{C\omega}\right)^2} = \sqrt{\SI{100}{\ohm}^2 + \SI{300}{\ohm}^2} = \sqrt{100\,000} \approx \SI{316}{\ohm}. \]
\item $\tan\varphi = \dfrac{L\omega - \frac{1}{C\omega}}{R} = \dfrac{300}{100} = 3$, d'où $\varphi \approx \SI{1,25}{\radian} \approx 71{,}6^{\circ}$. Comme $\varphi > 0$ ($L\omega > \dfrac{1}{C\omega}$), le circuit est \cle{inductif} : $u$ est en avance sur $i$.
\end{enumerate}
\item \textbf{Intensité.} $I = \dfrac{U}{Z} = \dfrac{\SI{30}{\volt}}{\SI{316}{\ohm}} \approx \SI{9,49e-2}{\ampere} \approx \SI{94,9}{\milli\ampere}$. D'où :
\[ i(t) = \SI{9,49e-2}{\ampere}\sqrt{2}\,\cos(\SI{1000}{s^{-1}} t - \SI{1,25}{\radian}). \]
\item \textbf{Diagramme de Fresnel.} $U_R = RI = \SI{100}{\ohm} \times \SI{0,0949}{\ampere} \approx \SI{9,5}{\volt}$ ; $U_L = L\omega I = \SI{500}{\ohm} \times \SI{0,0949}{\ampere} \approx \SI{47,4}{\volt}$ ; $U_C = \dfrac{I}{C\omega} = \SI{200}{\ohm} \times \SI{0,0949}{\ampere} \approx \SI{19,0}{\volt}$. À l'échelle 1 cm pour 5 V : $\overrightarrow{U_R}$ horizontal (1,9 cm), $\overrightarrow{U_L}$ vertical vers le haut (9,5 cm), $\overrightarrow{U_C}$ vertical vers le bas (3,8 cm). La somme $\vec{U}$ a pour composantes $(9{,}5\ ;\ 47{,}4 - 19{,}0 = 28{,}4)$ V, soit $U = \sqrt{9{,}5^2 + 28{,}4^2} \approx \SI{30}{\volt}$ : on retrouve bien la tension d'alimentation, et $Z = \dfrac{U}{I} \approx \dfrac{30}{0{,}0949} \approx \SI{316}{\ohm}$.
\item \textbf{Puissance.} $P = UI\cos\varphi = \SI{30}{\volt} \times \SI{0,0949}{\ampere} \times \cos 71{,}6^{\circ} \approx 30 \times 0{,}0949 \times 0{,}316 \approx \SI{0,90}{\watt}$. Vérification : $P = RI^2 = \SI{100}{\ohm} \times (\SI{0,0949}{\ampere})^2 \approx \SI{0,90}{\watt}$. La résistance de la bobine étant négligeable, la fraction dissipée dans la bobine est \textbf{nulle} : toute la puissance est dissipée dans le résistor (la bobine et le condensateur ne consomment pas de puissance moyenne).
\item \textbf{Résonance.}
\begin{enumerate}
\item $\omega_0 = \dfrac{1}{\sqrt{LC}} = \dfrac{1}{\sqrt{\SI{0,5}{\henry} \times \SI{5e-6}{\farad}}} = \dfrac{1}{\sqrt{2,5\times10^{-6}}} \approx \SI{632}{s^{-1}}$, d'où $N_0 = \dfrac{\omega_0}{2\pi} \approx \SI{100,6}{\hertz} \approx \SI{101}{\hertz}$.
\item À la résonance : $I_{\max} = \dfrac{U}{R} = \dfrac{\SI{30}{\volt}}{\SI{100}{\ohm}} = \SI{0,30}{\ampere}$ et $\varphi = 0$ (circuit résistif).
\item $Q = \dfrac{L\omega_0}{R} = \dfrac{\SI{0,5}{\henry} \times \SI{632}{s^{-1}}}{\SI{100}{\ohm}} \approx 3{,}16$. Surtension aux bornes du condensateur : $U_C = QU = 3{,}16 \times \SI{30}{\volt} \approx \SI{94,9}{\volt}$, soit plus de trois fois la tension d'alimentation : il faut un condensateur de tension de service suffisante.
\end{enumerate}
\end{enumerate}
\textbf{Barème indicatif (sur 10) :} 1 : 1,5 pt ; 2 : 2 pts ; 3 : 1 pt ; 4 : 2 pts ; 5 : 1,5 pt ; 6 : 2 pts.

\textbf{Points de vigilance :} citer la loi d'additivité des tensions avant d'écrire l'équation ; préciser que $\varphi$ est le déphasage de $u$ par rapport à $i$ ; justifier la nature inductive par le signe de $\varphi$ ou la comparaison $L\omega > \dfrac{1}{C\omega}$ ; à la résonance, écrire explicitement la condition $L\omega_0 = \dfrac{1}{C\omega_0}$ ; ne jamais additionner les tensions efficaces $U_R + U_L + U_C$.
\end{corrige}

\begin{corrige}[Exercice 10 — Type Bac : exploitation d'une courbe de résonance]
\begin{enumerate}
\item \textbf{Lecture graphique.}
\begin{enumerate}
\item Le maximum de la courbe est atteint pour $N_0 \approx \SI{2,25}{\kilo\hertz}$, avec $I_{\max} \approx \SI{60}{\milli\ampere}$.
\item $\dfrac{I_{\max}}{\sqrt{2}} = \dfrac{60}{1{,}41} \approx \SI{42,4}{\milli\ampere}$. La droite horizontale d'ordonnée \SI{42,4}{\milli\ampere} coupe la courbe en $N_1 \approx \SI{2,10}{\kilo\hertz}$ et $N_2 \approx \SI{2,42}{\kilo\hertz}$.
\item Bande passante : $\Delta N = N_2 - N_1 \approx \SI{0,32}{\kilo\hertz} = \SI{320}{\hertz}$. Facteur de qualité : $Q = \dfrac{N_0}{\Delta N} = \dfrac{2250}{320} \approx 7{,}0$.
\end{enumerate}
\item \textbf{Exploitation.}
\begin{enumerate}
\item À la résonance, $L\omega_0 = \dfrac{1}{C\omega_0}$, donc $Z = R$ et $I_{\max} = \dfrac{U}{R}$. D'où $R = \dfrac{U}{I_{\max}} = \dfrac{\SI{6}{\volt}}{\SI{0,060}{\ampere}} = \SI{100}{\ohm}$.
\item $N_0 = \dfrac{1}{2\pi\sqrt{LC}}$ donne $L = \dfrac{1}{4\pi^2 N_0^2 C} = \dfrac{1}{4\pi^2 \times (\SI{2250}{\hertz})^2 \times \SI{1e-7}{\farad}} \approx \dfrac{1}{2,0\times10^{1}} \approx \SI{5,0e-2}{\henry}$, soit $L \approx \SI{50}{\milli\henry}$.
\item $Q = \dfrac{1}{R}\sqrt{\dfrac{L}{C}} = \dfrac{1}{100}\sqrt{\dfrac{0,05}{10^{-7}}} = \dfrac{\sqrt{5\times10^{5}}}{100} \approx \dfrac{707}{100} \approx 7{,}1$. On retrouve, aux erreurs de lecture graphique près, la valeur obtenue à la question 1.c ($Q \approx 7{,}0$) : les deux méthodes sont cohérentes.
\end{enumerate}
\item \textbf{Interprétation.} Si $R$ est divisée par deux : $I_{\max} = \dfrac{U}{R}$ double (pic deux fois plus haut, $I_{\max} \approx \SI{120}{\milli\ampere}$) ; $Q = \dfrac{L\omega_0}{R}$ double, donc la bande passante $\Delta N = \dfrac{N_0}{Q}$ est divisée par deux : la résonance devient plus aiguë, le circuit plus \cle{sélectif}. Application : dans un récepteur radio, le circuit RLC d'accord sélectionne la fréquence de l'émetteur désiré (par exemple CRTV Radio ou une radio communautaire) en rejetant les stations voisines ; plus la résonance est aiguë, meilleure est la séparation des stations.
\end{enumerate}
\textbf{Barème indicatif (sur 10) :} 1 : 3 pts ; 2 : 4 pts ; 3 : 3 pts.

\textbf{Points de vigilance :} les lectures graphiques doivent être accompagnées des tracés (pointillés) sur la courbe ; justifier $I_{\max} = \dfrac{U}{R}$ par la condition de résonance $L\omega_0 = \dfrac{1}{C\omega_0}$ ; convertir les mA en A dans les calculs ; pour la question d'interprétation, raisonner qualitativement à partir des expressions de $I_{\max}$, $Q$ et $\Delta N$.
\end{corrige}

\begin{corrige}[Exercice 11 — Type Bac : mesures à l'oscilloscope]
\begin{enumerate}
\item \textbf{Mesures.}
\begin{enumerate}
\item La courbe de la voie A culmine à 2 divisions : $U_m = 2 \times \SI{5}{\volt} = \SI{10}{\volt}$, d'où $U = \dfrac{U_m}{\sqrt{2}} \approx \SI{7,07}{\volt}$. La courbe de la voie B culmine à 1,6 division : $U_{Rm} = 1{,}6 \times \SI{5}{\volt} = \SI{8}{\volt}$, d'où $U_R = \dfrac{8}{\sqrt{2}} \approx \SI{5,66}{\volt}$.
\item Une période occupe 5 divisions : $T = 5 \times \SI{0,5}{\milli\second} = \SI{2,5}{\milli\second}$. Fréquence : $N = \dfrac{1}{T} = \dfrac{1}{\SI{2,5e-3}{\second}} = \SI{400}{\hertz}$ ; pulsation : $\omega = 2\pi N = \SI{800\pi}{s^{-1}} \approx \SI{2513}{s^{-1}}$.
\item Le décalage temporel entre les deux maxima vaut $\Delta t = 0{,}5\ \text{div} \times \SI{0,5}{\milli\second/div} = \SI{0,25}{\milli\second}$. Le déphasage en valeur absolue :
\[ |\varphi| = 2\pi\dfrac{\Delta t}{T} = 2\pi \times \dfrac{0{,}25}{2{,}5} = \dfrac{2\pi}{10} \approx \SI{0,63}{\radian} \approx 36^{\circ}. \]
\end{enumerate}
\item \textbf{Exploitation.}
\begin{enumerate}
\item Aux bornes du résistor, $u_R = Ri$ : la tension $u_R$ est proportionnelle à $i$ et en phase avec elle ; la voie B donne donc l'image de $i(t)$ (au facteur $R$ près). Intensité efficace : $I = \dfrac{U_R}{R} = \dfrac{\SI{5,66}{\volt}}{\SI{50}{\ohm}} \approx \SI{0,113}{\ampere}$.
\item Impédance : $Z = \dfrac{U}{I} = \dfrac{\SI{7,07}{\volt}}{\SI{0,113}{\ampere}} \approx \SI{62,5}{\ohm}$.
\item Le maximum de la voie A est atteint \emph{avant} celui de la voie B : $u$ est en avance sur $i$, donc $\varphi = +36^{\circ} > 0$ : le circuit est \cle{inductif}.
\item Puissance moyenne : $P = UI\cos\varphi = \SI{7,07}{\volt} \times \SI{0,113}{\ampere} \times \cos 36^{\circ} \approx 7,07 \times 0,113 \times 0,809 \approx \SI{0,65}{\watt}$. Vérification : $P = RI^2 = \SI{50}{\ohm} \times (\SI{0,113}{\ampere})^2 \approx \SI{0,64}{\watt}$ — accord satisfaisant compte tenu des lectures.
\end{enumerate}
\item \textbf{Pour aller plus loin.} $\tan\varphi = \dfrac{L\omega - \frac{1}{C\omega}}{R}$, donc $L\omega - \dfrac{1}{C\omega} = R\tan\varphi = 50 \times \tan 36^{\circ} \approx \SI{36,3}{\ohm}$. Or $L\omega = \SI{0,1}{\henry} \times \SI{2513}{s^{-1}} \approx \SI{251,3}{\ohm}$, d'où $\dfrac{1}{C\omega} = 251,3 - 36,3 = \SI{215}{\ohm}$ et :
\[ C = \dfrac{1}{215 \times 2513} \approx \SI{1,85e-6}{\farad} \approx \SI{1,9}{\micro\farad}. \]
Pour observer la résonance d'intensité, il faut faire varier la fréquence du GBF jusqu'à ce que les deux courbes soient en phase ($\varphi = 0$, $\Delta t = 0$) : l'amplitude de $u_R$ est alors maximale.
\end{enumerate}
\textbf{Barème indicatif (sur 10) :} 1 : 4 pts ; 2 : 4 pts ; 3 : 2 pts.

\textbf{Points de vigilance :} ne pas oublier le facteur $\sqrt{2}$ entre amplitude et valeur efficace ; justifier que $u_R$ est l'image de $i$ par la loi d'Ohm $u_R = Ri$ ; la formule $|\varphi| = 2\pi\dfrac{\Delta t}{T}$ doit être écrite avant l'application numérique ; le sens du décalage (avance ou retard) se lit sur la position des maxima et doit être explicité pour conclure sur la nature du circuit.
\end{corrige}

\newpage

\coursec{Fiche bilan}

\begin{popfigure}
\begin{tikzpicture}[mindmap, grow cyclic, every node/.style={concept, execute at begin node=\hspace{0pt}}, concept color=popblue, root concept/.append style={concept color=popblue, font=\bfseries\small, minimum size=2.6cm}, level 1/.append style={level distance=3.4cm, sibling angle=60, font=\scriptsize}, level 2/.append style={level distance=2.2cm, sibling angle=45, font=\tiny}, text=white]
\node[root concept]{Oscillations forcées RLC série}
  child[concept color=poporange]{ node {Circuit RLC forcé}
    child { node {$u_R{+}u_L{+}u_C=u$} }
    child { node {$L\ddot{q}+R\dot{q}+\dfrac{q}{C}=u$} }
  }
  child[concept color=popgreen]{ node {Fresnel}
    child { node {$Z=\sqrt{R^2+\left(L\omega-\dfrac{1}{C\omega}\right)^2}$} }
    child { node {$\tan\varphi=\dfrac{L\omega-\frac{1}{C\omega}}{R}$} }
  }
  child[concept color=poppurple]{ node {Résonance}
    child { node {$\omega_0=\dfrac{1}{\sqrt{LC}}$} }
    child { node {$I_{\max}=\dfrac{U}{R}$} }
  }
  child[concept color=popgold]{ node {Bande passante}
    child { node {$\Delta\omega=\dfrac{R}{L}$} }
    child { node {$Q=\dfrac{\omega_0}{\Delta\omega}=\dfrac{L\omega_0}{R}$} }
  }
  child[concept color=poppink]{ node {Puissance}
    child { node {$P=UI\cos\varphi$} }
    child { node {$\cos\varphi=\dfrac{R}{Z}$} }
  }
  child[concept color=popblue!70]{ node {Oscilloscope}
    child { node {$\varphi=2\pi\,\dfrac{\Delta t}{T}$} }
  };
\end{tikzpicture}
\legende{Carte mentale de la leçon : les six compétences clés des oscillations électriques forcées.}
\end{popfigure}

\begin{bilan}
\textbf{Définitions clés}
\begin{itemize}
  \item \cle{Régime forcé} : le dipôle RLC série est alimenté par un générateur de tension sinusoïdale $u(t)=U_m\cos(\omega t)$ ; après le régime transitoire, $i(t)=I_m\cos(\omega t-\varphi)$ oscille à la \textbf{même pulsation} $\omega$ que le générateur.
  \item \cle{Impédance} $Z$ (en $\Omega$) : quotient des valeurs efficaces $Z=\dfrac{U}{I}$.
  \item \cle{Déphasage} $\varphi$ : avance de phase de $u$ sur $i$ ; $\varphi>0$ : circuit inductif ; $\varphi<0$ : circuit capacitif.
  \item \cle{Résonance d'intensité} : $I$ est maximale quand $\omega=\omega_0$.
  \item \cle{Bande passante à $-3$ dB} : intervalle $[\omega_1\,;\omega_2]$ où $I\geq \dfrac{I_{\max}}{\sqrt{2}}$.
  \item \cle{Facteur de qualité} $Q$ : mesure de l'acuité de la résonance (sans unité).
\end{itemize}

\textbf{Formules à connaître par cœur}
\begin{center}
\begin{tabularx}{\linewidth}{|l|X|}
\hline
\rowcolor{popblueL} \textbf{Grandeur} & \textbf{Expression} \\
\hline
Équation du circuit & $L\dfrac{di}{dt}+Ri+\dfrac{q}{C}=u(t)$, avec $i=\dfrac{dq}{dt}$ \\
\hline
Impédances & $Z_R=R$ ; \quad $Z_L=L\omega$ ; \quad $Z_C=\dfrac{1}{C\omega}$ \\
\hline
Impédance RLC & $Z=\sqrt{R^2+\left(L\omega-\dfrac{1}{C\omega}\right)^2}$ \\
\hline
Déphasage & $\tan\varphi=\dfrac{L\omega-\dfrac{1}{C\omega}}{R}$ \quad et \quad $\cos\varphi=\dfrac{R}{Z}$ \\
\hline
Résonance & $\omega_0=\dfrac{1}{\sqrt{LC}}$, \quad $f_0=\dfrac{1}{2\pi\sqrt{LC}}$, \quad $Z_{\min}=R$, \quad $I_{\max}=\dfrac{U}{R}$, \quad $\varphi=0$ \\
\hline
Bande passante & $\Delta\omega=\omega_2-\omega_1=\dfrac{R}{L}$ \\
\hline
Facteur de qualité & $Q=\dfrac{\omega_0}{\Delta\omega}=\dfrac{L\omega_0}{R}=\dfrac{1}{R}\sqrt{\dfrac{L}{C}}=\dfrac{1}{RC\omega_0}$ \\
\hline
Puissance moyenne & $P=UI\cos\varphi=RI^2$ \quad (seul $R$ consomme) \\
\hline
Oscilloscope & $\varphi=2\pi\,\dfrac{\Delta t}{T}$ \quad ($\Delta t$ : décalage temporel entre les deux courbes) \\
\hline
\end{tabularx}
\end{center}

\textbf{Méthodes express}
\begin{itemize}
  \item \textbf{Fresnel} : représenter $\vec{U}_R$ (en phase avec $i$), $\vec{U}_L$ (avance de $+\dfrac{\pi}{2}$), $\vec{U}_C$ (retard de $-\dfrac{\pi}{2}$) ; la somme vectorielle donne $\vec{U}$, d'où $Z$ et $\varphi$ par le théorème de Pythagore.
  \item \textbf{Courbe de résonance} : lire $I_{\max}$ et $\omega_0$ au sommet ; tracer la ligne $I=\dfrac{I_{\max}}{\sqrt{2}}$ pour lire $\omega_1$ et $\omega_2$ ; en déduire $\Delta\omega$ puis $Q$.
  \item \textbf{Signe de $\varphi$} : si $\omega>\omega_0$, le circuit est inductif ($\varphi>0$) ; si $\omega<\omega_0$, il est capacitif ($\varphi<0$).
  \item \textbf{Oscilloscope} : mesurer la période $T$ et le décalage $\Delta t$ entre les passages à zéro dans le même sens ; la courbe en avance est celle qui passe par zéro \emph{en premier}.
\end{itemize}
\end{bilan}

\begin{aretenir}[Checklist avant le Bac]
\begin{itemize}
  \item[$\square$] Je sais écrire l'équation différentielle du circuit RLC série en $q$ ou en $i$.
  \item[$\square$] Je connais les impédances $Z_R=R$, $Z_L=L\omega$, $Z_C=\dfrac{1}{C\omega}$ et leurs déphasages.
  \item[$\square$] Je sais construire la représentation de Fresnel et en déduire $Z$ et $\tan\varphi$.
  \item[$\square$] Je sais que la résonance a lieu pour $\omega_0=\dfrac{1}{\sqrt{LC}}$ avec $Z=R$ et $\varphi=0$.
  \item[$\square$] Je sais exploiter la courbe $I(\omega)$ : $I_{\max}$, bande passante à $\dfrac{I_{\max}}{\sqrt{2}}$.
  \item[$\square$] Je sais calculer $Q=\dfrac{\omega_0}{\Delta\omega}=\dfrac{L\omega_0}{R}$ et interpréter l'acuité.
  \item[$\square$] Je sais calculer la puissance moyenne $P=UI\cos\varphi$ et le facteur de puissance $\cos\varphi=\dfrac{R}{Z}$.
  \item[$\square$] Je sais mesurer un déphasage à l'oscilloscope : $\varphi=2\pi\dfrac{\Delta t}{T}$.
  \item[$\square$] Je vérifie toujours les unités : $Z$ en $\Omega$, $\omega$ en $\mathrm{rad\,s^{-1}}$, $P$ en W, $Q$ sans unité.
  \item[$\square$] Je distingue valeurs maximales et efficaces : $U_m=U\sqrt{2}$, $I_m=I\sqrt{2}$.
  \item[$\square$] Je n'oublie pas que seule la résistance consomme de la puissance moyenne.
  \item[$\square$] Je sais déterminer le caractère inductif ou capacitif du circuit selon le signe de $L\omega-\dfrac{1}{C\omega}$.
\end{itemize}
\end{aretenir}

\findecours

\end{document}
